\documentclass[xcolor=table,aspectratio=169]{beamer}

%\usepackage[english]{babel} % English language/hyphenation

\usepackage[defaultsans,scale=0.95]{opensans}
\usepackage{FiraMono}
\usepackage{fourier} % Use the Adobe Utopia font for the document - comment this line to return to the LaTeX default
\usefonttheme{professionalfonts} % https://tex.stackexchange.com/a/330700/235983
\DeclareMathAlphabet{\mathcal}{OMS}{cmsy}{m}{n}
\newcommand\bmmax{2} % used to fix issues with bbm and e.g. esint
\usepackage{bm} % for vec/mat
\usepackage[T1]{fontenc}

\usepackage{pdfpages}
\usepackage{graphicx}
\usepackage{caption}
\captionsetup[figure]{labelformat=empty}
\usepackage{subcaption}
\usepackage{multicol}

\usepackage{tikz}
\usepackage{wrapfig}

% COLORS
\usepackage{xcolor}
\colorlet{myred}{red!80!black}
\colorlet{myblue}{blue!80!black}
\colorlet{mygreen}{green!60!black}
\colorlet{myorange}{orange!70!red!60!black}
\colorlet{mydarkred}{red!30!black}
\colorlet{mydarkblue}{blue!40!black}
\colorlet{mydarkgreen}{green!30!black}

\colorlet{mylightblue}{blue!60!white}
\colorlet{mydarkyellow}{yellow!70!black}
\colorlet{myyellowred}{yellow!70!red}

\usepackage[backend=biber,style=numeric]{biblatex}
\addbibresource{template_for_qibin_bib.bib}

%\usepackage{tabu}  % \(\)-able tabulars
\usepackage{float} % H placement
\usepackage{hyperref}

\usepackage{amsmath, amsthm, amssymb, amsfonts}
% \usepackage{mbboard} % requires manual installation in texmf
\usepackage{esint}
\usepackage{mathtools} % boxed aligns, delimiters
\usepackage{listings}
\lstset{basicstyle=\ttfamily}

\usepackage{physics} % http://mirrors.ibiblio.org/CTAN/macros/latex/contrib/physics/physics.pdf
\usepackage[makeroom]{cancel}
% \usepackage{slashed}
\usepackage{siunitx}
\sisetup{inter-unit-product=\ensuremath{{}\cdot{}}}
\DeclareSIUnit\eVperc{\eV\per\clight}
\AtBeginDocument{\RenewCommandCopy\qty\SI}
\usepackage[version=4]{mhchem}
%\usepackage{hepnames}
\usepackage[italic]{heppennames}
\usepackage{xpatch} % https://tex.stackexchange.com/a/541095
\makeatletter
\xpatchcmd\@HepConStyle
{\edef\@upcode{\updefault}}
{\ifdefined\shapedefault\edef\@upcode{\shapedefault}\else\edef\@upcode{\updefault}\fi}
{}{}
\makeatother
\usepackage{tensor}
\usepackage{feynmp-auto}
\usepackage{listofitems}
\usetikzlibrary{arrows.meta} % for arrow size
\usepackage[outline]{contour} % glow around text

\addtobeamertemplate{frametitle}{}{\vspace{-15.5pt}}

%%%%%%%%%%%%%%%%%%%%%%%%%%%%%%%%%%%%%%

\newcommand{\fixme}{\text{{\Huge \(\mathbb{FIXME}\)}}}
\newcommand{\todo}{\text{{\huge \(\mathbb{TODO}\)}}}
\newcommand{\lboxed}[1]{\begin{array}{|l}
		\hline\raisebox{3pt}{\(\displaystyle\vphantom{#1}\)}
		\displaystyle#1\!\!\!\!\!
		\raisebox{-3pt}{\(\displaystyle\vphantom{#1}\)}\\\hline
	\end{array}
}
\newcommand{\rboxed}[1]{\begin{array}{r|}
		\hline\raisebox{3pt}{\(\displaystyle\vphantom{#1}\)}
		\!\!\!\!\!\displaystyle#1
		\raisebox{-3pt}{\(\displaystyle\vphantom{#1}\)}\\\hline
	\end{array}
}
\newcommand{\numberthis}{\addtocounter{equation}{1}\tag{\theequation}}% https://tex.stackexchange.com/a/42728

\newcommand{\Lagr}{\mathcal{L}}

\setlength\arrayrulewidth{1.0pt}

\usepackage{booktabs}
\usepackage{subcaption}
%\newcommand\blfootnote[1]{%
%	\begingroup
%	\renewcommand\thefootnote{}\footnote{#1}%
%	\addtocounter{footnote}{-1}%
%	\endgroup
%}
% Macro to skip item in enum
% https://tex.stackexchange.com/a/201162/235983
\newcommand{\skipitems}[1]{%
	\addtocounter{\@enumctr}{#1}%
}
% Macros to start and end Backup section with unnumbered slides
\newcommand{\backupbegin}{
    \newcounter{finalframe}
    \setcounter{finalframe}{\value{framenumber}}
    \setbeamertemplate{footline}{} % Remove the footline for Backup section
    \frame{\usebeamertemplate*{backup page}}
}

\newcommand{\backupend}{
    \setcounter{framenumber}{\value{finalframe}}
    \setbeamertemplate{footline}[frame number] % Restore original slide numbering
}
% Macro to add TOC entry for starred section
% https://tex.stackexchange.com/a/611664/235983
\makeatletter
\newcommand{\nonumsec}{%
	\addtocontents{toc}{\protect\beamer@sectionintoc{\the\c@section}{\secname}{\the\c@page}{\the\c@part}{0}}%
}
\makeatother

\def\thisframelogos{}

\newcommand{\framelogo}[1]{\def\thisframelogos{#1}}

\addtobeamertemplate{frametitle}{}{%
\begin{tikzpicture}[remember picture,overlay]
\node[anchor=north east,yshift=3.5pt] at (current page.north east) {%
    \foreach \img in \thisframelogos {%
        \hspace{.5ex}%
        \includegraphics[height=0.8cm, keepaspectratio]{\img}%
    }%
};
\end{tikzpicture}}

%%%%%%%%%%%%%%%%%%%%%%%%%%%%%%%%%%%%%%

\usetheme{Philadelphia}
%\usefonttheme[onlymath]{serif}
\hypersetup{colorlinks,allcolors=pennred}

\title[PV-Finder]{PV-Finder Updates}

\author[Qi Bin Lei]{\href{https://www.linkedin.com/in/qibinlei/}{Qi Bin~Lei}}

\institute[]{Stanford University}

\date{Compiled Slides 2024-2025}

% Replace these with the path to the logos
\titlegraphic{
    \hspace{5pt}
	\includegraphics[height=1cm]{logo-Stanford/Block_S_2_color.png}
	\includegraphics[trim={0 11.5cm 0 11.5cm},clip,height=1cm]{logo-ATLAS/ATLAS-Logo-Ref-Pantone.pdf}
    \includegraphics[height=1cm]{logo-iris-hep/Iris-hep-4-no-long-name.png}
}

\begin{document}

\begin{frame}[plain]
    \titlepage
\end{frame}
% \begin{frame}{Outline}
% 	\tableofcontents
% \end{frame}
% \begin{frame}{Title}
% 	\begin{itemize}
% 		\item An item
% 	\end{itemize}
% 	\begin{columns}
% 		\column{0.5\linewidth}
% 		Hello
% 		\column{0.5\linewidth}
% 		World
% 	\end{columns}
% \end{frame}
% \section{Section 1}
% \frame{\sectionpage}

\framelogo{"logo-ATLAS/ATLAS-Logo-White-transparent.png"}

\begin{frame}[plain]
    \centering
    TVPD October 2024
\end{frame}

\section{Introduction}

\begin{frame}{Background}
    \begin{columns}
        \begin{column}{.6\textwidth}
            \begin{itemize}
                \item B.A. Physics, University of Pennsylvania, Graduated Spring '24
                \item Post-Baccalaureate Position working with Dr. Rocky Garg on PV Finder
                \item Previous projects in the ATLAS Experiment under \href{https://live-sas-physics.pantheon.sas.upenn.edu/people/standing-faculty/evelyn-thomson}{Prof. Evelyn Thomson}
                \begin{itemize}
                    \item ML Jet Calibration (small-R Jets MCJES)
                    \item Chargino and Neutralino Analysis
                \end{itemize}
                \item Applying to Graduate School in the Fall
            \end{itemize}
        \end{column}

        \begin{column}{.4\textwidth}
            \begin{figure}
                \centering
                \includegraphics[width=.7\linewidth]{self_DRL.jpg}
            \end{figure}
        \end{column}
    \end{columns}
\end{frame}

\begin{frame}{Introduction}
    \begin{figure}
        \centering
        \begin{tikzpicture}
            % Place image
            \node[anchor=south west,inner sep=0] (image) at (0,0)
            {\includegraphics[width=0.6\linewidth]{oct25th/primary_vertices_run3.png}};

            % Draw an arrow for primary vertices
            \draw[->,mylightblue,thick] (-1,3) -- (2,2) node[anchor=south east, xshift=-60pt, yshift=25pt] {Primary Vertices};

            % Draw an arrow for tracks
            \draw[->,myyellowred,thick] (-1.5,1) -- (1,2) node[anchor=south east, xshift=-75pt, yshift=-30pt] {Tracks};

            % Draw an arrow for cones
            \draw[->,mydarkyellow,thick] (-2,5) -- (1,5) node[anchor=south east, xshift=-80pt] {Jet Cone};
        \end{tikzpicture}
        \caption{Run 3 event display with \(\left<\mu\right> = 65\)}
        \label{fig:run3_event_display}
    \end{figure}
\end{frame}

\begin{frame}{Tracks}
    \begin{columns}
        % Left column
        \column{0.3\linewidth}
        \begin{itemize}
            \item Gathering charged particle measurements readout from the detector
            \item Calculate trajectory estimation
            \item Used as inputs for reconstruction efforts and analysis
        \end{itemize}

        % Right column
        \column{0.7\linewidth}
        \begin{figure}
            \includegraphics[width=0.9\linewidth]{oct25th/coordinates_track.png}
            \caption{Reconstructed Track Data}
            \label{fig:coordinate_tracks}
        \end{figure}
    \end{columns}
\end{frame}

\begin{frame}{Clean, Merged, Split, Fake Prior Results}
    \begin{figure}
        \includegraphics[width=0.5\linewidth]{dec18th/clean_merged_split.png}
        \label{fig:clean_merged_split}
        \cite{ATLAS:2023hxv}
    \end{figure}
\end{frame}

\begin{frame}{Current Primary Vertex Algorithm: AVMF}
    % Talk about innate challenges that comes with using AVMF
    \begin{columns}
        % Text
        \column{0.4\linewidth}
        \begin{itemize}
            \item Global approach to vertex finding and fitting
            \item Finds and fits tracks based on compatibility to different primary vertices
            \item Limitations
            \begin{enumerate}
                \item Parallelization
                \item Time scale increases non linearly with number of PVs
            \end{enumerate}
        \end{itemize}

        % Text
        \column{0.6\linewidth}
        \begin{figure}
            \centering
            \includegraphics[width=0.9\linewidth]{oct25th/AVMF.png}
            \label{fig:AVMF}
            \cite{ATLAS:2019jmx}
        \end{figure}
    \end{columns}
\end{frame}


\begin{frame}{Data Statistics}
    \begin{itemize}
        \item Total Events: 51000
        \item Input Features: \(d_0, z_0, \sigma(d_0), \sigma(z_0), \sigma(d_0, z_0)\) (\(\sim 1000\) tracks per event)
        \begin{itemize}
            \item \(z_0 \in [-240\text{mm}, 240\text{mm}]\)
        \end{itemize}
        \item Output Label: KDE-A (12,000 Bins)
    \end{itemize}
    \begin{itemize}
        \item How should we make it easier for the neural network to learn?
        \begin{itemize}
            \item Padding input features make it all the same length
            \item Splitting KDE-A bins
        \end{itemize}
    \end{itemize}
\end{frame}


\begin{frame}{Masking}
    \begin{columns}
        \column{0.5\linewidth}
        \begin{itemize}
            \item A mask is used to identify valid tracks that should contribute to the final result (not padded value)
            \item The masking matrix, f2, is then multiplied with the output of the neural network to either
            \begin{enumerate}
                \item Zero out outputs of the neural network that do not have any track data associated with it
                \item Contribute values that will be summed to produce the predicted KDE bins
            \end{enumerate}
        \end{itemize}

        % Masking graphic
        \column{0.5\linewidth}
        \begin{figure}
            \centering
            \includegraphics[width=0.5\linewidth]{oct25th/creating_mask.png}
            \label{fig:masking_figure}
        \end{figure}
    \end{columns}
\end{frame}

\begin{frame}{Masking Output}
    \vspace{15pt}
    \begin{figure}
        \centering
        \includegraphics[width=0.8\linewidth]{oct25th/calculating_output(2).png}
        \label{fig:maskingtooutput}
    \end{figure}
\end{frame}

\begin{frame}{Loss Function}
    \begin{columns}
        % Equation and text
        \column{0.4\linewidth}
        \begin{itemize}
            \item Currently working with mean square error loss
            \begin{itemize}
                \item Will try out weighted mean square error loss to capture the peaks of the KDE better
            \end{itemize}
            \item Plateus for the first \(\sim 50-100\) epochs then drops exponentially
            \item MSE \(= \frac{1}{n} \sum (\text{KDE}_{\text{pred}} - \text{KDE}_{\text{truth}})^2\)
        \end{itemize}

        % Loss graph
        \column{0.6\linewidth}
        \begin{figure}
            \centering
            \includegraphics[width=1\linewidth]{oct25th/loss.png}
            \label{fig:loss}
        \end{figure}
    \end{columns}
\end{frame}

\begin{frame}[plain]
    \centering
    US LUA December 2024
\end{frame}


\begin{frame}{Towards HL-LHC}
    \begin{columns}
        % Text
        \column{0.4\linewidth}
        \begin{itemize}
            \item Current vertexing algorithms are iterative by nature
            % \begin{itemize}
            %     \item Preparation to account for more pileup events/detector effects is needed
            % \end{itemize}
            \item Motivations for ML
            \begin{itemize}
                \item Fast and accurate predictions
                \item Parallel training of events
                \item Integration of new variables
                \item Fine tuning for future upgrades/projects
            \end{itemize}
        \end{itemize}

        % HL-LHC Top-AntiTop Simultaion
        \column{0.6\linewidth}
        \begin{figure}
            \centering
            \includegraphics[width=0.9\linewidth]{oct25th/ATLAS_HLLHC_mu200_VP1_88vertices_angled_noID_4K_wLogo.png}
            \caption{Simulation of Top-Antitop Pair Production at HL-LHC [\href{http://cds.cern.ch/record/2846411/}{Link}]}
            \label{fig:HL-LHCSim}
        \end{figure}
    \end{columns}
\end{frame}

\section{PV-Finder/tracks to KDE}

\begin{frame}{PV-Finder Overview}
    \vspace{15pt}
    \begin{figure}
        \centering
        \includegraphics[width=1\linewidth]{dec18th/colored_trackstoHistsdiagram.png}
        \label{fig:PVFinderArchitecture}
    \end{figure}
\end{frame}

\begin{frame}{Data Pre-Processing}
    \vspace{15pt}
    \begin{figure}
        \centering
        \includegraphics[width=1\linewidth]{oct25th/data_preprocessing_pvfinder(1).png}
        \label{fig:data_preprocessing_pvfinder}
    \end{figure}
\end{frame}

\begin{frame}{Kernel Density Estimation}
    \begin{columns}
        % KDE Text
        \column{0.5\linewidth}
        \begin{itemize}
            \item Track's Gaussian Probability Density \\
            \(\mathbb{P}(r) = \frac{1}{2\pi\sqrt{|\Sigma|}} \exp{\alpha}\) \\
            \hspace{-25pt}
            \(\alpha = -\frac{1}{2}\left((d-d_0),(z-z_0)\right)^T \Sigma^{-1} \left((d-d_0),(z-z_0)\right)\) \\
            KDE-A = \(\sum_{\text{tracks}} \mathbb{P}(r)\) \\
            KDE-B = \(\sum_{\text{tracks}} \mathbb{P}(r)^2\)
            \item KDEs are composed of 12,000 bins
            \item Current analytical KDE production:
            \begin{itemize}
                \item Evaluates the Gaussian PDF for all    particle track
                \item \(\sim\)50 million evaluations per track per event
            \end{itemize}
        \end{itemize}

        % KDE Graphic
        \column{0.5\linewidth}
        \begin{figure}
            \centering
            \includegraphics[width=1\linewidth]{oct25th/KDE.png}
            \caption{From ATLAS Public Note \cite{ATLAS:2023hxv}}
            \label{fig:KDE}
        \end{figure}
    \end{columns}
\end{frame}

\begin{frame}{Tracks to KDE Architecture}
    \vspace{15pt}
    \begin{figure}
        \centering
        \includegraphics[width=1\linewidth]{oct25th/pipeline_trackstoKDE(2).png}
        \label{fig:pipeline_trackstokde}
    \end{figure}
\end{frame}

\begin{frame}{Example of Tracks to KDE Output (1)}
    \begin{figure}
        \centering
        \includegraphics[width=0.6\linewidth]{oct25th/trackstoKDE_results1.png}
        \label{fig:example_KDE_output_1}
    \end{figure}
\end{frame}

\begin{frame}{Example of Tracks to KDE Output (2)}
    \begin{figure}
        \centering
        \includegraphics[width=0.6\linewidth]{oct25th/trackstoKDE_results2.png}
        \label{fig:example_KDE_output_2}
    \end{figure}
\end{frame}

\begin{frame}{KDE to Hists Architecture}
    \begin{figure}
        \centering
        \includegraphics[width=1\linewidth]{dec18th/pipeline_KDEtoHists(1).png}
        \label{fig:KDEtoHists}
    \end{figure}
\end{frame}

\begin{frame}{Primary Vertex Information}
    \begin{figure}
        \centering
        \includegraphics[width=0.9\linewidth]{dec18th/primary_vertex_information.png}
        \caption{Target Truth Information \cite{ATLAS:2023hxv}}
        \label{fig:truthHistograms}
    \end{figure}
\end{frame}

\begin{frame}{Previous Work}
    \begin{columns}
        \column{0.4\linewidth}
        \begin{itemize}
            \item Trained using all KDEs produced from the combinatorial method
            \item Efficiency \(= \frac{\text{\# reconstructed vertices}}{\text{\# truth vertices}}\)
            \item Prior work noted in pubnote \cite{ATLAS:2023hxv}
        \end{itemize}

        \column{0.6\linewidth}
        \begin{figure}
            \centering
            \includegraphics[width=0.9\linewidth]{dec18th/current_results.png}
            \label{fig:KDEtoHistsResults}
        \end{figure}
    \end{columns}
\end{frame}

\begin{frame}{Example End-to-end KDE to Hists Output}
    \begin{columns}
        \column{0.4\linewidth}
        \begin{itemize}
            \item Training on predicted KDE-A
            \begin{itemize}
                \item Using prior outputs from tracks to KDE
                \item \(\sim 3\) days of training
            \end{itemize}
            \item Efficiency: \(\sim 62\%\)
            \item False Positive Rate: \(5.5\)
            \begin{itemize}
                \item Only improvements from here with more information
            \end{itemize}
        \end{itemize}

        \column{0.6\linewidth}
        \begin{figure}
            \centering
            \includegraphics[width=0.9\linewidth]{dec18th/ExamplePVFinderOutput4.png}
            \label{fig:example_Hists_output_1}
        \end{figure}
    \end{columns}
\end{frame}

\section{Summary}

\begin{frame}{Summary}
    \begin{itemize}
        \item Machine learning provides a viable avenue for primary vertex identificaiton
        \item NN-KDE network trained on one feature (KDE-A) is performing well
        \item Upcoming Work:
        \begin{itemize}
            \item Train NN-KDE on other KDEs (KDE-B-z, KDE-A-x, KDE-A-y)
            \item Training UNet and UNet++ architectures on NN-KDEs
            \item Compare preliminary efficiency with AMVF \cite{ATLAS:2023hxv}
        \end{itemize}
    \end{itemize}

    \centering
    \vspace{15pt}
    \textit{Thanks for listening!}
\end{frame}

\begin{frame}{Relevant Links and Directories}
    \begin{itemize}
        \item Ananya's prior work on tracks to KDE
        \begin{itemize}
            \item \href{https://github.com/qibinlei/TracksToKDE/tree/main}{Link to Gitlab}
            \item \href{https://docs.google.com/document/d/1oTlUo06nhkKBF059y9mCXGqZ46PPD3TuHFIe3EkTN-g/edit}{Link to Google Doc Documentation}
        \end{itemize}
        \item LHCb PVFinder
        \begin{itemize}
            \item \href{https://gitlab.cern.ch/sakar/pv-finder_v2/-/tree/master}{Link to Gitlab}
        \end{itemize}
        \item ATLAS PV Finder
        \begin{itemize}
            \item \href{https://gitlab.cern.ch/qlei/atlas_pvfinder}{Link to Gitlab}
        \end{itemize}
    \end{itemize}
\end{frame}

\begin{frame}{KDE to Hists Loss}
    \begin{figure}
        \centering
        \includegraphics[width=0.8\linewidth]{dec18th/loss_kde_a_unetplusplus.png}
        \label{fig:KDEHistsLoss}
    \end{figure}
\end{frame}

\begin{frame}{KDE to Hists Efficiency}
    \begin{figure}
        \centering
        \includegraphics[width=0.8\linewidth]{dec18th/efficiency.png}
        \label{fig:efficiency}
    \end{figure}
\end{frame}

\begin{frame}{KDE to Hists FPR}
    \begin{figure}
        \centering
        \includegraphics[width=0.8\linewidth]{dec18th/false_positive_rate.png}
        \label{fig:FPR}
    \end{figure}
\end{frame}

\begin{frame}{KDE to Hists Results}
    \begin{figure}
        \centering
        \includegraphics[width=0.6\linewidth]{dec18th/ExamplePVFinderOutput2.png}
        \label{fig:example_Hists_output_2}
    \end{figure}
\end{frame}

\begin{frame}[plain]
    \centering
    Goals Winter Quarter 2025
\end{frame}

% Goals Winter Quarter 2025
\begin{frame}{Recap from Last Quarter}
    \begin{itemize}
        \item The End Goal for PV Finder
        \begin{itemize}
            \item Deployment for general use in primary vertex reconstruction for the ATLAS experiment
            \item Provide a competitive approach to the existing AVMF method
            \item Publish a pubnote on results
        \end{itemize}
        \item Make Substantial Progress on the PV Finder Project
        \begin{itemize}
            \item Learn from existing literature from LHCb and other associated work \checkmark
            \item Start with tracks to KDE by building ontop of Ananya's work \checkmark
            \begin{itemize}
                \item Learn about the tools used in this project that are new to me (e.g. Pytorch, MLFlow, etc.) \checkmark
                \item Replicate prior results to serve as a baseline for improvement \checkmark
            \end{itemize}
            \item Connect framework onto existing KDE to hists \checkmark
            \begin{itemize}
                \item Revamping Gitlab and associated files to accomodate for more options (In Progress)
                \item Exploring training options and ways to efficiently improve the overall network (In Progress)
            \end{itemize}
        \end{itemize}
    \end{itemize}
\end{frame}

\begin{frame}{Focus on This Quarter}
    \begin{itemize}
        \item Start producing end-to-end results from tracks to primary vertex locations
        \begin{itemize}
            \item Separate training on all of the different KDEs (e.g. A-z, A-x, A-y, B-z) (In Progress)
            \item Combined neural network trainings (both DNN and UNet)
        \end{itemize}
        \item Start working on Muon Collider software project
        \item Personal Goals
        \begin{itemize}
            \item Start running again
            \item Explore Bay Area
        \end{itemize}
    \end{itemize}
\end{frame}

\begin{frame}{KDE B z Loss}
    \begin{figure}
        \centering
        \includegraphics[width=0.9\linewidth]{jan17th/KDEBz_jan17th_loss.png}
        \caption{Currently training for another 1000 epochs}
        \label{fig:KDE_B_z_example_jan17th_loss}
    \end{figure}
\end{frame}

\begin{frame}{KDE A x Loss}
    \begin{figure}
        \centering
        \includegraphics[width=0.9\linewidth]{jan17th/KDEAx_jan17th_loss.png}
        \caption{Starting to overfit}
        \label{fig:KDE_A_x_example_jan17th_loss}
    \end{figure}
\end{frame}

\begin{frame}{KDE A x Progress}
    \begin{figure}
        \centering
        \includegraphics[width=0.9\linewidth]{jan17th/KDEAx_jan17th.png}
        \label{fig:KDE_A_x_example_jan17th}
    \end{figure}
\end{frame}

\begin{frame}[plain]
    \centering
    January 24th Updates
\end{frame}

\begin{frame}{Training split KDE to split hists}
    \begin{itemize}
        \item Having the tracks to KDE output 100 KDE bins means that you have to either
        \begin{enumerate}
            \item Wait in the intermediary to recombine each 100 KDE bins until you reach 12000 bins to feed into next step
            \begin{itemize}
                \item Doesn't allow for randomization of inputs/outputs
                \item Requires a strict batch size
            \end{itemize}
            \item Feed the 100 KDE bins into CNN and expect 100 predicted histogram bins output
            \begin{itemize}
                \item Already what LHCb does
                \item Just need to retrain UNet/UNet++
            \end{itemize}
        \end{enumerate}
    \end{itemize}
\end{frame}

\begin{frame}{Preliminary UNet Training}
    \begin{itemize}
        \item Need to double check the event counting after spliting (may be why the false positive rate is really low)
        \item BUT, matches what we expect from prior UNet trainings
    \end{itemize}
    \begin{figure}
        \centering
        \includegraphics[width=0.7\linewidth]{jan23rd/preliminary_unet_training_split.png}
        \caption{Currently training for 500 epochs}
        \label{fig:prelim_unet_jan23rd}
    \end{figure}
\end{frame}

\begin{frame}{Current Separate End-to-End Training}
    \begin{itemize}
        \item Take the best results of KDE-A-z, KDE-B-z, KDE-A-x, KDE-A-y and feed it into a UNet
        \item Quick sanity test on if the current DNN works (spoiler: don't think so)
        \item Epoch 0: train=1.56262, val=1.00523, took 481.18 s                                                                         Validation Found 102533 of 274117, added 26008 (eff 37.40\%) (0.0212 FP/event)
        \item Fine... but this is just one epoch? Empirically, prior trainings on just a predicted KDE-A-z has resulted in (next slide, efficiency capped at 60\%)
    \end{itemize}
\end{frame}

\begin{frame}{KDE to Hists Efficiency}
    \begin{figure}
        \centering
        \includegraphics[width=0.8\linewidth]{dec18th/efficiency.png}
        \label{fig:efficiency}
    \end{figure}
\end{frame}

\begin{frame}{Proposal: Try POCA Variables}
    \begin{itemize}
        \item I think reco track variables are going to be the bottleneck - lets try POCA variables
        \item There are 14 variables we can use (can also add reco track parameters)
        \begin{itemize}
            \item poca x, poca y, poca z, poca z min, poca z max
            \item poca minoraxis1 x, poca minoraxis1 y, poca minoraxis1 z
            \item poca minoraxis2 x, poca minoraxis2 y, poca minoraxis2 z
            \item poca majoraxis x, poca majoraxis y, poca majoraxis z
        \end{itemize}
        \item LHCb uses a projection technique noted in \href{https://gitlab.cern.ch/sakar/pv-finder_v2/-/blob/master/tools/ellipsoid.py?ref_type=heads}{GitLab Link}, \href{https://math.stackexchange.com/questions/1865188/how-to-prove-the-parallel-projection-of-an-ellipsoid-is-an-ellipse}{Math Stack Exchange Link}
        \begin{itemize}
            \item Reduces 9 variables into 6 variables (major and minor axis) \(\rightarrow\) (A,B,C,D,E,F)
        \end{itemize}
    \end{itemize}
\end{frame}

\begin{frame}{A Quick Peak into the Math}
    \begin{itemize}
        \item General form of an ellipsoid:
    \end{itemize}
    \begin{align*}
        Ax^2+By^2+Cz^2+2(Dxy+Exz+Fyz)=1
    \end{align*}
    \begin{itemize}
        \item Axis Transformation
    \end{itemize}
    \begin{align*}
        \left( \vec{x} \cdot \hat{u}_1 \right) ^2 +
        \left( \vec{x} \cdot \hat{u}_2 \right) ^2 +
        \left( \vec{x} \cdot \hat{u}_3 \right) ^2 = 1
    \end{align*}
    \begin{itemize}
        \item Where \(\hat{u}_1, \hat{u}_2\) represents the minor axis and \(\hat{u}_3\) represents the major axis. Expanding out the terms we are able to get can equation for each letter.
    \end{itemize}
    \begin{align*}
        A = u_{1x}^2 + u_{2x}^2 + u_{3x}^2
    \end{align*}
\end{frame}

\begin{frame}{Infinities Found in POCA Variables}
    \begin{itemize}
        \item Need to go through each variable, remove each infinity, and make sure every POCA is the same length
    \end{itemize}
    \begin{figure}
        \centering
        \includegraphics[width=0.6\linewidth]{jan23rd/infinities_in_poca.png}
        \label{fig:infinities_in_poca}
    \end{figure}
\end{frame}

\begin{frame}{Interval Splitting? Track Selection?}
    \begin{itemize}
        \item Interval Splitting
        \begin{itemize}
            \item Option 1: Calculating z sigma and allowing for an interval extension of \(\pm 3 \sigma_{z}\)
            \begin{align*}
                \sigma_{z} = \frac{1}{\sqrt{C_{\text{POCA}}}}
            \end{align*}
            \item Option 2: Hard coded parameter of \(\pm 2.5\) mm (how LHCb does it)
        \end{itemize}
        \item Track Selection
        \begin{itemize}
            \item Set a maximum allowed \(\sigma_{x}, \sigma_{y}, \sigma_{z}\) values and if exceeded, remove track
            \item Should I implement something along the lines of this?
        \end{itemize}
    \end{itemize}
\end{frame}

\begin{frame}{Up Next:}
    \begin{itemize}
        \item Create POCA variable H5 files and start training KDEs on them
        \begin{itemize}
            \item Spending time plucking out the infinities and adding an epsilon terms
            \item Hopefully have preliminary predicted KDE results by next week
        \end{itemize}
        \item Overhaul cleaning up the Git repository and making it function similarly to LHCb
        \item Schedule meeting with Simon and Michael after POCA runs
    \end{itemize}
\end{frame}

\begin{frame}[plain]
    \centering
    January 30th Updates
\end{frame}

\begin{frame}{Overview of This Week's Work}
    \begin{itemize}
        \item End-to-end tracks to hist split with initialized weights using all KDEs
        \begin{itemize}
            \item Will resume training for another 100 epochs - had to prioritize other runs
            \item Epoch 39: train=1.31555, val=1.33709, took 506.91 s   Validation Found 170066 of 274117, added 53649 (eff 62.04\%) (0.0438 FP/event) 1.3155532292451846
        \end{itemize}
        \item POCA to KDE-A-z
        \begin{itemize}
            \item Plateau'd around 589 Loss, training on reconstructed track was able to reach 350
            \item Not sure why this is the case - will try again with non transformed variables
        \end{itemize}
        \item Changed UNet Runs (next slides)
        \item Trying out different loss functions for tracks-to-KDE (next slides)
        \begin{itemize}
            \item Normalized Chisq Loss Function (LHCb)
            \item KL Divergence
        \end{itemize}
        \item QoL Misc Stuff
    \end{itemize}
\end{frame}

\begin{frame}{End-to-end Tracks to Hists Preliminary Training (All KDEs)}
    \begin{figure}
        \centering
        \includegraphics[width=0.8\linewidth]{jan30th/end-to-end_t2h_sample.png}
    \end{figure}
\end{frame}

\begin{frame}{End-to-end Poca to Hists Preliminary Training (KDE-A)}
    \begin{figure}
        \centering
        \includegraphics[width=0.8\linewidth]{jan30th/poca_t2kde_sample.png}
    \end{figure}
\end{frame}

\begin{frame}{UNet Architecture}
    \begin{figure}
        \centering
        \includegraphics[width=0.8\linewidth]{jan30th/unet_architecture.png}
    \end{figure}
\end{frame}

\begin{frame}{UNet Architecture Code}
    \begin{figure}
        \centering
        \includegraphics[width=0.8\linewidth]{jan30th/unet_alt_maybe(1).jpg}
    \end{figure}
\end{frame}

\begin{frame}{Alt Loss Functions for Tracks to KDE}
    \begin{itemize}
        \item Normalized Chisq (MSE-Like Loss Function):
        \begin{align*}
            \sigma &= 0.01 \\
            \chi^2 &= \left( \frac{y_{\text{pred}} - y_{\text{truth}}}{\sigma} \right)^2 \\
            \chi^2_{\text{norm}} &= \frac{\chi^2}{\text{nBins} * \text{nBatch}}
        \end{align*}
        \item Kullback–Leibler divergence (Use to describe the distance between two distributions):
        \begin{align*}
            \Lagr_{\text{KR}} \left( y_{\text{pred}}, y_{\text{true}} \right) = \frac{1}{n}\sum y_{\text{true}} \cdot \log(\frac{y_{\text{true}}}{y_{\text{pred}}})
        \end{align*}
    \end{itemize}
\end{frame}

\begin{frame}{QoL: Cleaning Up Repository}
    \begin{itemize}
        \item Merged trackstoKDE and model folders - they all just live in model now
        \item All trainings are now just done through these scripts
    \end{itemize}
    \begin{figure}
        \centering
        \includegraphics[width=1\linewidth]{jan30th/gitlab_scripts_example.png}
    \end{figure}
\end{frame}

\begin{frame}{QoL: Config Files}
    \begin{columns}
        \begin{column}{.5\textwidth}
            \begin{figure}
                \centering
                \includegraphics[width=1\linewidth]{jan30th/config_file_example.png}
            \end{figure}
        \end{column}

        \begin{column}{.5\textwidth}
            \begin{figure}
                \centering
                \includegraphics[width=1\linewidth]{jan30th/script_to_generate_yaml_file.png}
            \end{figure}
        \end{column}
    \end{columns}
\end{frame}

\begin{frame}{Up Next:}
    \begin{itemize}
        \item Meeting with Simon and Michael
        \begin{itemize}
            \item Comment on UNet Architecture
            \item Making intermittent feature channels of UNet larger?
            \begin{itemize}
                \item Interesting paper!: \href{https://arxiv.org/pdf/1905.11946}{EfficientNet: Rethinking Model Scaling for Convolutional Neural Networks}
            \end{itemize}
            \item Kernel size differences? 25? 7? 5?
        \end{itemize}
        \item Training runs/results on:
        \begin{itemize}
            \item tracks-to-KDE run with Norm Chisq Loss Function
            \item UNet with "fixed" architecture
            \item End-to-end tracks-to-hist longer run
            \item Maybe redo POCA runs?
        \end{itemize}
    \end{itemize}
\end{frame}

\begin{frame}[plain]
    \centering
    Feb 7th Updates
\end{frame}

\begin{frame}{Overview of This Week's Work}
    \begin{itemize}
        \item Reduced GPU Availability in sneezy
        \item Chisq Loss Function KDE-A-z Training
        \item Ideas:
        \begin{itemize}
            \item What went wrong in prior end-to-end tracks-to-hist?
            \item Grokking?
            \item L2 Regularization?
        \end{itemize}
        \item Next Week: Group Meeting Tuesday (02/11) to Go Over APS Rehearsal Slides
        \item Next Week: Meeting Wednesday (02/12) with Simon and Mike
        \begin{itemize}
            \item Get clarifying ideas on where to head towards next
        \end{itemize}
    \end{itemize}
\end{frame}

\begin{frame}{Chisq Loss Function 500 Epochs}
    \begin{figure}
        \centering
        \includegraphics[width=0.6\linewidth]{feb7th/sample_chisq_output.png}
    \end{figure}
\end{frame}

\begin{frame}{LHCb DNN for End-to-end}
    \begin{figure}
        \centering
        \includegraphics[width=0.9\linewidth]{feb7th/lhcb_t2k.jpg}
    \end{figure}
\end{frame}

\begin{frame}{Ideas/Musings}
    \begin{itemize}
        \item (Perhaps) What went wrong?
        \begin{itemize}
            \item Bug: Accidentally switched output activation function from Softplus to LeakyRelu KDE-A-z
            \item Need to keep training the DNN - never plateaus (?)
        \end{itemize}
        \item Normalizing KDEs?
        \begin{itemize}
            \item LHCb divides the output KDE-A-z by 2500 so that the values are close to 0 and 2
            \item Prevents random gradient explosions
            \begin{itemize}
                \item Especially with large KDE-B-z values that reach up to \(\sim100,000\)
            \end{itemize}
            \item Adjusting loss function to reduce weight of KDE-B-z
            \begin{itemize}
                \item \( \Lagr_{\text{total}} = \alpha \Lagr_{\text{MSE, KDE-A-z}} + \beta \Lagr_{\text{MSE, KDE-B-z}} + \Lagr_{\text{MSE, KDE-A-x}} + \Lagr_{\text{MSE, KDE-A-y}} \)
                \item \( \alpha = 1e-4, \beta = 1e-8\)
            \end{itemize}
        \end{itemize}
        \item Grokking? (for KDE x and y) \href{https://arxiv.org/pdf/2201.02177}{Link}
        \begin{itemize}
            \item Paper discussing neural networks that generalize far after overfitting
            \item Maybe we just need to let the DNN run for a thousand more epochs (?)
            \item Uses a couple interesting techniques with weight decay, AdamW, drop out, etc
        \end{itemize}
        \item Lets have talk with LHCb then come back to this later!
    \end{itemize}
\end{frame}

\begin{frame}{Loss Plots KDE x and y}
    \begin{columns}
        \begin{column}{0.5\textwidth}
            \begin{figure}
                \centering
                \includegraphics[width=1\linewidth]{feb7th/lossKDE_A_x.png}
            \end{figure}
        \end{column}

        \begin{column}{0.5\textwidth}
            \begin{figure}
                \centering
                \includegraphics[width=1\linewidth]{feb7th/loss_KDE_A_y.png}
            \end{figure}
        \end{column}
    \end{columns}
\end{frame}

\begin{frame}{Grokking Paper}
    \begin{figure}
        \centering
        \includegraphics[width=1\linewidth]{feb7th/grokking_paper.png}
    \end{figure}
\end{frame}

\begin{frame}{Guiding Questions for Meeting with Mike and Simon}
    \begin{itemize}
        \item Data Processing
        \begin{itemize}
            \item How did you choose what to normalize KDE-A-z by? (value is 2500)
            \item What about KDE-B-z? (x and y are untouched)
        \end{itemize}
        \item tracks-to-KDE
        \begin{itemize}
            \item (Create slides giving an update of current progress)
            \item When training just tracks-to-KDE, convergence is slow and takes thousands of epochs, any recommendations to speed it up?
            \item Chisq Loss Function? K-L Divergence?
            \item For your end-to-end version, you output (batch size, 4,100) for the DNN section. Should I just train to that then connect them together?
        \end{itemize}
        \item KDE-to-hist
        \begin{itemize}
            \item Skip Connections
            \item Performance of UNet with increasing number of channels by 2?
            \item Both Batch Normalization and Dropout in UNet? (Variance Shifting Issue? \href{https://arxiv.org/abs/1801.05134}{Paper})
        \end{itemize}
    \end{itemize}
\end{frame}


\begin{frame}[plain]
    \centering
    Rehearsal APS Global Physics Summit 2025
\end{frame}

\begin{frame}{The ATLAS Experiment: Inner Detector}
    \begin{figure}
        \centering
        \includegraphics[width=0.8\linewidth]{APS2025/inner_detector_atlas.png}
        \label{fig:TheATLASITK}
    \end{figure}
\end{frame}

\begin{frame}{Track-level Information}
    \begin{columns}
        % Left column
        \column{0.3\linewidth}
        \begin{itemize}
            \item Gathering charged particle measurements readout from the detector
            \item Calculate trajectory estimation
            \item Used as inputs for reconstruction efforts and analysis
        \end{itemize}

        % Right column
        \column{0.7\linewidth}
        \begin{figure}
            \includegraphics[width=0.9\linewidth]{oct25th/coordinates_track.png}
            \caption{Reconstructed Track Data}
            \label{fig:coordinate_tracks}
        \end{figure}
    \end{columns}
\end{frame}

\begin{frame}{Current Primary Vertex Algorithm: AVMF}
    % Talk about innate challenges that comes with using AVMF
    \begin{columns}
        % Text
        \column{0.4\linewidth}
        \begin{itemize}
            \item Global approach to vertex finding and fitting \cite{sanderswood2019developmentatlasprimaryvertex}
            \item Finds and fits tracks based on compatibility to different primary vertices
            \item Limitations
            \begin{enumerate}
                \item Parallelization
                \item Time scale increases non linearly with number of PVs
            \end{enumerate}
        \end{itemize}

        % Text
        \column{0.6\linewidth}
        \begin{figure}
            \centering
            \includegraphics[width=0.9\linewidth]{oct25th/AVMF.png}
            \label{fig:AVMF}
            \cite{ATLAS:2019jmx}
        \end{figure}
    \end{columns}
\end{frame}

\begin{frame}{HL-LHC}
    \begin{columns}
        % Text
        \column{0.4\linewidth}
        \begin{itemize}
            \item Current vertexing algorithms are iterative by nature
            \item HL-LHC's order-of-magnitude increase in data exponentially extends processing time
            \item Motivations for ML
            \begin{itemize}
                \item Capability of learning complex functions
                \item Parallel training of events
                \item Integration of new variables
                \item Fine tuning for future upgrades/projects
            \end{itemize}
        \end{itemize}

        % HL-LHC Top-AntiTop Simultaion
        \column{0.6\linewidth}
        \begin{figure}
            \centering
            \includegraphics[width=0.9\linewidth]{oct25th/ATLAS_HLLHC_mu200_VP1_88vertices_angled_noID_4K_wLogo.png}
            \caption{\href{http://cds.cern.ch/record/2846411/}{Simulation of Top-Antitop Pair Production at HL-LHC}}
            \label{fig:HL-LHCSim}
        \end{figure}
    \end{columns}
\end{frame}

\section{PV-Finder/tracks to KDE}

\begin{frame}{PV-Finder Overview}
    \vspace{25pt}
    \begin{figure}
        \centering
        \includegraphics[width=1\linewidth]{APS2025/tracks_pv_hybrid_diagram.png}
        \label{fig:oldPVFinderArchitecture}
    \end{figure}
\end{frame}

\begin{frame}{Kernel Density Estimation (1)}
    \begin{columns}
        % KDE Text
        \column{0.5\linewidth}
        \begin{itemize}
            \item Track's Gaussian Probability Density \\
            \(\mathbb{P}(r) = \frac{1}{2\pi\sqrt{|\Sigma|}} \exp{\alpha}\) \\
            \hspace{-25pt}
            \(\alpha = -\frac{1}{2}\left((d-d_0),(z-z_0)\right)^T \Sigma^{-1} \left((d-d_0),(z-z_0)\right)\) \\
            KDE-A = \(\sum_{\text{tracks}} \mathbb{P}(r)\) \\
            KDE-B = \(\sum_{\text{tracks}} \mathbb{P}(r)^2\) \\
            XMax = x position with max \(\mathbb{P}(r)\) \\
            YMax = y position with max \(\mathbb{P}(r)\)
            \item KDEs are composed of 12,000 bins
            \begin{itemize}
                \item \(z_0 \in \left[-240\text{mm}, 240\text{mm}\right]\)
                \item Each 100 KDE bins \(\equiv\) 4mm z range
            \end{itemize}
        \end{itemize}

        % KDE Graphic
        \column{0.5\linewidth}
        \begin{figure}
            \centering
            \includegraphics[width=1\linewidth]{oct25th/KDE.png}
            \caption{From ATLAS Public Note \cite{ATLAS:2023hxv}}
            \label{fig:KDE}
        \end{figure}
    \end{columns}
\end{frame}

\begin{frame}{Kernel Density Estimation (2)}
    \begin{columns}
        % KDE Text
        \column{0.5\linewidth}
        \begin{itemize}
            \item Current analytical KDE production:
            \begin{itemize}
                \item Evaluates the Gaussian PDF for all    particle track
                \item \(\sim\)50 million evaluations per track per event
            \end{itemize}
            \item Motivations for machine learning:
            \begin{itemize}
                \item Bypasses need to evaluate over large grid of points
                \item Parallelization
            \end{itemize}
        \end{itemize}

        % KDE Graphic
        \column{0.5\linewidth}
        \begin{figure}
            \centering
            \includegraphics[width=1\linewidth]{oct25th/KDE.png}
            \caption{From ATLAS Public Note \cite{ATLAS:2023hxv}}
            \label{fig:KDE}
        \end{figure}
    \end{columns}
\end{frame}

\begin{frame}{PV-Finder Updated}
    \vspace{25pt}
    \begin{figure}
        \centering
        \includegraphics[width=1\linewidth]{APS2025/tracks_pv_end_to_end_diagram.png}
        \label{fig:oldPVFinderArchitecture}
    \end{figure}
\end{frame}

\begin{frame}{Tracks to KDE Architecture}
    \vspace{25pt}
    \begin{figure}
        \centering
        \includegraphics[width=1\linewidth]{oct25th/pipeline_trackstoKDE(2).png}
        \label{fig:pipeline_trackstokde}
    \end{figure}
\end{frame}

\begin{frame}{Example Predicted KDE Outputs (Work in Progress)}
    \begin{figure}
        \centering
        \includegraphics[width=0.6\linewidth]{APS2025/sample_tracks_KDE.png}
        \label{fig:example_KDE_output_2}
    \end{figure}
\end{frame}

\begin{frame}{KDE to hist (UNet) Architecture}
    \begin{figure}
        \centering
        \includegraphics[width=0.8\linewidth]{APS2025/UNet.pdf}
        \label{fig:pipeline_kdetohist}
    \end{figure}
\end{frame}

\begin{frame}{Example End-to-end KDE to Hists Output (Work in Progress)}
    \begin{columns}
        \column{0.4\linewidth}
        \begin{itemize}
            \item Training on predicted KDE-A
            \begin{itemize}
                \item Using prior outputs from tracks to KDE
                \item \(\sim 3\) days of training
                \item Only improvements from here with more information
            \end{itemize}
            \item Note: Will update with work within the coming weeks
        \end{itemize}

        \column{0.6\linewidth}
        \begin{figure}
            \centering
            \includegraphics[width=0.9\linewidth]{dec18th/ExamplePVFinderOutput4.png}
            \label{fig:example_Hists_output_1}
        \end{figure}
    \end{columns}
\end{frame}

\begin{frame}{Comparison with AVMF: Efficiency (Work in Progress)}
    \begin{columns}
        \column{0.4\linewidth}
        \begin{itemize}
            \item Trained using all KDEs produced from the combinatorial method
            \item Efficiency \(= \frac{\text{\# reconstructed vertices}}{\text{\# truth vertices}}\)
            \item Prior work noted in pubnote \cite{ATLAS:2023hxv}
            \item Note: Will also update with work within the coming weeks
        \end{itemize}

        \column{0.6\linewidth}
        \begin{figure}
            \centering
            \includegraphics[width=0.9\linewidth]{dec18th/current_results.png}
            \label{fig:KDEtoHistsResults}
        \end{figure}
    \end{columns}
\end{frame}

\section{Summary}

\begin{frame}{Summary}
    \begin{itemize}
        \item Machine learning provides a viable avenue for primary vertex identificaiton
        \item Proof of concept of hybrid tracks-to-hist demonstrated in pubnote \cite{ATLAS:2023hxv}
        \item NN kernel generation has the potential to decrease false positive rates
        \item Future Plans:
        \begin{itemize}
            \item Continue optimizing the NN architecture for kernel generation
            \item Integration into ATHENA
        \end{itemize}
    \end{itemize}

    \centering
    \vspace{15pt}
    \textit{Thanks for listening!}
\end{frame}

\begin{frame}{Relevant Links and Directories}
    \begin{itemize}
        \item Ananya's prior work on tracks to KDE
        \begin{itemize}
            \item \href{https://github.com/qibinlei/TracksToKDE/tree/main}{Link to Gitlab}
            \item \href{https://docs.google.com/document/d/1oTlUo06nhkKBF059y9mCXGqZ46PPD3TuHFIe3EkTN-g/edit}{Link to Google Doc Documentation}
        \end{itemize}
        \item LHCb PVFinder
        \begin{itemize}
            \item \href{https://gitlab.cern.ch/sakar/pv-finder_v2/-/tree/master}{Link to Gitlab}
        \end{itemize}
        \item ATLAS PV Finder
        \begin{itemize}
            \item \href{https://gitlab.cern.ch/qlei/atlas_pvfinder}{Link to Gitlab}
        \end{itemize}
    \end{itemize}
\end{frame}

\section{Bibliography}
\begin{frame}[allowframebreaks]{Bibliography}
    \nocite{*}
    \printbibliography
\end{frame}

%%%%% Backup %%%%%%
\section*{Backup}
\nonumsec
% \backupbegin
% \begin{frame}{Network Structure Readout}
% 	\begin{wrapfig}{r}{0.5\textwidth}
% 	    \centering
% 	    \includegraphics[width=.6\textwidth]{graphs/DenseLayers.png}
%         \caption*{Dense Layers}
% 	\end{wrapfig}
% \end{frame}

\begin{frame}{Abstract}
    Extensive research has been conducted on deep neural networks (DNNs) for the identification and localization of primary vertices (PVs) in proton-proton collision data from ATLAS. Previous studies focused on locating primary vertices in simulated ATLAS data using a hybrid methodology. This approach began with the derivation of kernel density estimators (KDEs) from the ensemble of charged track parameters, employing an analytical probability density estimation technique. These KDEs were subsequently utilized as input for two neural network (NN) architectures, namely UNet and UNet++, alongside the truth PV positions extracted in the form of target histograms from the simulated data. Through these investigations, a proof-of-concept was demonstrated, achieving performance comparable to the ATLAS Adaptive Multi-Vertex Finder (AMVF) algorithm, while also enhancing the vertex position resolution.

The current studies transition from analytical KDE computation to a fully NN-based implementation, presenting an end-to-end primary vertex finder algorithm driven by neural networks. A comprehensive analysis of this approach, including a comparative assessment of its performance against the AMVF algorithm, will be presented.
\end{frame}

\begin{frame}{Data Statistics}
    \begin{itemize}
        \item Hard scatter simulation of \(t \bar{t}\) decaying semi-leptonically
        \item Total Events: 51000
        \item Input Features: \(d_0, z_0, \sigma(d_0), \sigma(z_0), \sigma(d_0, z_0)\) (\(\sim 1000\) tracks per event)
        \begin{itemize}
            \item \(z_0 \in [-240\text{mm}, 240\text{mm}]\)
        \end{itemize}
        \item Output Label: KDE-A (12,000 Bins)
    \end{itemize}
\end{frame}

\begin{frame}{Data Pre-Processing}
    \vspace{15pt}
    \begin{figure}
        \centering
        \includegraphics[width=1\linewidth]{oct25th/data_preprocessing_pvfinder(1).png}
        \label{fig:data_preprocessing_pvfinder}
    \end{figure}
\end{frame}

\begin{frame}{Masking}
    \begin{columns}
        \column{0.5\linewidth}
        \begin{itemize}
            \item A mask is used to identify valid tracks that should contribute to the final result (not padded value)
            \item The masking matrix, f2, is then multiplied with the output of the neural network to either
            \begin{enumerate}
                \item Zero out outputs of the neural network that do not have any track data associated with it
                \item Contribute values that will be summed to produce the predicted KDE bins
            \end{enumerate}
        \end{itemize}

        % Masking graphic
        \column{0.5\linewidth}
        \begin{figure}
            \centering
            \includegraphics[width=0.5\linewidth]{oct25th/creating_mask.png}
            \label{fig:masking_figure}
        \end{figure}
    \end{columns}
\end{frame}

\begin{frame}{Masking Output}
    \vspace{15pt}
    \begin{figure}
        \centering
        \includegraphics[width=0.8\linewidth]{oct25th/calculating_output(2).png}
        \label{fig:maskingtooutput}
    \end{figure}
\end{frame}

\begin{frame}{More Technical Details of DNN}
    \begin{figure}
        \centering
        \includegraphics[width=0.8\linewidth]{APS2025/DNN_specifications.png}
    \end{figure}
\end{frame}

\begin{frame}{More Technical Details of UNet}
    \begin{figure}
        \centering
        \includegraphics[width=0.8\linewidth]{APS2025/UNet_specifications.png}
    \end{figure}
\end{frame}

\begin{frame}{Loss Functions for Tracks to KDE}
    \begin{itemize}
        \item MSE:
        \begin{align*}
            \Lagr_{\text{MSE}} \left( y_{\text{pred}}, y_{\text{true}} \right) = \frac{1}{n}\sum \left(y_{\text{true}} - y_{\text{pred}} \right)^2
        \end{align*}
        \item Normalized Chisq (MSE-Like Loss Function):
        \begin{align*}
            \sigma &= 0.01 \\
            \chi^2 &= \left( \frac{y_{\text{pred}} - y_{\text{truth}}}{\sigma} \right)^2 \\
            \Lagr_{\text{Chisq}} = \chi^2_{\text{norm}} &= \frac{\chi^2}{\text{nBins} * \text{nBatch}}
        \end{align*}
        \item K-L Divergence
        \begin{align*}
            \Lagr_{\text{KR}} \left( y_{\text{pred}}, y_{\text{true}} \right) = \frac{1}{n}\sum y_{\text{true}} \cdot \log(\frac{y_{\text{true}}}{y_{\text{pred}}})
        \end{align*}
    \end{itemize}
\end{frame}

\begin{frame}{Loss Functions for UNet}
    \begin{align*}
        \Lagr_{\text{UNet}} &= - \sum_{i=1}^{12000} \ln(z_i) \\
        z_i &= \frac{2r_i}{r_i + r_i^{-1}} \\
        r_i &= \frac{t_i + \epsilon}{p_i + \epsilon}
    \end{align*}

    \begin{itemize}
        \item \(t_i\) is the value at bin i of the target histogram
        \item \(p_i\) is the value at bin i of the predicted histogram
        \item \(\epsilon\) is used to handle edge cases for denominator
    \end{itemize}
\end{frame}

% \section{Bibliography}
% \begin{frame}[allowframebreaks]{Bibliography}
%     \nocite{*}
%     \printbibliography
% \end{frame}

%%%%%% Backup %%%%%%
% \section*{Backup}
% \nonumsec
% \backupbegin
% % \begin{frame}{Network Structure Readout}
% % 	\begin{wrapfig}{r}{0.5\textwidth}
% % 	    \centering
% % 	    \includegraphics[width=.6\textwidth]{graphs/DenseLayers.png}
% %         \caption*{Dense Layers}
% % 	\end{wrapfig}
% % \end{frame}

\begin{frame}[plain]
    \centering
    Feb 14th Updates
\end{frame}

\begin{frame}{Overview of This Week's Work}
    \begin{itemize}
        \item Group Meeting Tuesday (02/10) Rehearsal\(^2\) for APS Slides \checkmark
        \item ATLAS Internal Rehearsal (02/12) \checkmark
        \begin{itemize}
            \item Will update feedback! Really focusing on getting NN results done before anything else
        \end{itemize}
        \item Meeting with Simon and Mike (02/12) \checkmark
        \begin{itemize}
            \item Went well! Was able to clarify a lot of points of confusion I had and gives me a roadmap for upcoming work
            \item Frankenstein-ing DNN and UNet
            \item Split Analytical KDE to Split Hist Study
        \end{itemize}
    \end{itemize}
\end{frame}

\begin{enumerate}
    \item Initialize weights for the DNN using KDE-A-z trainings
    \item In the last layer of the DNN, use existing 100 nodes from prior output layer + create 700 extra nodes
    \begin{itemize}
        \item Using: torch.randn(100, 700) * 0.02
        \item New DNN Outputs: (8,100)
    \end{itemize}
    \item Initialize weights for UNet using prior trainings
    \item Connect DNN to UNet
\end{enumerate}
\begin{itemize}
    \item Note: Don't need to train KDE-B-z, KDE-A-x, KDE-A-y anymore
\end{itemize}

\begin{frame}{Frankensteining DNN and UNet (1)}
    \begin{figure}
        \centering
        \includegraphics[width=1\linewidth]{feb14th/frankenstein1.png}
        \label{fig:frankenstein1}
    \end{figure}
\end{frame}

\begin{frame}{Frankensteining DNN and UNet (2)}
    \begin{figure}
        \centering
        \includegraphics[width=1\linewidth]{feb14th/frankenstein2.png}
        \label{fig:frankenstein2}
    \end{figure}
\end{frame}

\begin{frame}{Frankensteining DNN and UNet (3)}
    \begin{figure}
        \centering
        \includegraphics[width=1\linewidth]{feb14th/frankenstein3.png}
        \label{fig:frankenstein3}
    \end{figure}
\end{frame}

\begin{frame}{Preliminary Run (Will Let Run Over Weekend)}
    \begin{figure}
        \centering
        \includegraphics[width=0.6\linewidth]{feb14th/prelim_endtoend.png}
    \end{figure}
\end{frame}

\begin{frame}[plain]
    \centering
    Feb 21st Updates
\end{frame}

\begin{frame}{Overview of This Week’s Work}
    \begin{itemize}
        \item Fix APS Slides
        \item Studies
        \begin{itemize}
            \item End-to-end UNet
            \item End-to-end UNet Double
            \item False positive rate
        \end{itemize}
        \item \(\Delta z_{\text{vtx-vtx}}\) issues?
        \item Make end-to-end result plots
    \end{itemize}
\end{frame}

\begin{frame}{UNet}
    \begin{itemize}
        \item Does suprisingly amazingly well, hits about an 82\% efficiency
    \end{itemize}
    \begin{figure}
        \centering
        \includegraphics[width=0.8\linewidth]{feb20th/unetloss.png}
    \end{figure}
\end{frame}

\begin{frame}{UNet Double}
    \begin{itemize}
        \item Immitated the original UNet with tons of inbetween CNN layers
        \item Overfits but is able to learn the training set really well
    \end{itemize}
    \begin{figure}
        \centering
        \includegraphics[width=0.8\linewidth]{feb20th/unetdoubleloss.png}
    \end{figure}
\end{frame}

\begin{frame}{\(\Delta z_{\text{vtx-vtx}}\) issues?}
    \begin{columns}
        \begin{column}{0.5\linewidth}
            \begin{figure}
                \centering
                \includegraphics[width=1\linewidth]{feb20th/amvf_pv_distances.png}
            \end{figure}
        \end{column}

        \begin{column}{0.5\linewidth}
            \begin{figure}
                \centering
                \includegraphics[width=1\linewidth]{feb20th/pred_t2h_pv_distances.png}
            \end{figure}
        \end{column}
    \end{columns}
\end{frame}

\begin{frame}{False Positive Rate (1)}
    \begin{itemize}
        \item These tails in the PV outputs might be why the \(\Delta z\) issues are happening???
    \end{itemize}
    \begin{figure}
        \centering
        \includegraphics[width=0.5\linewidth]{feb20th/problems_unet.png}
        \label{fig:tails_unet}
    \end{figure}
\end{frame}

\begin{frame}{False Positive Rate (2)}
    \begin{itemize}
        \item Maybe having larger interval extensions for tracks would help? (\(\pm 2.5\) mm)
        \item Ran a quick run using larger interval extensions and got...
    \end{itemize}
    \begin{figure}
        \centering
        \includegraphics[width=0.5\linewidth]{feb20th/problems_constint_unet.png}
        \label{fig:tails_unet}
    \end{figure}
\end{frame}

\begin{frame}{Preliminary Result Plot: Compare Clean}
    \begin{figure}
        \centering
        \includegraphics[width=0.6\linewidth]{feb20th/npv_compare_Clean.png}
    \end{figure}
\end{frame}

\begin{frame}{Preliminary Result Plot: Total}
    \begin{figure}
        \centering
        \includegraphics[width=0.6\linewidth]{feb20th/npv_compare_Total.png}
    \end{figure}
\end{frame}

\begin{frame}{Preliminary Result Plot: Avg Per Category}
    \begin{figure}
        \centering
        \includegraphics[width=0.5\linewidth]{feb20th/avg_per_category.png}
    \end{figure}
\end{frame}

\begin{frame}{Overview of This Week's Work}
    \begin{itemize}
        \item \(\cancel{\text{Fix APS Slides}}\)
        \item Preliminary Results/Plots
        \item Upcoming Work/Summary
    \end{itemize}
\end{frame}

\begin{frame}{Approximate Training Times}
    \begin{itemize}
        \item \textbf{To Do: Get an time estimate of per event evaluation}
        \item 100 Bins + \(\theta + \phi\)
        \begin{itemize}
            \item tracks-to-kde: 3 minutes per epoch for 200 epochs (\(\sim\) 10 hours)
            \item KDE-to-hist: 3 minutes per epoch for 50 epochs (\(\sim\) 2.5 hours)
            \item tracks-to-hist: 5 minutes per epoch for 100 epochs (\(\sim\) 8 hours)
            \item \textbf{Total: 21 hours of training time}
        \end{itemize}
        \item 400 Bins + \(\theta + \phi\)
        \begin{itemize}
            \item tracks-to-kde: 68 seconds per epoch for 150 epochs (\(\sim\) 3 hours)
            \item KDE-to-hist: 74 seconds per epoch for 50 epochs (\(\sim\) 1 hour)
            \item tracks-to-hist: 155 seconds per epoch for 50 epochs (\(\sim\) 2 hours)
            \item \textbf{Total: 6 hours of training time (less than a work day!)}
        \end{itemize}
    \end{itemize}
\end{frame}

\begin{frame}{Results of Preliminary Runs}
    \centering
    \begin{tabular}{c c c}
        \hline
        Model & Efficiency & False Positive Rate (\# / event) \\
        \hline
        T2H UNet 400 Bins & 90.4\% & 0.97 \\
        T2H UNet 400 Bins + \(\theta + \phi\) & 91.4\% & 0.96 \\
        T2H UNet 400 Bins + \(\theta + \phi\) + 8 Latent & 90.9\% & 0.81 \\
        \textbf{T2H UNet 100 Bins + \(\theta + \phi\)} & \textbf{91.8\%} & \textbf{0.73} \\
        \textbf{T2H UNet 1000 Bins + 8 Latent} & \textbf{92.8\%} & \textbf{0.81} \\
        AMVF & 93.9\% & 0.8 \\
        \hline
    \end{tabular}

    \vspace{7pt}

    \begin{tabular}{c c}
        \hline
        Method & \(\sigma_{\text{vtx-vtx}}\) (mm) \\
        \hline
        T2H UNet 400 bins & \(0.29 \pm 0.03\) \\
        T2H UNet 400 bins + \(\theta + \phi\) & \(0.30 \pm 0.02\) \\
        T2H UNet 400 Bins + \(\theta + \phi\) + 8 Latent & \(0.31 \pm 0.01\) \\
        T2H UNet 100 Bins + \(\theta + \phi\) & \(0.33 \pm 0.01\) \\
        T2H UNet 1000 Bins + 8 Latent & \(0.37 \pm 0.01\) \\
        AMVF & \(0.76 \pm 0.02\) \\
        \hline
    \end{tabular}
\end{frame}

\begin{frame}{Upcoming Work/Summary (1)}
    \begin{itemize}
        \item Upcoming Runs:
        \begin{itemize}
            \item UNet++ Runs
            \item Hyperparameter Tuning (???)
            \begin{itemize}
                \item \href{https://github.com/google-research/tuning_playbook}{Google Deep Learning Tuning Playbook}
                \item Can probably squeeze out 2-3\% Efficiency and 0.1-0.2 lower FPR
                \item Possibly make faster with a larger batch size or output size or both
            \end{itemize}
        \end{itemize}
        \item Start Preliminary Reading/Work on GNN
        \begin{itemize}
            \item \href{https://distill.pub/2021/gnn-intro/}{Nice intro reading!}
            \item \href{https://indico.nikhef.nl/event/4875/contributions/20317/attachments/8200/11960/PVFinder_EuCAIFCon24.pdf}{LHCb's Graph Neural Network (Poster)}
        \end{itemize}
        \item Start Muon Collider Research (woo!)
        \begin{itemize}
            \item Created Nersc account and started reading through documentation
        \end{itemize}
    \end{itemize}
\end{frame}

\begin{frame}{Upcoming Work/Summary (2)}
    \begin{itemize}
        \item Upcoming Presentations:
        \begin{itemize}
            \item \(\cancel{\text{Rehearsal\(^3\) March 4th}}\)
            \item TVPD Update (March ???)
            \item APS Joint Meeting 2025 (March 18th)
            \item Poster (Start assembling graphics)
        \end{itemize}
        \item Misc:
        \begin{itemize}
            \item Doing a complete audit of code to ensure reproducability and correctness
            \item Push to GitLab
        \end{itemize}
    \end{itemize}
\end{frame}

\begin{frame}[plain]
    \centering
    FINAL SLIDES APS Global Physics Summit 2025
\end{frame}

\section{Introduction}

% \begin{frame}[plain]{Outline}
% 	\tableofcontents
% \end{frame}

\begin{frame}{The ATLAS Experiment: Inner Detector}
    \begin{figure}
        \centering
        \includegraphics[width=0.8\linewidth]{APS2025/inner_detector_atlas.png}
        \label{fig:TheATLASITK}
    \end{figure}
\end{frame}

\begin{frame}{Track-level Information}
    \begin{columns}
        % Left column
        \column{0.3\linewidth}
        \begin{itemize}
            \item Tracks are within a magnetic field which creates a helical path
            \item \(d_0\) and \(z_0\) are position parameters in the transverse and longitudinal plane
        \end{itemize}

        % Right column
        \column{0.7\linewidth}
        \begin{figure}
            \includegraphics[width=0.9\linewidth]{oct25th/coordinates_track.png}
            \caption{Reconstructed Track Data}
            \label{fig:coordinate_tracks}
        \end{figure}
    \end{columns}
\end{frame}

\begin{frame}{Current Primary Vertex Algorithm: AMVF}
    % Talk about innate challenges that comes with using AVMF
    \begin{columns}
        % Text
        \column{0.4\linewidth}
        \begin{itemize}
            \item Global approach to vertex finding and fitting
            \item Finds and fits tracks based on compatibility to different primary vertices
            \item Limitations
            \begin{enumerate}
                \item Parallelization
                \item Time scale increases non linearly with number of PVs
            \end{enumerate}
        \end{itemize}

        % Text
        \column{0.6\linewidth}
        \begin{figure}
            \centering
            \includegraphics[width=0.9\linewidth]{oct25th/AVMF.png}
            \label{fig:AVMF}
                \href{http://cds.cern.ch/record/2670380/files/ATL-PHYS-PUB-2019-015.pdf?version=1}{[ATL-PHYS-PUB-2019-015]}
        \end{figure}
    \end{columns}
\end{frame}

\begin{frame}{The Large Hadron Collider}
    \begin{figure}
        \centering
        \includegraphics[width=1\linewidth]{APS2025/hllhc_timeline_lum_crosS_3.png}
        \label{fig:LHC_plans}
    \end{figure}
\end{frame}

\begin{frame}{The Large Hadron Collider Event Display}
    \hspace{10pt}
    \begin{figure}
        \centering
        \includegraphics[width=1\linewidth]{APS2025/HLLHC_event_display_ex.png}
        \label{fig:event_display_comparison}
    \end{figure}
\end{frame}

% \begin{columns}
%     % Text
%     \column{0.4\linewidth}
%     \begin{itemize}
%         \item Current vertexing algorithms are iterative by nature
%         \item HL-LHC's order-of-magnitude increase in data exponentially extends processing time
%         \item Motivations for ML
%         \begin{itemize}
%             \item Capability of learning complex functions
%             \item Parallel training of events
%             \item Integration of new variables
%             \item Fine tuning for future upgrades/projects
%         \end{itemize}
%     \end{itemize}

%     % HL-LHC Top-AntiTop Simultaion
%     \column{0.6\linewidth}
%     \begin{figure}
%         \centering
%         \includegraphics[width=0.9\linewidth]{oct25th/ATLAS_HLLHC_mu200_VP1_88vertices_angled_noID_4K_wLogo.png}
%         \caption{\href{http://cds.cern.ch/record/2846411/}{Simulation of Top-Antitop Pair Production at HL-LHC}}
%         \label{fig:HL-LHCSim}
%     \end{figure}
% \end{columns}

\section{PV-Finder/tracks to KDE}

\begin{frame}{PV-Finder Overview}
    \vspace{25pt}
    \begin{figure}
        \centering
        \includegraphics[width=1\linewidth]{APS2025/analyticalkde_pv.png}
        \label{fig:oldPVFinderArchitecture}
    \end{figure}
\end{frame}

\begin{frame}{Kernel Density Estimation}
    \begin{columns}
        % KDE Text
        \column{0.5\linewidth}
        \begin{itemize}
            \item Track's Gaussian Probability Density \\
            \(\mathbb{P}(r) = \frac{1}{2\pi\sqrt{|\Sigma|}} \exp{\alpha}\) \\
            \hspace{-25pt}
            \(\alpha = -\frac{1}{2}\left((d-d_0),(z-z_0)\right)^T \Sigma^{-1} \left((d-d_0),(z-z_0)\right)\) \\
            KDE-A = \(\sum_{\text{tracks}} \mathbb{P}(r)\) \\
            KDE-B = \(\sum_{\text{tracks}} \mathbb{P}(r)^2\)
            \item KDEs are composed of 12,000 bins
            \begin{itemize}
                \item \(z_0 \in \left[-240\text{mm}, 240\text{mm}\right]\)
                \item Each 100 KDE bins \(\equiv\) 4mm z range
            \end{itemize}
            \item Analytical production requires \textbf{\(\sim\)50 million evaluations} per track per event
        \end{itemize}

        % KDE Graphic
        \column{0.5\linewidth}
        \begin{figure}
            \centering
            \includegraphics[width=1\linewidth]{APS2025/KDE_example_z_axis.png}
            \href{https://cds.cern.ch/record/2858348/files/ATL-PHYS-PUB-2023-011.pdf?version=1}{[ATL-PHYS-PUB-2023-011]}
            \label{fig:KDE}
        \end{figure}
    \end{columns}
\end{frame}

\begin{frame}{Prior PV-Finder Results}
    \begin{columns}
        \begin{column}{0.4\linewidth}
            \begin{itemize}
                \item \textbf{Efficiency:} fraction of reconstructable truth vertices that are correctly assigned to reconstructed vertices
                \item \textbf{False Positive Rate:} fraction of reconstructed primary vertices that do not match truth vertices \\
            \href{https://cds.cern.ch/record/2858348/files/ATL-PHYS-PUB-2023-011.pdf?version=1}{[ATL-PHYS-PUB-2023-011]}
            \end{itemize}
        \end{column}

        \begin{column}{0.6\linewidth}
            \begin{figure}
                \centering
                \includegraphics[width=0.8\linewidth]{APS2025/highres_current_results.png}
                \label{fig:KDEtoHistsResults}
            \end{figure}
        \end{column}
    \end{columns}
\end{frame}

% \begin{frame}[noframenumbering]{Kernel Density Estimation}
%     \begin{columns}
%         % KDE Text
%         \column{0.5\linewidth}
%         \begin{itemize}
%             \item Current analytical KDE production:
%             \begin{itemize}
%                 \item Evaluates the Gaussian PDF for all    particle track
%                 \item \(\sim\)50 million evaluations per track per event
%             \end{itemize}
%             \item Motivations for machine learning:
%             \begin{itemize}
%                 \item Bypasses need to evaluate over large grid of points
%                 \item Group tracks that lie within a 100 KDE bin range
%                 \item Parallelization
%             \end{itemize}
%         \end{itemize}

%         % KDE Graphic
%         \column{0.5\linewidth}
%         \begin{figure}
%             \centering
%             \includegraphics[width=1\linewidth]{oct25th/KDE.png}
%             \href{https://cds.cern.ch/record/2858348/files/ATL-PHYS-PUB-2023-011.pdf?version=1}{[ATL-PHYS-PUB-2023-011]}
%             \label{fig:KDE}
%         \end{figure}
%     \end{columns}
% \end{frame}

\begin{frame}{PV-Finder Before}
    \vspace{5pt}
    \centering
    This network performs nicely but is limited by execution time
    \vspace{25pt}
    \begin{figure}
        \centering
        \includegraphics[width=1\linewidth]{APS2025/analyticalkde_pv.png}
        \label{fig:oldPVFinderArchitecture}
    \end{figure}
\end{frame}

\begin{frame}[noframenumbering]{PV-Finder After}
    \vspace{5pt}
    \centering
    A DNN that can learn the KDE mapping can help speed up the process
    \vspace{25pt}
    \begin{figure}
        \centering
        \includegraphics[width=1\linewidth]{APS2025/end_to_end_trackspv.png}
        \label{fig:oldPVFinderArchitecture}
    \end{figure}
\end{frame}

\begin{frame}{Example Predicted KDE Outputs}
    \begin{figure}
        \centering
        \includegraphics[width=0.6\linewidth]{APS2025/KDE_A_pred_example_workinprogress.png}
        \label{fig:example_KDE_output_2}
    \end{figure}
\end{frame}

\begin{frame}{Example End-to-end tracks-to-hist Output}
    \begin{columns}
        \column{0.3\linewidth}
        \begin{itemize}
            \item Initialize weights to be from separate trainings of the DNN and UNet
            \item \(\sim 3\) hours of training
        \end{itemize}

        \column{0.7\linewidth}
        \begin{figure}
            \centering
            \includegraphics[width=0.85\linewidth]{APS2025/embed_pred_KDE_A_pv_pred_pv_final.png}
            \label{fig:example_Hists_output_1}
        \end{figure}
    \end{columns}
\end{frame}

\begin{frame}{Comparison with AMVF: Efficiency}
    \begin{columns}
        \column{0.3\linewidth}
        \begin{itemize}
            \item DNN helps drastically reduce the false positive rates in PV-Finder architecture
            \item Achieves comparable efficiency and false positive rates to AMVF
            \item Much faster!
        \end{itemize}

        \column{0.7\linewidth}
        \begin{figure}
            \centering
            \includegraphics[width=0.9\linewidth]{APS2025/ntrks_fpr_eff_fixed_2.png}
            \label{fig:efffprresults}
        \end{figure}
    \end{columns}
\end{frame}

\section{Summary}

\begin{frame}{Summary}
    \begin{itemize}
        \item Machine learning provides a viable avenue for primary vertex identificaiton
        \item Proof of concept of end-to-end tracks-to-hist has been demonstrated
        \begin{itemize}
            \item NN kernel generation improves efficiency and decreases the false positive rate
        \end{itemize}
        \item Future Plans:
        \begin{itemize}
            \item Continue optimizing the end-to-end architecture for better performance
            \item Test architecture on HL-LHC simulated conditions
            \item Exploration of other architectures (e.g. GNN)
            \item Integration into ATHENA
        \end{itemize}
    \end{itemize}

    \centering
    \vspace{15pt}
    \textit{Thanks for listening!} \\
    \vspace{10pt}
    \textit{Special thanks to IRIS-HEP and Simon Akar/Mike Sokoloff from LHCb for their guidance!}
\end{frame}

\begin{frame}{Relevant Links and Directories}
    \begin{itemize}
        \item Ananya's prior work on tracks to KDE
        \begin{itemize}
            \item \href{https://github.com/qibinlei/TracksToKDE/tree/main}{Link to Gitlab}
            \item \href{https://docs.google.com/document/d/1oTlUo06nhkKBF059y9mCXGqZ46PPD3TuHFIe3EkTN-g/edit}{Link to Google Doc Documentation}
        \end{itemize}
        \item LHCb PVFinder
        \begin{itemize}
            \item \href{https://gitlab.cern.ch/sakar/pv-finder_v2/-/tree/master}{Link to Gitlab}
        \end{itemize}
        \item ATLAS PV Finder
        \begin{itemize}
            \item \href{https://gitlab.cern.ch/qlei/atlas_pvfinder}{Link to Gitlab}
        \end{itemize}
    \end{itemize}
\end{frame}

\begin{frame}{Abstract}
    Extensive research has been conducted on deep neural networks (DNNs) for the identification and localization of primary vertices (PVs) in proton-proton collision data from ATLAS. Previous studies focused on locating primary vertices in simulated ATLAS data using a hybrid methodology. This approach began with the derivation of kernel density estimators (KDEs) from the ensemble of charged track parameters, employing an analytical probability density estimation technique. These KDEs were subsequently utilized as input for two neural network (NN) architectures, namely UNet and UNet++, alongside the truth PV positions extracted in the form of target histograms from the simulated data. Through these investigations, a proof-of-concept was demonstrated, achieving performance comparable to the ATLAS Adaptive Multi-Vertex Finder (AMVF) algorithm, while also enhancing the vertex position resolution.

The current studies transition from analytical KDE computation to a fully NN-based implementation, presenting an end-to-end primary vertex finder algorithm driven by neural networks. A comprehensive analysis of this approach, including a comparative assessment of its performance against the AMVF algorithm, will be presented.
\end{frame}

\begin{frame}{Other Track Parameters}
    \centering
    \begin{align*}
        \text{Track Parameters: }\vec{x} = \left( d_0, z_0, \phi, \theta, \frac{q}{p}, t \right)
    \end{align*}
    \vspace{15pt}
    \begin{figure}
        \centering
        \includegraphics[width=0.8\linewidth]{APS2025/parameters.png}
        \label{fig:trackparameters_other}
        \caption{ACTS Docs - Tracking in a nutshell}
    \end{figure}
\end{frame}

\begin{frame}{Data Statistics}
    \begin{itemize}
        \item Full detector hard scatter simulation of \(t \bar{t}\) decaying semi-leptonically with \(\left< \mu\right> = 60\)
        \item Total Events: 51000
        \item Input Features: \(d_0, z_0, \sigma(d_0), \sigma(z_0), \sigma(d_0, z_0)\) (\(\sim 1000\) tracks per event)
        \begin{itemize}
            \item \(z_0 \in [-240\text{mm}, 240\text{mm}]\)
        \end{itemize}
        \item Output Label: KDE-A (12,000 Bins)
        \item Full Noted In: \href{https://cds.cern.ch/record/2858348/files/ATL-PHYS-PUB-2023-011.pdf?version=1}{ATL-PHYS-PUB-2023-011}
    \end{itemize}
\end{frame}

\begin{frame}{Data Pre-Processing}
    \vspace{15pt}
    \begin{figure}
        \centering
        \includegraphics[width=1\linewidth]{oct25th/data_preprocessing_pvfinder(1).png}
        \label{fig:data_preprocessing_pvfinder}
    \end{figure}
\end{frame}

\begin{frame}{Masking}
    \begin{columns}
        \column{0.5\linewidth}
        \begin{itemize}
            \item A mask is used to identify valid tracks that should contribute to the final result (not padded value)
            \item The masking matrix, f2, is then multiplied with the output of the neural network to either
            \begin{enumerate}
                \item Zero out outputs of the neural network that do not have any track data associated with it
                \item Contribute values that will be summed to produce the predicted KDE bins
            \end{enumerate}
        \end{itemize}

        % Masking graphic
        \column{0.5\linewidth}
        \begin{figure}
            \centering
            \includegraphics[width=0.5\linewidth]{oct25th/creating_mask.png}
            \label{fig:masking_figure}
        \end{figure}
    \end{columns}
\end{frame}

\begin{frame}{Masking Output}
    \vspace{15pt}
    \begin{figure}
        \centering
        \includegraphics[width=0.8\linewidth]{oct25th/calculating_output(2).png}
        \label{fig:maskingtooutput}
    \end{figure}
\end{frame}

% \begin{frame}{More Technical Details of DNN}
%     \begin{figure}
%         \centering
%         \includegraphics[width=0.8\linewidth]{APS2025/DNN_specifications.png}
%     \end{figure}
% \end{frame}

\begin{frame}{Tracks to KDE Architecture}
    \vspace{25pt}
    \begin{figure}
        \centering
        \includegraphics[width=1\linewidth]{APS2025/trackstoKDE_simplified_400.png}
        \label{fig:pipeline_trackstokde}
    \end{figure}
\end{frame}

\begin{frame}{KDE to hist (UNet) Architecture}
    \begin{columns}
        \begin{column}{0.4\linewidth}
            \begin{itemize}
                \item Effective for feature extraction
                \item Downsampling reduces dimensionality, helping emphasize important data features
                \item Skip connections propogate prior information
            \end{itemize}
        \end{column}

        \begin{column}{0.6\linewidth}
            \begin{figure}
                \centering
                \includegraphics[width=1\linewidth]{APS2025/UNet.pdf}
                \label{fig:pipeline_kdetohist}
            \end{figure}
        \end{column}
    \end{columns}
\end{frame}

\begin{frame}{More Technical Details of Neural Networks}
    \begin{figure}
        \centering
        \includegraphics[width=0.8\linewidth]{APS2025/hyperparameters_unet.png}
    \end{figure}
\end{frame}

\begin{frame}{Loss Function for Tracks to KDE}
    \begin{itemize}
        \item MSE:
        \begin{align*}
            \Lagr_{\text{MSE}} \left( y_{\text{pred}}, y_{\text{true}} \right) = \frac{1}{n}\sum \left(y_{\text{true}} - y_{\text{pred}} \right)^2
        \end{align*}
        \item Normalized Chisq (MSE-Like Loss Function) (Used by LHCb):
        \begin{align*}
            \sigma &= 0.01 \\
            \chi^2 &= \left( \frac{y_{\text{pred}} - y_{\text{truth}}}{\sigma} \right)^2 \\
            \Lagr_{\text{Chisq}} = \chi^2_{\text{norm}} &= \frac{\chi^2}{\text{nBins} * \text{nBatch}}
        \end{align*}
    \end{itemize}
\end{frame}

\begin{frame}{Loss Functions for UNet}
    \begin{align*}
        \Lagr_{\text{UNet}} &= - \sum_{i=1}^{12000} \ln(z_i) \\
        z_i &= \frac{2r_i}{r_i + r_i^{-1}} \\
        r_i &= \frac{t_i + \epsilon}{p_i + \epsilon}
    \end{align*}

    \begin{itemize}
        \item \(t_i\) is the value at bin i of the target histogram
        \item \(p_i\) is the value at bin i of the predicted histogram
        \item \(\epsilon\) is used to handle edge cases for denominator
    \end{itemize}
\end{frame}

\begin{frame}{\(\Delta z\) Resolution and \(\sigma_{\text{vtx-vtx}}\) Calculation}
    \begin{columns}
        \begin{column}{0.6\linewidth}
            \begin{figure}
                \centering
                \includegraphics[width=0.9\linewidth]{APS2025/deltaz(1).png}
            \end{figure}
        \end{column}

        \begin{column}{0.4\linewidth}
            \begin{figure}
                \centering
                \includegraphics[width=1\linewidth]{APS2025/methods_sigma_vtx_vtx2.png}
            \end{figure}
        \end{column}
    \end{columns}
\end{frame}

\begin{frame}{Comparison with AVMF: Vertex Classification Scheme}
    \begin{figure}
        \centering
        \includegraphics[width=0.75\linewidth]{APS2025/vertex_classification_aps2025.png}
    \end{figure}
\end{frame}

\begin{frame}{Number of Primary Vertives vs \(\mu\)}
    \begin{figure}
        \centering
        \includegraphics[width=0.7\linewidth]{APS2025/npv_compare_Total.png}
    \end{figure}
\end{frame}

\begin{frame}{Vertex Resolution vs Tracks}
    \begin{figure}
        \centering
        \includegraphics[width=0.7\linewidth]{APS2025/vertex_resolution_tracks.pdf}
    \end{figure}
\end{frame}

\begin{frame}{AMVF CPU Time Per Event}
    \begin{figure}
        \centering
        \includegraphics[width=0.7\linewidth]{APS2025/acts_pvfinder_timing.png}
    \end{figure}
\end{frame}

\begin{frame}[plain]
    \centering
    TVPD April 4th, 2025
\end{frame}

\begin{frame}{PV-Finder Recap}
    \vspace{25pt}
    \begin{figure}
        \centering
        \includegraphics[width=1\linewidth]{april4th/before_pipeline.png}
        \label{fig:before_pipeline}
    \end{figure}
\end{frame}

\begin{frame}{Tracks to KDE: Data Pre-Processing Updated}
    \begin{figure}
        \centering
        \includegraphics[width=1\linewidth]{april4th/updated_data_preprocessing.png}
        \label{fig:updated_data_preprocessing}
    \end{figure}
\end{frame}

\begin{frame}{Tracks to KDE: Architecture Updated}
    \vspace{25pt}
    \begin{figure}
        \centering
        \includegraphics[width=1\linewidth]{april4th/mlp_pipeline_graph.png}
        \label{fig:updated_pipeline}
    \end{figure}
\end{frame}

\begin{frame}{Example Predicted KDE Outputs}
    \begin{figure}
        \centering
        \includegraphics[width=0.6\linewidth]{APS2025/KDE_A_pred_example_workinprogress.png}
        \label{fig:example_KDE_output_2}
    \end{figure}
\end{frame}

\begin{frame}{A Brief Review of UNet}
    \begin{columns}
        \begin{column}{0.4\linewidth}
            \begin{itemize}
                \item Effective for feature extraction
                \item Downsampling reduces dimensionality, helping emphasize important data features
                \item Skip connections propogate prior information
                \item \(\sim 279,000\) Parameters
            \end{itemize}
        \end{column}

        \begin{column}{0.6\linewidth}
            \begin{figure}
                \centering
                \includegraphics[width=1\linewidth]{april4th/UNet.png}
                \label{fig:pipeline_kdetohist}
            \end{figure}
        \end{column}
    \end{columns}
\end{frame}

\begin{frame}{A Brief Review of UNet++}
    \begin{columns}
        \begin{column}{0.4\linewidth}
            \begin{itemize}
                \item Just more skip connections and upsamplings
                \item \(\sim 796,000\) Parameters
            \end{itemize}
        \end{column}

        \begin{column}{0.6\linewidth}
            \begin{figure}
                \centering
                \includegraphics[width=1\linewidth]{april4th/unetplusplus_figure.png}
            \end{figure}
        \end{column}
    \end{columns}
\end{frame}

\begin{frame}{New PV-Finder}
    \vspace{5pt}
    \centering
    Connecting the MLP and UNet resolves the need to have an intermediary representation
    \begin{figure}
        \centering
        \includegraphics[width=1\linewidth]{april4th/question_updated_tracks_to_pv.png}
        \label{fig:oldPVFinderArchitecture}
    \end{figure}
\end{frame}

\begin{frame}{Connecting the MLP and UNet (1)}
    \vspace{5pt}
    \begin{figure}
        \centering
        \includegraphics[width=0.9\linewidth]{april4th/redo_frankenstein_1.png}
    \end{figure}
\end{frame}

\begin{frame}{Connecting the MLP and UNet (2)}
    \vspace{5pt}
    \begin{figure}
        \centering
        \includegraphics[width=0.9\linewidth]{april4th/redo_frankenstein_2.png}
    \end{figure}
\end{frame}

\begin{frame}{Connecting the MLP and UNet (3)}
    \vspace{5pt}
    \begin{figure}
        \centering
        \includegraphics[width=0.9\linewidth]{april4th/redo_frankenstein_3.png}
    \end{figure}
\end{frame}

\begin{frame}{Connecting the MLP and UNet (4)}
    \vspace{5pt}
    \begin{figure}
        \centering
        \includegraphics[width=1\linewidth]{april4th/redo_frankenstein_4.png}
    \end{figure}
\end{frame}

\begin{frame}{New PV-Finder}
    \vspace{5pt}
    \centering
    ...and we are done! Don't need KDE-A-x, KDE-A-y, or KDE-B-z anymore!
    \begin{figure}
        \centering
        \includegraphics[width=1\linewidth]{april4th/updated_tracks_to_pv.png}
        \label{fig:oldPVFinderArchitecture}
    \end{figure}
\end{frame}

\begin{frame}{Example End-to-end tracks-to-hist Output}
    \begin{figure}
        \centering
        \includegraphics[width=0.6\linewidth]{APS2025/embed_pred_KDE_A_pv_pred_pv_final.png}
        \label{fig:example_Hists_output_1}
    \end{figure}
\end{frame}

\begin{frame}{A Brief Intermission}
    \centering
    \textbf{Any questions on the architecture updates before I move on?}
\end{frame}

\begin{frame}[noframenumbering]{A Brief Intermission}
    \begin{itemize}
        \item Now that we have a preliminary end-to-end PV-Finder architecture there are many parameters to choose and adjust
        \begin{itemize}
            \item UNet vs UNet++?
            \item Binnings?
            \begin{itemize}
                \item 100? \textbf{400? 1000?}
            \end{itemize}
            \item Architecture Add-ons?
            \begin{itemize}
                \item \textbf{Dropout?}
                \item Layer Depth?
            \end{itemize}
        \end{itemize}
        \item New results shown in upcoming part of presentation
        \begin{enumerate}
            \item (Sanity Check/Review) (400 Bins) MLP+UNet vs Prior Results
            \item (1000 Bins) MLP+UNet with dropout vs MLP+UNet++ with dropout
        \end{enumerate}
    \end{itemize}
\end{frame}

\begin{frame}[plain]
    \centering
    \textbf{(400 Bins) MLP+UNet vs Prior Results}
\end{frame}

\begin{frame}{400 Bins}
    \vspace{5pt}
    \begin{figure}
        \centering
        \includegraphics[width=1\linewidth]{april4th/frankenstein_400bins.png}
    \end{figure}
\end{frame}

\begin{frame}{\(\Delta z\) Resolution and \(\sigma_{\text{vtx-vtx}}\) Calculation}
    \begin{columns}
        \begin{column}{0.6\linewidth}
            \begin{figure}
                \centering
                \includegraphics[width=0.9\linewidth]{APS2025/deltaz(1).png}
            \end{figure}
        \end{column}

        \begin{column}{0.4\linewidth}
            \begin{figure}
                \centering
                \includegraphics[width=1\linewidth]{APS2025/methods_sigma_vtx_vtx2.png}
            \end{figure}
        \end{column}
    \end{columns}
\end{frame}

\begin{frame}{Comparison with AVMF: Vertex Classification Scheme}
    \begin{figure}
        \centering
        \includegraphics[width=0.75\linewidth]{APS2025/vertex_classification_aps2025.png}
    \end{figure}
\end{frame}

\begin{frame}{Comparison with AMVF: Efficiency}
    \centering
    \begin{figure}
        \centering
        \includegraphics[width=0.7\linewidth]{APS2025/ntrks_fpr_eff_fixed_2.png}
        \label{fig:efffprresults}
    \end{figure}
\end{frame}

\begin{frame}[plain]
    \centering
    \textbf{(1000 Bins) MLP+UNet with Dropout vs MLP+UNet++ with Dropout}
\end{frame}

\begin{frame}{\(\Delta z\) Resolution and \(\sigma_{\text{vtx-vtx}}\) Calculation}
    \begin{columns}
        \begin{column}{0.6\linewidth}
            \begin{figure}
                \centering
                \includegraphics[width=0.9\linewidth]{april4th/deltaz(2).png}
            \end{figure}
        \end{column}

        \begin{column}{0.4\linewidth}
            \begin{figure}
                \centering
                \includegraphics[width=1\linewidth]{april4th/sigma_vtx_vtx_dropout.png}
            \end{figure}
        \end{column}
    \end{columns}
\end{frame}

\begin{frame}{Dropout UNet vs UNet++ Vertex Classification}
    \begin{figure}
        \centering
        \includegraphics[width=0.5\linewidth]{april4th/avg_per_category(1).png}
    \end{figure}
\end{frame}

\begin{frame}{Dropout UNet vs UNet++ Efficiency}
    \centering
    UNet and UNet++ perform basically identically!
    \begin{figure}
        \centering
        \includegraphics[width=1\linewidth]{april4th/droput_unet_unetplusplus_eff.png}
    \end{figure}
\end{frame}

\begin{frame}{Merge/Split Case: PV-Finder UNet/UNet++ Output (1)}
    \begin{figure}
        \centering
        \includegraphics[width=0.6\linewidth]{april4th/edgecase_t2h_unet_unetpp_1.png}
    \end{figure}
\end{frame}

\begin{frame}{Merge Case: PV-Finder UNet/UNet++ Output (2)}
    \begin{figure}
        \centering
        \includegraphics[width=0.6\linewidth]{april4th/edgecase_t2h_unet_unetpp_2.png}
    \end{figure}
\end{frame}

\begin{frame}{CPU Time Per Event}
    \begin{itemize}
        \item UNet: 5.3 ms
        \item UNet++: 5.5 ms
    \end{itemize}
    \begin{figure}
        \centering
        \includegraphics[width=0.6\linewidth]{APS2025/acts_pvfinder_timing.png}
    \end{figure}
\end{frame}

\begin{frame}{Conclusion}
    \begin{itemize}
        \item Lots of great progress has been made on PV-Finder!
        \begin{itemize}
            \item Including the binning max/min helped the MLP learn the KDE function
            \item Combining the MLP and UNet/UNet++ has drastically reduced the FPR (as intended)
            \item Dropout has helped the architectures outperform AMVF
        \end{itemize}
        \item Upcoming Work!
        \begin{itemize}
            \item GNN for Track-Vertex Association
            \item Further Studies
            \begin{itemize}
                \item More variables (e.g. \(p_T, \theta, \phi \))
                \item Optimizing the end-to-end architecture (e.g. more/less layers)
            \end{itemize}
            \item Test architecture on HL-LHC simulated conditions
        \end{itemize}
    \end{itemize}
\end{frame}

\begin{frame}{Relevant Links and Directories}
    \begin{itemize}
        \item LHCb PV-Finder
        \begin{itemize}
            \item \href{https://gitlab.cern.ch/sakar/pv-finder_v2/-/tree/master}{Link to Gitlab}
        \end{itemize}
        \item ATLAS PV-Finder
        \begin{itemize}
            \item \href{https://gitlab.cern.ch/qlei/atlas_pvfinder}{Link to Gitlab}
        \end{itemize}
        \item Prior Presentations:
        \begin{itemize}
            \item TVPD Oct 2024: \href{https://drive.google.com/file/d/17qCnL57fT7LCKyQ1Rz-soJySaNCVeMRV/view?usp=sharing}{Slides}
            \item US LUA 2024: \href{https://drive.google.com/file/d/1HbAjlKBw6nej-oqOksUL_HOfKY8xsd2r/view?usp=sharing}{Slides}
            \item APS Joint Meeting 2025: \href{https://drive.google.com/file/d/1j6pnqzuXcP0OWXoUXkwElH8lAbv9GtRq/view?usp=drive_link}{Slides}
        \end{itemize}
    \end{itemize}
\end{frame}

\begin{frame}{Overview of Parameter Space}
    \centering
    \begin{tabular}{ c | c | c | c }
        \multicolumn{4}{c}{General Parameters} \\
        \hline
        Models & MLP + UNet & MLP + UNet++ &  \\
        \hline
        Bin Size & 100 & 400 & 1000 \\
        \hline
        Extra Inputs & \(\theta, \phi, p_T\) & POCA Ellipsoid & \\
    \end{tabular}

    \vspace{5pt}
    \begin{tabular}{ c | c | c }
        \multicolumn{3}{c}{Model Specific Parameters} \\
        \hline
        UNet & \# of Convo+BN+ReLu=1 & Dropout Rate=0.25  \\
        \hline
        UNet++ & \# of Convo+BN+ReLu=1 & Dropout Rate=0.25 \\
    \end{tabular}

    \vspace{5pt}
    \begin{tabular}{ c | c }
        \multicolumn{2}{c}{Fitting Parameters} \\
        \hline
        Integral Threshold & 0.2 to 0.7 \\
        \hline
        Integral Width & 3 to 5 \\
    \end{tabular}
\end{frame}

\begin{frame}{400 Bins MLP+UNet: Number of Primary Vertives vs \(\mu\)}
    \begin{figure}
        \centering
        \includegraphics[width=0.7\linewidth]{APS2025/npv_compare_Total.png}
    \end{figure}
\end{frame}

\begin{frame}{1000 Bins MLP+UNet/UNet++: Number of Primary Vertives vs \(\mu\)}
    \begin{figure}
        \centering
        \includegraphics[width=0.7\linewidth]{april4th/npv_compare_Total(1).png}
    \end{figure}
\end{frame}

\begin{frame}{List of Other Results}
    \centering
    \begin{tabular}{c c c}
        \hline
        Model & Efficiency & False Positive Rate (\# / event) \\
        \hline
        T2H UNet 400 Bins & 90.4\% & 0.97 \\
        T2H UNet 400 Bins + \(\theta + \phi\) & 91.4\% & 0.96 \\
        T2H UNet 400 Bins + \(\theta + \phi\) + 8 Latent & 90.9\% & 0.81 \\
        \textbf{T2H UNet 100 Bins + \(\theta + \phi\)} & \textbf{91.8\%} & \textbf{0.73} \\
        \textbf{T2H UNet 1000 Bins + 8 Latent} & \textbf{92.8\%} & \textbf{0.81} \\
        AMVF & 93.9\% & 0.8 \\
        \hline
    \end{tabular}

    \vspace{7pt}

    \begin{tabular}{c c}
        \hline
        Method & \(\sigma_{\text{vtx-vtx}}\) (mm) \\
        \hline
        T2H UNet 400 bins & \(0.29 \pm 0.03\) \\
        T2H UNet 400 bins + \(\theta + \phi\) & \(0.30 \pm 0.02\) \\
        T2H UNet 400 Bins + \(\theta + \phi\) + 8 Latent & \(0.31 \pm 0.01\) \\
        T2H UNet 100 Bins + \(\theta + \phi\) & \(0.33 \pm 0.01\) \\
        T2H UNet 1000 Bins + 8 Latent & \(0.37 \pm 0.01\) \\
        AMVF & \(0.76 \pm 0.02\) \\
        \hline
    \end{tabular}
\end{frame}

\begin{frame}{Overfitting Musings (1)}
    \begin{itemize}
        \item Why does adding dropout drastically improve efficiency?
        \item Why are the UNet results in the pubnote much worse than UNet++?
        \item Pubnote Parameters of Interest:
        \begin{itemize}
            \item Epochs: 500
            \item Number of Parameters:
            \begin{itemize}
                \item UNet: 279,105
                \item UNet++: 796,417
            \end{itemize}
            \item Dropout Rate:
            \begin{itemize}
                \item UNet: \textbf{0\%}
                \item UNet++: 25\%
            \end{itemize}
        \end{itemize}
        \item I think UNet++ was more robust against overfitting causing it to perform better
    \end{itemize}
\end{frame}

\begin{frame}{Overfitting Musings (2)}
    \centering
    Here is the loss plot of the pubnote's training for UNet (overfitting)
    \begin{figure}
        \centering
        \includegraphics[width=1\linewidth]{april4th/overtrained_unet.png}
    \end{figure}
\end{frame}

\begin{frame}{Overfitting Musings (3)}
    \centering
    Here is the loss plot of the pubnote's training for UNet++ (less overfitting)
    \begin{figure}
        \centering
        \includegraphics[width=1\linewidth]{april4th/overtrained_unetplusplus.png}
    \end{figure}
\end{frame}

\begin{frame}{Upcoming Work: PV-Finder + GNN Fitter Idea}
    \begin{figure}
        \centering
        \includegraphics[width=0.55\linewidth]{april4th/pipeline_pvfinder_gnn.png}
    \end{figure}
\end{frame}

\begin{frame}{GNN Concept}
    \begin{figure}
        \centering
        \includegraphics[width=0.55\linewidth]{april4th/gnn_concept.png}
    \end{figure}
\end{frame}

\begin{frame}{Forming Graphs}
    \begin{figure}
        \centering
        \includegraphics[width=0.6\linewidth]{april4th/forming_edges_gnn.png}
    \end{figure}
\end{frame}

\begin{frame}{Scoring Edges in Graphs}
    \begin{figure}
        \centering
        \includegraphics[width=0.6\linewidth]{april4th/scoring_edges_gnn.png}
    \end{figure}
\end{frame}

\begin{frame}{Choosing Edges in Graphs}
    \begin{figure}
        \centering
        \includegraphics[width=0.6\linewidth]{april4th/choosing_edges_gnn.png}
    \end{figure}
\end{frame}

\begin{frame}[plain]
    \centering
    Tompkins Group Meeting April 15th
\end{frame}

\begin{frame}{Preview}
    \begin{itemize}
        \item Going over new PV-Finder results
        \item Going over, PV-Fitter, a GNN used to associate tracks and vertices
    \end{itemize}
\end{frame}

\begin{frame}{PV-Finder}
    \vspace{5pt}
    \centering
    Connecting the MLP and UNet resolves the need to have an intermediary representation
    \begin{figure}
        \centering
        \includegraphics[width=1\linewidth]{APS2025/tracks_to_pv.png}
        \label{fig:oldPVFinderArchitecture}
    \end{figure}
\end{frame}

\begin{frame}{A Brief Review of UNet}
    \begin{columns}
        \begin{column}{0.4\linewidth}
            \begin{itemize}
                \item Effective for feature extraction
                \item Downsampling reduces dimensionality, helping emphasize important data features
                \item Skip connections propogate prior information
                \item \(\sim 279,000\) Parameters
            \end{itemize}
        \end{column}

        \begin{column}{0.6\linewidth}
            \begin{figure}
                \centering
                \includegraphics[width=1\linewidth]{april4th/UNet.png}
                \label{fig:pipeline_kdetohist}
            \end{figure}
        \end{column}
    \end{columns}
\end{frame}

\begin{frame}{A Brief Review of UNet++}
    \begin{columns}
        \begin{column}{0.4\linewidth}
            \begin{itemize}
                \item Just more skip connections and upsamplings
                \item \(\sim 796,000\) Parameters
            \end{itemize}
        \end{column}

        \begin{column}{0.6\linewidth}
            \begin{figure}
                \centering
                \includegraphics[width=1\linewidth]{april4th/unetplusplus_figure.png}
            \end{figure}
        \end{column}
    \end{columns}
\end{frame}

\begin{frame}{Example End-to-end tracks-to-hist Output}
    \begin{figure}
        \centering
        \includegraphics[width=0.6\linewidth]{APS2025/embed_pred_KDE_A_pv_pred_pv_final.png}
        \label{fig:example_Hists_output_1}
    \end{figure}
\end{frame}

\begin{frame}[plain]
    \centering
    \textbf{(1000 Bins) MLP+UNet with Dropout vs MLP+UNet++ with Dropout}
\end{frame}

\begin{frame}{\(\Delta z\) Resolution and \(\sigma_{\text{vtx-vtx}}\) Calculation}
    \begin{columns}
        \begin{column}{0.6\linewidth}
            \begin{figure}
                \centering
                \includegraphics[width=0.9\linewidth]{april4th/deltaz(2).png}
            \end{figure}
        \end{column}

        \begin{column}{0.4\linewidth}
            \begin{figure}
                \centering
                \includegraphics[width=1\linewidth]{april4th/sigma_vtx_vtx_dropout.png}
            \end{figure}
        \end{column}
    \end{columns}
\end{frame}

\begin{frame}{Dropout UNet vs UNet++ Efficiency}
    \centering
    UNet and UNet++ perform basically identically!
    \begin{figure}
        \centering
        \includegraphics[width=1\linewidth]{april4th/droput_unet_unetplusplus_eff.png}
    \end{figure}
\end{frame}

\begin{frame}{CPU Time Per Event}
    \begin{itemize}
        \item UNet: 5.3 ms
        \item UNet++: 5.5 ms
        \item Note: Not an exact apples to apples comparison (different dataset)
    \end{itemize}
    \begin{figure}
        \centering
        \includegraphics[width=0.55\linewidth]{APS2025/acts_pvfinder_timing.png}
    \end{figure}
\end{frame}

\begin{frame}{Merge/Split Case: PV-Finder UNet/UNet++ Output (1)}
    \begin{figure}
        \centering
        \includegraphics[width=0.6\linewidth]{april4th/edgecase_t2h_unet_unetpp_1.png}
    \end{figure}
\end{frame}

\begin{frame}{Merge Case: PV-Finder UNet/UNet++ Output (2)}
    \begin{figure}
        \centering
        \includegraphics[width=0.6\linewidth]{april4th/edgecase_t2h_unet_unetpp_2.png}
    \end{figure}
\end{frame}

\begin{frame}{PV-Finder + GNN Fitter Idea}
    \begin{figure}
        \centering
        \includegraphics[width=0.55\linewidth]{april4th/pipeline_pvfinder_gnn.png}
    \end{figure}
\end{frame}

\begin{frame}{GNN Concept}
    \begin{figure}
        \centering
        \includegraphics[width=0.55\linewidth]{april4th/gnn_concept.png}
    \end{figure}
\end{frame}

\begin{frame}{Forming Graphs}
    \begin{figure}
        \centering
        \includegraphics[width=0.6\linewidth]{april4th/forming_edges_gnn.png}
    \end{figure}
\end{frame}

\begin{frame}{Scoring Edges in Graphs}
    \begin{figure}
        \centering
        \includegraphics[width=0.6\linewidth]{april4th/scoring_edges_gnn.png}
    \end{figure}
\end{frame}

\begin{frame}{Choosing Edges in Graphs}
    \begin{figure}
        \centering
        \includegraphics[width=0.6\linewidth]{april4th/choosing_edges_gnn.png}
    \end{figure}
\end{frame}

\begin{frame}{GNN Pipeline}
    \begin{figure}
        \centering
        \includegraphics[width=1\linewidth]{april15th/gnn_architecture.png}
    \end{figure}
\end{frame}

\begin{frame}{MLP Embedding Layer}
    \begin{figure}
        \centering
        \includegraphics[width=0.7\linewidth]{april15th/mlp_embedding_layer_32.png}
    \end{figure}
\end{frame}

\begin{frame}{GNN Pipeline After Embedding}
    \begin{figure}
        \centering
        \includegraphics[width=1\linewidth]{april15th/gnn_architecture_after_embedding.png}
    \end{figure}
\end{frame}

\begin{frame}{GraphConv Layer}
    \begin{figure}
        \centering
        \includegraphics[width=0.8\linewidth]{april15th/graphconv_example.png}
    \end{figure}
\end{frame}

\begin{frame}{GraphConv Layer with Embedding}
    \begin{figure}
        \centering
        \includegraphics[width=0.8\linewidth]{april15th/embedding_gcn_pv.png}
    \end{figure}
\end{frame}

\begin{frame}{Heterogeneous GraphConv Layer with Embedding}
    \begin{figure}
        \centering
        \includegraphics[width=0.8\linewidth]{april15th/heterogeneousgraphconv_embedding.png}
    \end{figure}
\end{frame}

\begin{frame}{GNN Pipeline After GCN}
    \begin{figure}
        \centering
        \includegraphics[width=1\linewidth]{april15th/gnn_architecture_after_gcn(1).png}
    \end{figure}
\end{frame}

\begin{frame}{MLP Edge Predictor Layer}
    \begin{figure}
        \centering
        \includegraphics[width=1\linewidth]{april15th/mlp_edge_predictor_layer.png}
    \end{figure}
\end{frame}

\begin{frame}{End-to-end GNN Pipeline}
    \begin{figure}
        \centering
        \includegraphics[width=0.9\linewidth]{april15th/endtoend_gnn_pipeline.png}
    \end{figure}
\end{frame}


\begin{frame}{Assessing Performance}
    \begin{itemize}
        \item "track passing" means scoring greater than 0.5 by GNN
        \begin{itemize}
            \item Note: can have harder cutoffs, like score greater than 0.6, 0.7, . . .
        \end{itemize}
        \item Inclusivity of tracks within a certain range of PV
        \begin{itemize}
            \item \( \abs{z_{0, \text{track}} - z_{\text{pv}}} <= 5\text{mm} \text{ or } 10\text{mm}\)
        \end{itemize}
        \item \textbf{AUC, FPR, FNR, TPR, TNR}
        \item \textbf{Efficiency: } (track is truth associated to the supplied vertex AND track passes GNN) / (track is truth associated to the supplied vertex)
        \item \textbf{Impurity: } (track passes GNN AND is NOT associated to supplied vertex) / (track passes GNN)
        \item Borrowed from \href{https://indico.cern.ch/event/1025152/contributions/4304219/attachments/2220591/3761247/210406_Basso_ATLAS-Weekly_TTVA-for-Run3.pdf}{Link}
    \end{itemize}
\end{frame}

\begin{frame}{Example Event: ROC/AUC (10mm)}
    \begin{columns}
        \begin{column}{0.3\linewidth}
            \begin{itemize}
                \item \# of tracks: 3986
                \item \# of PVs: 50
                \item TPR: 96.79\%
                \item FPR: 11.02\%
                \item FNR: 3.21\%
                \item TNR: 88.98\%
            \end{itemize}
        \end{column}

        \begin{column}{0.7\linewidth}
            \begin{figure}
                \centering
                \includegraphics[width=.95\linewidth]{april15th/sample_ROC_GNN.png}
            \end{figure}
        \end{column}
    \end{columns}
\end{frame}

\begin{frame}{Efficiency vs \(\eta\) with Varying Thresholds (5mm Range)}
    \centering
    \begin{figure}
        \centering
        \includegraphics[width=0.65\linewidth]{april15th/efficiency_vs_eta_multiple_thresholds_5.png}
    \end{figure}
\end{frame}

\begin{frame}{Impurity vs \(\eta\) with Varying Thresholds (5mm Range)}
    \begin{figure}
        \centering
        \includegraphics[width=0.65\linewidth]{april15th/impurity_vs_eta_multiple_thresholds_5.png}
    \end{figure}
\end{frame}

\begin{frame}{Efficiency vs \(\eta\) with Varying Thresholds (10mm Range)}
    \centering
    \begin{figure}
        \centering
        \includegraphics[width=0.65\linewidth]{april15th/efficiency_vs_eta_multiple_thresholds_10.png}
    \end{figure}
\end{frame}

\begin{frame}{Impurity vs \(\eta\) with Varying Thresholds (10mm Range)}
    \begin{figure}
        \centering
        \includegraphics[width=0.65\linewidth]{april15th/impurity_vs_eta_multiple_thresholds_10.png}
    \end{figure}
\end{frame}

\begin{frame}{Conclusion}
    \begin{itemize}
        \item Lots of great progress has been made on PV-Finder!
        \begin{itemize}
            \item Combining the MLP and UNet/UNet++ has drastically reduced the FPR (as intended)
            \item Dropout has helped the architectures outperform AMVF
        \end{itemize}
        \item Interesting preliminary results on GNN for Track-Vertex Association
        \item Upcoming Work!
        \begin{itemize}
            \item Further Studies
            \begin{itemize}
                \item More variables (e.g. \(p_T, \theta, \phi \))
                \item Optimizing the end-to-end architecture (e.g. more/less layers)
            \end{itemize}
            \item GNN comparison with AMVF
            \item Test architecture on HL-LHC simulated conditions
            \item Start Muon Collider Software/Analysis!
        \end{itemize}
    \end{itemize}
\end{frame}

\begin{frame}{GNN Loss Function}
    \begin{itemize}
        \item Binary Cross Entropy Loss Function
    \end{itemize}
    \begin{align*}
        \mathcal{L}_{\text{BCE}}(y_{\text{pred}}, y_{\text{truth}}) = y_{\text{truth}} \log(y_{\text{pred}}) + (1 - y_{\text{truth}}) \log ((1- y_{\text{pred}}))
    \end{align*}
\end{frame}

\begin{frame}[plain]
    \centering
    PV-Finder Meeting with LHCb 05/20/25
\end{frame}

\begin{frame}{Updates}
    \begin{itemize}
        \item Next Meeting: June 3rd, 8am PST
        \item (No Success) Test modifying PV-Finder loss function
        \begin{itemize}
            \item Looks about the same as before
        \end{itemize}
        \item Current Studies:
        \begin{itemize}
            \item Incorporating peak information into GNN
            \item Having only one association per track vs threshold
            \item Focal Loss
        \end{itemize}
        \item Other:
        \begin{itemize}
            \item Python KDE Generation Script for ACTS
        \end{itemize}
    \end{itemize}
\end{frame}

\begin{frame}{Recap: GNN Pipeline}
    \begin{figure}
        \centering
        \includegraphics[width=0.9\linewidth]{april15th/endtoend_gnn_pipeline.png}
    \end{figure}
\end{frame}

\begin{frame}{Recap: Assessing Performance}
    \begin{itemize}
        \item Track is associated with a primary vertex if
        \begin{itemize}
            \item Method 1: Score greater than 0.5
            \begin{itemize}
                \item Note: can have harder cutoffs, like score greater than 0.6, 0.7, . . .
            \end{itemize}
            \item Method 2: Take the highest scoring edge for each track
        \end{itemize}
        \item Inclusivity of tracks within a certain range of PV
        \begin{itemize}
            \item \( \abs{z_{0, \text{track}} - z_{\text{pv}}} <= 5\text{mm} \text{ or } 10\text{mm}\)
        \end{itemize}
        \item \textbf{Efficiency: } (track is truth associated to the supplied vertex AND track passes GNN) / (track is truth associated to the supplied vertex)
        \item \textbf{Impurity: } (track passes GNN AND is NOT associated to supplied vertex) / (track passes GNN)
        \item Borrowed from \href{https://indico.cern.ch/event/1025152/contributions/4304219/attachments/2220591/3761247/210406_Basso_ATLAS-Weekly_TTVA-for-Run3.pdf}{Link}
    \end{itemize}
\end{frame}

\begin{frame}[plain]
    \centering
    Including PV Histogram Height vs No Height Studies
\end{frame}

\begin{frame}{5mm Include vs Exclude PV Histogram Height: Efficiency}
    \begin{figure}
        \centering
        \includegraphics[width=0.65\linewidth]{may16th/5mm_inc_vs_no_height_efficiency_vs_eta_multiple_thresholds.png}
    \end{figure}
\end{frame}

\begin{frame}{5mm Include vs Exclude PV Histogram Height: Impurity}
    \begin{figure}
        \centering
        \includegraphics[width=0.65\linewidth]{may16th/5mm_inc_vs_no_height_impurity_vs_eta_multiple_thresholds.png}
    \end{figure}
\end{frame}

\begin{frame}[plain]
    \centering
    1 Edge Limit Per Track Studies
\end{frame}

\begin{frame}{5mm Include vs Exclude Height 1 Edge Limit: Efficiency}
    \begin{figure}
        \centering
        \includegraphics[width=0.65\linewidth]{may16th/combined_efficiency_vs_eta_multiple_thresholds_and_1edgelimit.png}
    \end{figure}
\end{frame}

\begin{frame}{5mm Include vs Exclude Height 1 Edge Limit: Impurity}
    \begin{figure}
        \centering
        \includegraphics[width=0.65\linewidth]{may16th/combined_impurity_vs_eta_multiple_thresholds_and_1edgelimit.png}
    \end{figure}
\end{frame}

\begin{frame}[plain]
    \centering
    5mm Include Heights Loss Function Studies
\end{frame}

\begin{frame}{GNN Loss Function}
    \begin{itemize}
        \item Notation definition, for \(y\in \left\{\pm 1\right\}\) (truth class) and \( p \in \left[ 0, 1 \right] \) (predicted class),
        \begin{align*}
            p_t = \left\{\begin{array}{lr}
                p, & \text{if } y=1\\
                1-p, & \text{otherwise.}
                \end{array}\right.
        \end{align*}
        \item Binary Cross Entropy Loss Function
    \end{itemize}
    \begin{align*}
        CE(p,y) = \left\{\begin{array}{lr}
            -\log(p_t), & \text{if } y=1\\
            -\log(p_t), & \text{otherwise.}
            \end{array}\right.
        % \mathcal{L}_{\text{BCE}}(y_{\text{pred}}, y_{\text{truth}}) = y_{\text{truth}} \log(y_{\text{pred}}) + (1 - y_{\text{truth}}) \log ((1- y_{\text{pred}}))
    \end{align*}
    \begin{itemize}
        \item Focal Loss Function (\href{https://arxiv.org/pdf/1708.02002}{1708.02002})
        \begin{itemize}
            \item For a tunable focusing parameter (\(\gamma\)) and a balancing variable (\(\alpha\)),
        \end{itemize}
    \end{itemize}
    \begin{align*}
        \alpha_t &= \left\{\begin{array}{lr}
                \alpha, & \text{if } y=1\\
                1-\alpha, & \text{otherwise.}
                \end{array}\right. \\
        FL(p_t) &= -\alpha_t \left( 1 - p_t \right)^\gamma \log(p_t)
    \end{align*}
    \begin{itemize}
        \item Allows us to deal with large class imbalances and ignore well understood examples
    \end{itemize}
\end{frame}

\begin{frame}{BCE vs Focal Loss (\(\alpha = 0.25, \gamma = 4\)): Efficiency}
    \centering
    Focal Loss bins with Threshold of 0.9 and 0.7 are empty
    \begin{figure}
        \centering
        \includegraphics[width=0.65\linewidth]{may16th/combined_efficiency_vs_eta_multiple_thresholds_and_1edgelimit_BCE_vs_Focal.png}
    \end{figure}
\end{frame}

\begin{frame}{BCE vs Focal Loss (\(\alpha = 0.25, \gamma = 4\)): Efficiency}
    \centering
    Focal Loss bins with Threshold of 0.9 and 0.7 are empty
    \begin{figure}
        \centering
        \includegraphics[width=0.65\linewidth]{may16th/combined_impurity_vs_eta_multiple_thresholds_and_1edgelimit_BCE_vs_Focal.png}
    \end{figure}
\end{frame}

\begin{frame}{Discussion on Focal Loss}
    \begin{itemize}
        \item The outputs of GNN using Focal Loss results is mainly composed of values <0.5
        \item Method 2 (choosing highest probable edge) still works and shows that BCE/Focal performs about the same
        \item Will revisit and do thorough study using different \(\alpha\) and \(\gamma\) parameters
    \end{itemize}
\end{frame}

\begin{frame}{Other: Python Script for KDE Generation}
    \begin{itemize}
        \item LHCb has a general python script for POCA variables to KDE
        \item This is a general python script for track variables to KDE
        \item 26 seconds for 10 events \(\rightarrow\) 36.8 hours for 51,000 events
    \end{itemize}
    \begin{figure}
        \centering
        \includegraphics[width=0.55\linewidth]{may16th/old_vs_new_script_KDE_gen.png}
    \end{figure}
\end{frame}

\begin{frame}{Concluding Remarks}
    \begin{itemize}
        \item Upcoming Work:
        \begin{itemize}
            \item Incorporating bins nearby each peak (\(\pm 5\) bins) into PV information
            \item Using AMVF metrics and \(p_T\) reweighting for proper reco track-vertex studies
            \begin{itemize}
                \item Begin studies on using predicted PV location from PV-Finder
            \end{itemize}
            \item Incorporate GNN and relevant scripts into GitLab
            \item Continue pipeline generalization for ACTS implementation
            \begin{itemize}
                \item Eventually want to try on OpenDetector data
            \end{itemize}
        \end{itemize}
        \item Idea Musings:
        \begin{itemize}
            \item Having narrower PV histogram distributions?
            \begin{itemize}
                \item Instead of \(\left[-5,6\right]\) you can have \(\left[-2,3\right]\)?
            \end{itemize}
            \item Incorporation of PV x and y location into output of PV-Finder?
            \begin{itemize}
                \item Use as additional input information into GNN
                \item Possibly useful for displaced vertices and tracks (e.g. \(B\) decays)?
            \end{itemize}
        \end{itemize}
    \end{itemize}
\end{frame}

\begin{frame}{Focal Loss Expanded}
    \begin{align*}
        \mathcal{L}_{\text{FL}}(y_{\text{pred}}, y_{\text{truth}}) = y_{\text{truth}} (1-y_{\text{pred}})^\gamma \log(y_{\text{pred}}) + (1 - y_{\text{truth}}) (y_{\text{pred}})^\gamma \log ((1- y_{\text{pred}}))
    \end{align*}
\end{frame}

\begin{frame}[plain]
    \centering
    Friday Meeting 05/24/2025
\end{frame}

\begin{frame}[plain]
    \centering
    Vertex Classification
\end{frame}

\begin{frame}{ATLAS Vertex Classification (1)}
    \begin{itemize}
        \item Given a track, \(j\), and edge score (or weight), \(w_j\), of how compatible a track is to a particular reconstructed vertex, \(V_{reco}\)
        \item You can calculate
        \begin{align*}
            W_{total, V_{reco}} = \sum_{j \in V_{reco}} w_j
        \end{align*}
        \item For each truth proton-proton interaction/primary vertex (\(P_A, P_B, P_C...\)) that contributed tracks to \(V_{reco}\), sum the weights of the tracks
        \begin{align*}
            P_A: W_{P_A\_in\_V_{reco}} = \sum_{j\in V_{reco} \text{ and } j\in P_A} w_j
        \end{align*}
        \item Do the same for all truth primary vertices that has tracks potentially associated with reconstructed primary vertex
    \end{itemize}
\end{frame}

\begin{frame}{ATLAS Vertex Classification (2)}
    \begin{itemize}
        \item Find maximum contribution from single interaction
        \begin{align*}
            W_{max \_ assoc} = max(W_{P_A\_in\_V_{reco}}, W_{P_B\_in\_V_{reco}}, ...)
        \end{align*}
        \item Clean Classification:
        \begin{align*}
            \frac{W_{max \_ assoc}}{W_{total, V_{reco}}} \geq 0.70
        \end{align*}
        \item Merged Classification:
        \begin{align*}
            \frac{W_{max \_ assoc}}{W_{total, V_{reco}}} < 0.70
        \end{align*}
    \end{itemize}
\end{frame}

\begin{frame}{ATLAS Vertex Classification (3)}
    \begin{itemize}
        \item Split Classification:
        \begin{itemize}
            \item A true primary vertex contributes largest track weight to two or more reconstructed primary vertices
            \item Reconstructed vertex with largest \(\sum p_T^2\) is either clean or merged and the other is split
            \item (Haven't implemented yet, just being lazy about it)
        \end{itemize}
        \item Fake Classification:
        \begin{itemize}
            \item Fake tracks contribute more weight to reconstructed primary vertex than tracks from true primary vertices
        \end{itemize}
    \end{itemize}
\end{frame}

\begin{frame}{Re: Methods for Evaluation}
    \begin{itemize}
        \item Since we can just take the outputs of the GNN as just weights there are multiple ways to adjust the weights
        \item Method 1 (None): No adjustment
        \item Method 2 (Round): Threshold of having a score be greater than 0.5
        \begin{itemize}
            \item Equivalent to rounding the edges
        \end{itemize}
        \item Method 3 (MaxScore): Take the highest scoring edge for each track and set weight as one
    \end{itemize}
\end{frame}

\begin{frame}{Case Study: Looking at Merged Cases}
    \begin{figure}
        \centering
        \includegraphics[width=0.6\linewidth]{may22nd/example_merged_track_assoc.png}
        \caption{Gut check: height of PV in histogram is equivalent to number of tracks associated with it}
    \end{figure}
\end{frame}

\begin{frame}{Case Study: None on Example PV 25}
    \begin{itemize}
        \item PV 25 has a weight sum of 8.98883 and a merged classification
        \item Truth Vertex Weight Score:
        \begin{itemize}
            \item Truth PV 13: 2.51281076669693
            \item Truth PV 25: 5.682659983634949
            \item Truth PV 35: 0.3927352369064465
            \item Truth PV 45: 0.09750588983297348
            \item Truth PV 34: 0.3031177520751953
        \end{itemize}
        \item Division Results:
        \begin{itemize}
            \item PV 25 / Truth PV 13 = 0.279
            \item PV 25 / Truth PV 25 = 0.632 (barely misses 0.7)
            \item PV 25 / Truth PV 35 = 0.0437
            \item PV 25 / Truth PV 45 = 0.01
            \item PV 25 / Truth PV 34 = 0.034
        \end{itemize}
        \item Notice that Truth PV 13 and PV 25 score the highest
        \item Generally Method 1 results in the majority of PVs being merged
    \end{itemize}
\end{frame}

\begin{frame}{Case Study: Round on Example PV 25}
    \begin{itemize}
        \item PV 25 has a weight sum of 10 and a merged classification
        \begin{itemize}
            \item Truth PV 13: 4
            \item Truth PV 25: 6
            \item Truth PV 35: 0
            \item Truth PV 45: 0
            \item Truth PV 34: 0
        \end{itemize}
        \item Division Results:
        \begin{itemize}
            \item PV 25 / Truth PV 13 = 0.4
            \item PV 25 / Truth PV 25 = 0.6
            \item PV 25 / Truth PV 35 = 0
            \item PV 25 / Truth PV 45 = 0
            \item PV 25 / Truth PV 34 = 0
        \end{itemize}
    \end{itemize}
\end{frame}

\begin{frame}{Case Study: MaxScore on Example PV 25}
    \begin{itemize}
        \item PV 25 has a weight sum of 12 and a merged classification
        \begin{itemize}
            \item Truth PV 13: 4
            \item Truth PV 25: 6
            \item Truth PV 35: 1
            \item Truth PV 45: 0
            \item Truth PV 34: 1
        \end{itemize}
        \item Division Results:
        \begin{itemize}
            \item PV 25 / Truth PV 13 = 0.33
            \item PV 25 / Truth PV 25 = 0.5
            \item PV 25 / Truth PV 35 = 0.083
            \item PV 25 / Truth PV 45 = 0
            \item PV 25 / Truth PV 34 = 0.083
        \end{itemize}
    \end{itemize}
\end{frame}

\begin{frame}{Case Study: Results for Entire Event}
    \centering
    \begin{tabular}{c | c | c | c}
        \hline
        blank & Clean Rate & Merge Rate & Fake Rate \\
        \hline
        Method 1: None & 45.6\% & 54.4\% & 0\% \\
        Method 2: Round & 63.2\% & 36.8\% & 0\% \\
        Method 3: MaxScore & 64.9\% & 35.1\% & 0\% \\
        \hline
    \end{tabular}
\end{frame}

\begin{frame}{Concluding Remarks}
    \begin{itemize}
        \item Upcoming Work:
        \begin{itemize}
            \item Rewrite code from Jupyter Notebook to Python Scripts
            \item Finish up Vertex Classification metrics (haven't implemented split)
            \item \(p_T\) reweighting efficiency calculations
            \item Begin studies on using predicted PV location from PV-Finder
        \end{itemize}
        \item Unsure how to improve performance...
        \begin{itemize}
            \item Input Data?: \(d_0, z_0, \Sigma, \theta, \phi\)
            \item Label Data?: z location and histogram height
        \end{itemize}
        \item Have tested:
        \begin{itemize}
            \item Graph Attention Layers (performs same), Including PV x and y (performs slightly better),...
        \end{itemize}
    \end{itemize}
\end{frame}

\begin{frame}{Method 1 None: Probability Distribution}
    \begin{figure}
        \centering
        \includegraphics[width=0.65\linewidth]{may22nd/method1_probdist.png}
    \end{figure}
\end{frame}

\begin{frame}{Method 2 Round: Probability Distribution}
    \begin{figure}
        \centering
        \includegraphics[width=0.65\linewidth]{may22nd/method2_probdist.png}
    \end{figure}
\end{frame}

\begin{frame}{Method 3 MaxScore: Probability Distribution}
    \begin{figure}
        \centering
        \includegraphics[width=0.65\linewidth]{may22nd/method3_probdist.png}
    \end{figure}
\end{frame}

\begin{frame}[plain]
    \centering
    PV-Finder Meeting with Rocky/Lauren 06/06/2025
\end{frame}

\begin{frame}{Overview}
    \centering
    \textit{Putting it all together!}
    \begin{itemize}
        \item Current Studies:
        \begin{itemize}
            \item Primary Vertex Classification
            \item Application:
            \begin{itemize}
                \item Truth Primary Vertices
                \item \textbf{Reco Primary Vertices}
            \end{itemize}
            \item Generating Plots of Interest:
            \begin{itemize}
                \item Number of Reco PVs vs Number of Truth PVs
                \item Number of Reco PVs per Event vs Vertex Class at \(\left< \mu \right>=60\)
                \item Average Number of Track Associations vs \(\mu\)
            \end{itemize}
            \item \textbf{Example Merged/Split Cases}
            \item \textbf{Hard Scatter Classifications}
        \end{itemize}
    \end{itemize}
\end{frame}

% \begin{frame}{Recap: GNN Pipeline}
%     \begin{figure}
%         \centering
%         \includegraphics[width=0.9\linewidth]{april15th/endtoend_gnn_pipeline.png}
%     \end{figure}
% \end{frame}

% \begin{frame}{Recap: Assessing Performance}
%     \begin{itemize}
%         \item Track is associated with a primary vertex if
%         \begin{itemize}
%             \item Method 1: Score greater than 0.5
%             \begin{itemize}
%                 \item Note: can have harder cutoffs, like score greater than 0.6, 0.7, . . .
%             \end{itemize}
%             \item Method 2: Take the highest scoring edge for each track
%         \end{itemize}
%         \item Inclusivity of tracks within a certain range of PV
%         \begin{itemize}
%             \item \( \abs{z_{0, \text{track}} - z_{\text{pv}}} <= 5\text{mm} \text{ or } 10\text{mm}\)
%             \item \textbf{All Edges: Potential edge between all tracks and PVs}
%         \end{itemize}
%         \item \textbf{Efficiency: } (track is truth associated to the supplied vertex AND track passes GNN) / (track is truth associated to the supplied vertex)
%         \item \textbf{Impurity: } (track passes GNN AND is NOT associated to supplied vertex) / (track passes GNN)
%         \item Borrowed from \href{https://indico.cern.ch/event/1025152/contributions/4304219/attachments/2220591/3761247/210406_Basso_ATLAS-Weekly_TTVA-for-Run3.pdf}{Link}
%     \end{itemize}
% \end{frame}

% \begin{frame}[plain]
%     \centering
%     Vertex Classification
% \end{frame}

\begin{frame}{Re: Changes and Methods for Evaluation}
    \begin{itemize}
        \item Every track will now have a potential edge with every PV
        \begin{itemize}
            \item No more \( \abs{z_{0, \text{track}} - z_{\text{pv}}} <= 5\text{mm} \text{ or } 10\text{mm}\)
        \end{itemize}
        \item Outputs of the GNN can be interpreted as weights and evaluated in many ways
        \item Method 1 (None): No adjustment
        \begin{itemize}
            \item Will not do because evaluation time will become $\mathcal{O}(n_{\text{vertex}} \cdot n_{\text{tracks}})$
        \end{itemize}
        \item Method 2 (Threshold): Threshold of having a score be greater than 0.7, 0.85, or 0.9
        \item Method 3 (MaxScore): Highest scoring edge for each track is 1, rest is 0
    \end{itemize}
\end{frame}

\begin{frame}{Vertex Classification (1)}
    \begin{figure}
        \centering
        \includegraphics[width=0.7\linewidth]{june3rd/vertex_classification_1.png}
    \end{figure}
\end{frame}

\begin{frame}{Vertex Classification (2)}
    \begin{figure}
        \centering
        \includegraphics[width=0.7\linewidth]{june3rd/vertex_classification_2.png}
    \end{figure}
\end{frame}

\begin{frame}{Vertex Classification (3)}
    \begin{figure}
        \centering
        \includegraphics[width=0.7\linewidth]{june3rd/vertex_classification_3.png}
    \end{figure}
\end{frame}

\begin{frame}{Vertex Classification (4)}
    \begin{figure}
        \centering
        \includegraphics[width=0.7\linewidth]{june3rd/vertex_classification_4.png}
    \end{figure}
\end{frame}

\begin{frame}{Vertex Classification (5)}
    \begin{figure}
        \centering
        \includegraphics[width=0.7\linewidth]{june3rd/vertex_classification_5.png}
    \end{figure}
\end{frame}

\begin{frame}{Vertex Classification (6)}
    \begin{figure}
        \centering
        \includegraphics[width=0.7\linewidth]{june3rd/vertex_classification_6.png}
    \end{figure}
\end{frame}

\begin{frame}{Vertex Classification (7)}
    \begin{figure}
        \centering
        \includegraphics[width=0.7\linewidth]{june3rd/vertex_classification_7.png}
    \end{figure}
\end{frame}

% \begin{frame}{Average Number of Truth Vertices per \(\left< \mu \right>\)}

% \end{frame}

\begin{frame}[plain]
    \centering
    Vertex Classification on Reco Vertex
\end{frame}

\begin{frame}{Start Connecting the Pipelines}
    \begin{figure}
        \centering
        \includegraphics[width=0.55\linewidth]{april4th/pipeline_pvfinder_gnn.png}
    \end{figure}
\end{frame}

\begin{frame}{Steps Outlined}
    \begin{itemize}
        \item Train PV-Finder
        \begin{itemize}
            \item Using current best PV-Finder run (UNet with Dropout Model)
            \item (Old Metrics) Efficiency: 94.56\%
            \item (Old Metrics) FPR: 0.72
        \end{itemize}
        \item Calculate predicted primary vertex location
        \begin{itemize}
            \item Current peak finding algorithm (\href{https://gitlab.cern.ch/qlei/atlas_pvfinder/-/blob/master/model/efficiency_res_optimized_atlas.py?ref_type=heads}{Link})
            \item scipy.find\_peaks (\href{https://docs.scipy.org/doc/scipy/reference/generated/scipy.signal.find_peaks.html}{Link})
        \end{itemize}
        \item Construct graphs from reconstructed PV information
        \item Evaluate using GNN model trained on truth information
        \begin{itemize}
            \item Using current best GNN run (GraphConv Model)
            \item Achieves results shown in prior slides
        \end{itemize}
    \end{itemize}
\end{frame}

\begin{frame}{Reco Vertices with Classification at \(\left< \mu \right> =60\) (Thr=0.95)}
    \begin{columns}
        \begin{column}{0.55\textwidth}
            \begin{itemize}
                \item Total Reconstructed PVs: 62457
                \item Total Clean: 49500, Rate: 0.7925
                \item Total Merged: 8243, Rate: 0.1320
                \item Total Split: 3838, Rate: 0.0615
                \item Total Fake: 876, Rate: 0.0140
            \end{itemize}
        \end{column}

        \begin{column}{0.45\textwidth}
            \begin{figure}
                \centering
                \includegraphics[width=1\linewidth]{june5th/reco_pv_avg_bar_chart_thr95_recovertex.png}
            \end{figure}
        \end{column}
    \end{columns}
\end{frame}

% \begin{frame}{Reco Vertices with Classification at \(\left< \mu \right> =60\) MaxScore}
%     \begin{columns}
%         \begin{column}{0.55\textwidth}
%             \begin{itemize}
%                 \item Total Reconstructed PVs: 62457
%                 \item Total Clean: 49500, Rate: 0.7925
%                 \item Total Merged: 8243, Rate: 0.1320
%                 \item Total Split: 3838, Rate: 0.0615
%                 \item Total Fake: 876, Rate: 0.0140
%             \end{itemize}
%         \end{column}

%         \begin{column}{0.45\textwidth}
%             \begin{figure}
%                 \centering
%                 \includegraphics[width=1\linewidth]{example-image-a}
%             \end{figure}
%         \end{column}
%     \end{columns}
% \end{frame}

\begin{frame}{Comparison of Reco Vertex Classification: Thresholds}
    \begin{columns}
        \begin{column}{0.5\textwidth}
            \begin{figure}
                \centering
                \includegraphics[width=1\linewidth]{june5th/stacked_pv_classification_5_80_thr9_recovertex.png}
                \caption{Thr = 0.9}
            \end{figure}
        \end{column}

        \begin{column}{0.5\textwidth}
            \begin{figure}
                \centering
                \includegraphics[width=1\linewidth]{june5th/stacked_pv_classification_5_80_thr85_recovertex.png}
                \caption{Thr = 0.85}
            \end{figure}
        \end{column}
    \end{columns}
\end{frame}

\begin{frame}{Comparison of Reco Vertex Classification: Thresholds}
    \begin{columns}
        \begin{column}{0.5\textwidth}
            \begin{figure}
                \centering
                \includegraphics[width=1\linewidth]{june5th/stacked_pv_classification_5_80_thr7_recovertex.png}
                \caption{Thr = 0.7}
            \end{figure}
        \end{column}

        \begin{column}{0.5\textwidth}
            \begin{figure}
                \centering
                \includegraphics[width=1\linewidth]{june5th/stacked_pv_classification_5_80_thr95_recovertex.png}
                \caption{Thr = 0.95}
            \end{figure}
        \end{column}
    \end{columns}
\end{frame}

% \begin{frame}{Comparison of Reco Vertex Classification: MaxScore}
%     \begin{columns}
%         \begin{column}{0.5\textwidth}
%             \begin{figure}
%                 \centering
%                 \includegraphics[width=1\linewidth]{example-image-a}
%                 \caption{MaxScore}
%             \end{figure}
%         \end{column}

%         \begin{column}{0.5\textwidth}
%             \begin{figure}
%                 \centering
%                 \includegraphics[width=1\linewidth]{june5th/stacked_pv_classification_5_80_thr95_recovertex.png}
%                 \caption{Thr = 0.95}
%             \end{figure}
%         \end{column}
%     \end{columns}
% \end{frame}

\begin{frame}{A Look into GNN Associations and Thresholds}
    \begin{itemize}
        \item Classification rates look fine... but how do the track-vertex associations look?
        \begin{itemize}
            \item What is the "quality" of reconstructed vertices?
            \item How is the "quality" impacted by the threshold?
        \end{itemize}
        \item What is causing the increase in merged/split vertices?
        \item Are there tracks that are not being used due to threshold?
        \item Lets look at test cases!
    \end{itemize}
    \begin{figure}
        \centering
        \includegraphics[width=0.4\linewidth]{this_is_fine.jpg}
    \end{figure}
\end{frame}

\begin{frame}{Case Study: Looking at Merged/Split Reco PVs}
    \begin{columns}
        \begin{column}{0.4\textwidth}
            \begin{itemize}
                \item Total Reco PVs: 42
                \item Total Truth PVs: 59
                \item Thr=0.85
                \begin{itemize}
                    \item Num Clean: 22
                    \item Num Merged: 12
                    \item Num Split: 8
                    \item Num Fake: 0
                \end{itemize}
            \end{itemize}
        \end{column}

        \begin{column}{0.6\textwidth}
            \begin{figure}
                \centering
                \includegraphics[width=1\linewidth]{june5th/casestudy_example_reco_truth_hist_dist.png}
            \end{figure}
        \end{column}
    \end{columns}
\end{frame}

\begin{frame}{Case Study: Two Overlapping Vertices}
    \begin{columns}
        \begin{column}{0.4\textwidth}
            \begin{itemize}
                \item Thr=0.85
                \item Reco PV 26: 59 Assoc tracks
                \begin{itemize}
                    \item \(W_{sum} = 59\)
                    \item Missing Track 168
                    \item Output Probs: 0.834
                \end{itemize}
                \item Truth PV 16: 15 Assoc tracks
                \begin{itemize}
                    \item \(W_{1} = 41\)
                    \item \(W_{1}/W_{sum} = 0.69\)
                \end{itemize}
                \item Truth PV 39: 45 Assoc tracks
                \begin{itemize}
                    \item \(W_{2} = 14\)
                    \item \(W_{2}/W_{sum} = 0.24\)
                \end{itemize}
                \item Classification: Merged
            \end{itemize}
        \end{column}

        \begin{column}{0.6\textwidth}
            \begin{figure}
                \centering
                \includegraphics[width=1\linewidth]{june5th/casestudy_example_reco_truth_hist_dist_16_39.png}
            \end{figure}
        \end{column}
    \end{columns}
\end{frame}

\begin{frame}{Case Study: Two Overlapping Vertices}
    \begin{columns}
        \begin{column}{0.4\textwidth}
            \begin{itemize}
                \item Thr=0.95
                \item Reco PV 26: 51 Assoc tracks
                \begin{itemize}
                    \item \(W_{sum} = 51\)
                \end{itemize}
                \item Truth PV 16: 15 Assoc tracks
                \begin{itemize}
                    \item \(W_{1} = 13\)
                    \item \(W_{1}/W_{sum} = 0.25\)
                \end{itemize}
                \item Truth PV 39: 45 Assoc tracks
                \begin{itemize}
                    \item \(W_{2} = 38\)
                    \item \(W_{2}/W_{sum} = 0.74\)
                \end{itemize}
                \item Classification: Clean
            \end{itemize}
        \end{column}

        \begin{column}{0.6\textwidth}
            \begin{figure}
                \centering
                \includegraphics[width=1\linewidth]{june5th/casestudy_example_reco_truth_hist_dist_16_39.png}
            \end{figure}
        \end{column}
    \end{columns}
\end{frame}

\begin{frame}{Case Study: Three Nearby Vertices (1)}
    \begin{columns}
        \begin{column}{0.4\textwidth}
            \begin{itemize}
                \item Thr=0.85
                \item Reco PV 24: 33 Assoc tracks
                \begin{itemize}
                    \item \(\sum p_T^2 = 4e7\)
                \end{itemize}
                \item Truth PV 31: 13 Assoc tracks
                \begin{itemize}
                    \item \(W_{1} = 12\)
                    \item \(W_{1}/W_{sum} = 0.4\)
                \end{itemize}
                \item Truth PV 38: 10 Assoc tracks
                \begin{itemize}
                    \item \(W_{2} = 10\)
                    \item \(W_{2}/W_{sum} = 0.3\)
                \end{itemize}
                \item Truth PV 51: 6 Assoc tracks
                \begin{itemize}
                    \item \(W_{3} = 5\)
                    \item \(W_{3}/W_{sum} = 0.15\)
                \end{itemize}
                \item Classification: \textbf{Split}
            \end{itemize}
        \end{column}

        \begin{column}{0.6\textwidth}
            \begin{figure}
                \centering
                \includegraphics[width=1\linewidth]{june5th/casestudy_example_reco_truth_hist_dist_31_38_51_1.png}
            \end{figure}
        \end{column}
    \end{columns}
\end{frame}

\begin{frame}{Case Study: Three Nearby Vertices (2)}
    \begin{columns}
        \begin{column}{0.4\textwidth}
            \begin{itemize}
                \item Thr=0.85
                \item Reco PV 25: 42 Assoc tracks
                \begin{itemize}
                    \item \(\sum p_T^2 = 6e7\)
                \end{itemize}
                \item Truth PV 31: 13 Assoc tracks
                \begin{itemize}
                    \item \(W_{1} = 12\)
                    \item \(W_{1}/W_{sum} = 0.4\)
                \end{itemize}
                \item Truth PV 38: 10 Assoc tracks
                \begin{itemize}
                    \item \(W_{2} = 10\)
                    \item \(W_{2}/W_{sum} = 0.3\)
                \end{itemize}
                \item Truth PV 51: 6 Assoc tracks
                \begin{itemize}
                    \item \(W_{3} = 5\)
                    \item \(W_{3}/W_{sum} = 0.15\)
                \end{itemize}
                \item Classification: \textbf{Merged}
            \end{itemize}
        \end{column}

        \begin{column}{0.6\textwidth}
            \begin{figure}
                \centering
                \includegraphics[width=1\linewidth]{june5th/casestudy_example_reco_truth_hist_dist_31_38_51_2.png}
            \end{figure}
        \end{column}
    \end{columns}
\end{frame}

\begin{frame}{Case Study: Three Nearby Vertices (3)}
    \begin{columns}
        \begin{column}{0.4\textwidth}
            \begin{itemize}
                \item Thr=0.95
                \item Reco PV 24: 24 Assoc tracks
                \begin{itemize}
                    \item \(\sum p_T^2 = 4e7\)
                \end{itemize}
                \item Truth PV 31: 13 Assoc tracks
                \begin{itemize}
                    \item \(W_{1} = 12\)
                    \item \(W_{1}/W_{sum} = 0.5\)
                \end{itemize}
                \item Truth PV 38: 10 Assoc tracks
                \begin{itemize}
                    \item \(W_{2} = 6\)
                    \item \(W_{2}/W_{sum} = 0.25\)
                \end{itemize}
                \item Truth PV 51: 5 Assoc tracks
                \begin{itemize}
                    \item \(W_{3} = 5\)
                    \item \(W_{3}/W_{sum} = 0.21\)
                \end{itemize}
                \item Classification: \textbf{Merged}
            \end{itemize}
        \end{column}

        \begin{column}{0.6\textwidth}
            \begin{figure}
                \centering
                \includegraphics[width=1\linewidth]{june5th/casestudy_example_reco_truth_hist_dist_31_38_51_1.png}
            \end{figure}
        \end{column}
    \end{columns}
\end{frame}

\begin{frame}{Case Study: Three Nearby Vertices (4)}
    \begin{columns}
        \begin{column}{0.4\textwidth}
            \begin{itemize}
                \item Thr=0.95
                \item Reco PV 25: 22 Assoc tracks
                \begin{itemize}
                    \item \(\sum p_T^2 = 3e7\)
                \end{itemize}
                \item Truth PV 31: 13 Assoc tracks
                \begin{itemize}
                    \item \(W_{1} = 0.4\)
                    \item \(W_{1}/W_{sum} = 0.4\)
                \end{itemize}
                \item Truth PV 38: 10 Assoc tracks
                \begin{itemize}
                    \item \(W_{2} = 0.4\)
                    \item \(W_{2}/W_{sum} = 0.3\)
                \end{itemize}
                \item Classification: \textbf{Merged}
            \end{itemize}
        \end{column}

        \begin{column}{0.6\textwidth}
            \begin{figure}
                \centering
                \includegraphics[width=1\linewidth]{june5th/casestudy_example_reco_truth_hist_dist_31_38_51_2.png}
            \end{figure}
        \end{column}
    \end{columns}
\end{frame}

\begin{frame}{Number of Associations Per Tracks Varying Threshold}
    \begin{columns}
        \begin{column}{0.5\textwidth}
            \begin{figure}
                \centering
                \includegraphics[width=1\linewidth]{june5th/casestudy_num_track_assocs_thr85.png}
                \caption{Thr = 0.85}
            \end{figure}
        \end{column}

        \begin{column}{0.5\textwidth}
            \begin{figure}
                \centering
                \includegraphics[width=1\linewidth]{june5th/casestudy_num_track_assocs_thr95.png}
                \caption{Thr = 0.95}
            \end{figure}
        \end{column}
    \end{columns}
\end{frame}

\begin{frame}{Number of Associations Per Track vs \( \mu \) Varying Threshold}
    \begin{columns}
        \begin{column}{0.5\textwidth}
            \begin{figure}
                \centering
                \includegraphics[width=1\linewidth]{june5th/num_track_assoc_vs_mu_thr85.png}
                \caption{Thr = 0.85}
            \end{figure}
        \end{column}

        \begin{column}{0.5\textwidth}
            \begin{figure}
                \centering
                \includegraphics[width=1\linewidth]{june5th/num_track_assoc_vs_mu_thr95.png}
                \caption{Thr = 0.95}
            \end{figure}
        \end{column}
    \end{columns}
\end{frame}

\begin{frame}{Re: A Look into GNN Associations and Thresholds}
    \begin{itemize}
        \item What is causing the increase in merged/split vertices?
        \begin{itemize}
            \item Seems to be that distinguishing tracks for nearby vertices are tough
        \end{itemize}
        \item Are there tracks that are not being used due to threshold?
        \begin{itemize}
            \item Tracks are being used even in high thresholds
        \end{itemize}
        \item Are there other ways on evaluating performance?
        \begin{itemize}
            \item \textbf{Hard Scatter Vertex Classification}
        \end{itemize}
    \end{itemize}
    \begin{figure}
        \centering
        \includegraphics[width=0.4\linewidth]{this_is_fine.jpg}
    \end{figure}
\end{frame}

\begin{frame}{Concluding Remarks}
    \begin{itemize}
        \item How to continue improving performance...? What is the next direction...?
        \begin{itemize}
            \item Start applying cuts to the output of GNN (?)
        \end{itemize}
        \item Developing tools to continue thoroughly analyze PV-Finder + GNN outputs
        \begin{itemize}
            \item Work on HS Vertex Classification and Metrics
            \item TTVA Efficiency Calculations on HS Vertices
        \end{itemize}
        \item Muon Collider Studies
        \begin{itemize}
            \item Track Selection Studies
            \item Jet \(p_T\) Studies
            \item MUSIC Detector
        \end{itemize}
        \item Upcoming Presentations:
        \begin{itemize}
            \item Biweekly LHCb Meeting Moved to Wednesday June 11th
            \item TVPD (?)
        \end{itemize}
    \end{itemize}

    \centering
    \textit{Fin.}
\end{frame}


\begin{frame}[plain]
    \centering
    Vertex Classification on Truth Vertex
\end{frame}

\begin{frame}{Truth Vertices with Classification at \(\left< \mu \right> =60\) (Thr=0.85)}
    \begin{columns}
        \begin{column}{0.55\textwidth}
            \begin{itemize}
                \item Total Reconstructed PVs: 80498
                \item Total Clean: 37428, Rate: 0.4650
                \item Total Merged: 19552, Rate: 0.2429
                \item Total Split: 23508, Rate: 0.2920
                \item Total Fake: 10, Rate: 0.0001
                \item Lots of merged and split but this is evaluating on truth primary vertices which include PVs that would never be reconstructed
            \end{itemize}
        \end{column}

        \begin{column}{0.45\textwidth}
            \begin{figure}
                \centering
                \includegraphics[width=1\linewidth]{june5th/reco_pv_avg_bar_chart_thr85.png}
            \end{figure}
        \end{column}
    \end{columns}
\end{frame}

\begin{frame}{Truth Vertices with Classification at \(\left< \mu \right> =60\) (Thr=0.9)}
    \begin{columns}
        \begin{column}{0.55\textwidth}
            \begin{itemize}
                \item Total Reconstructed PVs: 80498
                \item Total Clean: 41588, Rate: 0.5166
                \item Total Merged: 17989, Rate: 0.2235
                \item Total Split: 20913, Rate: 0.2598
                \item Total Fake: 8, Rate: 0.0001
            \end{itemize}
        \end{column}

        \begin{column}{0.45\textwidth}
            \begin{figure}
                \centering
                \includegraphics[width=1\linewidth]{june5th/reco_pv_avg_bar_chart_thr9.png}
            \end{figure}
        \end{column}
    \end{columns}
\end{frame}

\begin{frame}{Comparison of Truth Vertex Classification: Thresholds}
    \begin{columns}
        \begin{column}{0.5\textwidth}
            \begin{figure}
                \centering
                \includegraphics[width=1\linewidth]{june5th/stacked_pv_classification_5_80_thr9.png}
                \caption{Thr = 0.9}
            \end{figure}
        \end{column}

        \begin{column}{0.5\textwidth}
            \begin{figure}
                \centering
                \includegraphics[width=1\linewidth]{june5th/stacked_pv_classification_5_80_thr85.png}
                \caption{Thr = 0.85}
            \end{figure}
        \end{column}
    \end{columns}
\end{frame}

\begin{frame}{Comparison of Truth Vertex Classification: Threshold}
    \begin{columns}
        \begin{column}{0.5\textwidth}
            \begin{figure}
                \centering
                \includegraphics[width=1\linewidth]{june5th/stacked_pv_classification_5_80_thr7.png}
                \caption{Thr = 0.7}
            \end{figure}
        \end{column}

        \begin{column}{0.5\textwidth}
            \begin{figure}
                \centering
                \includegraphics[width=1\linewidth]{example-image-a}
                \caption{MaxScore (Work in Progress, my algo is slow)}
            \end{figure}
        \end{column}
    \end{columns}
\end{frame}

\begin{frame}{Truth Vertices with Classification at \(\left< \mu \right> =60\) (Thr=0.7)}
    \begin{columns}
        \begin{column}{0.55\textwidth}
            \begin{itemize}
                \item Total Reconstructed PVs: 80498
                \item Total Clean: 25016, Rate: 0.3108
                \item Total Merged: 26647, Rate: 0.3310
                \item Total Split: 28824, Rate: 0.3581
                \item Total Fake: 11, Rate: 0.0001
            \end{itemize}
        \end{column}

        \begin{column}{0.45\textwidth}
            \begin{figure}
                \centering
                \includegraphics[width=1\linewidth]{june5th/reco_pv_avg_bar_chart_thr7.png}
            \end{figure}
        \end{column}
    \end{columns}
\end{frame}

\begin{frame}{Proper Definitions of Different HS Classifications}
    \begin{itemize}
        \item Clean/Matched: "event contains a CLEAN reconstructed vertex originating from true HS interaction, AND the HS interaction does not contribute more than 50\% of the accumulated track weight to any other vertex"
        \item LowPU: "event contains a MERGED vertex with at least 50\% of the accumulated track weight coming from the simulated HS interaction"
        \item HighPU: "event contains a MERGED vertex with its main contribution coming from a simulated pile-up interaction, and in which the HS interaction contributes between 1\% and 50\% of the accumulated track weight"
        \item PurePU: "event does not contain any reconstructed vertex with at least 1\% accumulated track weight from HS interaction"
    \end{itemize}
\end{frame}

\begin{frame}[plain]
    \centering
    PV-Finder Meeting with Rocky/Lauren 06/25/2025
\end{frame}


\begin{frame}{Overview Musings for PV-Finder}
    \begin{itemize}
        \item (In Progress) GNN Training on Reco Vertex
        \item \textbf{HS Vertex Quality and Metrics}
        \item \textbf{SLAC ATLAS Group Meeting Friday (06/24), 8:30am}
        \begin{itemize}
            \item Send draft by end of day
            \item Fix plots from prior TVPD presentation
        \end{itemize}
        \item Note: Couple days where sneezy was down last week
    \end{itemize}
\end{frame}

\begin{frame}{Recap: HS Vertex Quality}
    \begin{itemize}
        \item HS Vertex Quality Classification:
        \begin{itemize}
            \item \textbf{Clean/Matched}: \textbf{CLEAN} reco vertex matching true HS interaction AND \textbf{track weights don't contribute \(> 50 \%\) to other vertices}
            \item \textbf{LowPU}: \textbf{MERGED} reco vertex with \textbf{at least 50\%} of the accumulated \textbf{track weight} coming \textbf{from the simulated HS interaction}
            \item \textbf{HighPU}: \textbf{MERGED} reco vertex the HS interaction contributing between \textbf{1\% and 50\% of the accumulated track weight}"
            \item \textbf{PurePU}: No reco vertex with at least 1\% accumulated track weight from HS interaction"
        \end{itemize}
        \item Edge Case:
        \begin{itemize}
            \item What to do when, there is a split reco vertex?
            \item A reco vertex that correctly matches 60 tracks from HS vertex (Clean)
            \item Another reco vertex that correctly matches 59 tracks from HS vertex (Split)
            \item Due to this, have 67 events "unclassified" in 2550 total events
        \end{itemize}
    \end{itemize}
\end{frame}

\begin{frame}{HS Vertex Quality Category Plot}
    \begin{columns}
        \begin{column}{0.5\textwidth}
            \begin{figure}
                \centering
                \includegraphics[width=1\linewidth]{june25th/fig_07_amvf.png}
                \caption{AMVF Pub Note: \href{https://cds.cern.ch/record/2670380/files/ATL-PHYS-PUB-2019-015.pdf}{Link}}
            \end{figure}
        \end{column}

        \begin{column}{0.5\textwidth}
            \begin{figure}
                \centering
                \includegraphics[width=1\linewidth]{june25th/rhs_vertex_quality_category.png}
                \caption{Thr = 0.95}
            \end{figure}
        \end{column}
    \end{columns}
\end{frame}

\begin{frame}{HS Vertex Reconstruction Efficiency}
    \begin{columns}
        \begin{column}{0.3\textwidth}
            \begin{itemize}
                \item Fraction of events which the HS vertex is reconstructed and classified as Clean/Matched, LowPU, or HighPu
                \item Local pile-up density is the number of vertices within a \(\pm\) 2mm window of a HS vertex
            \end{itemize}
        \end{column}

        \begin{column}{0.7\textwidth}
            \begin{figure}
                \centering
                \includegraphics[width=0.85\linewidth]{june25th/vertex_reconstruction_efficiency_vs_local_pu_density.png}
                \caption{Thr = 0.95}
            \end{figure}
        \end{column}
    \end{columns}
\end{frame}

\begin{frame}{Comparison: HS Vertex Reconstruction Efficiency}
    \begin{columns}
        \begin{column}{0.5\textwidth}
            \begin{figure}
                \centering
                \includegraphics[width=1\linewidth]{june25th/fig_08a_amvf_reco.png}
                \caption{AMVF Pub Note (Fig. 8): \href{https://cds.cern.ch/record/2670380/files/ATL-PHYS-PUB-2019-015.pdf}{Link}}
            \end{figure}
        \end{column}

        \begin{column}{0.5\textwidth}
            \begin{figure}
                \centering
                \includegraphics[width=1\linewidth]{june25th/vertex_reconstruction_efficiency_vs_local_pu_density.png}
                \caption{Thr = 0.95}
            \end{figure}
        \end{column}
    \end{columns}
\end{frame}

\begin{frame}{HS Vertex Selection Efficiency}
    \begin{columns}
        \begin{column}{0.3\textwidth}
            \begin{itemize}
                \item Fraction of events where the reconstructed vertex with highest \(\sum p_T^2\) is the one contianing the largest total weight from true HS tracks
            \end{itemize}
        \end{column}

        \begin{column}{0.7\textwidth}
            \begin{figure}
                \centering
                \includegraphics[width=0.85\linewidth]{june25th/vertex_selection_efficiency_vs_local_pu_density.png}
                \caption{Thr = 0.95}
            \end{figure}
        \end{column}
    \end{columns}
\end{frame}

\begin{frame}{Comparison: HS Vertex Selection Efficiency}
    \begin{columns}
        \begin{column}{0.5\textwidth}
            \begin{figure}
                \centering
                \includegraphics[width=1\linewidth]{june25th/fig_09a_amvf_selection.png}
                \caption{AMVF Pub Note (Fig. 9): \href{https://cds.cern.ch/record/2670380/files/ATL-PHYS-PUB-2019-015.pdf}{Link}}
            \end{figure}
        \end{column}

        \begin{column}{0.5\textwidth}
            \begin{figure}
                \centering
                \includegraphics[width=1\linewidth]{june25th/vertex_selection_efficiency_vs_local_pu_density.png}
                \caption{Thr = 0.95}
            \end{figure}
        \end{column}
    \end{columns}
\end{frame}

\begin{frame}[plain]
    \centering
    SLAC ATLAS Group Meeting (06/26)
\end{frame}

\begin{frame}{Current Primary Vertex Algorithm: AVMF}
    % Talk about innate challenges that comes with using AVMF
    \begin{columns}
        % Text
        \column{0.4\linewidth}
        \begin{itemize}
            \item Global approach to vertex finding and fitting
            \item Finds and fits tracks based on compatibility to different primary vertices
            \item Limitations
            \begin{enumerate}
                \item Parallelization
                \item Time scale increases non linearly with number of PVs
            \end{enumerate}
        \end{itemize}

        % Text
        \column{0.6\linewidth}
        \begin{figure}
            \centering
            \includegraphics[width=0.9\linewidth]{oct25th/AVMF.png}
            \label{fig:AVMF}
            \caption{\href{https://cds.cern.ch/record/2670380/files/ATL-PHYS-PUB-2019-015.pdf}{[ATL-PHYS-PUB-2019-015]}}
        \end{figure}
    \end{columns}
\end{frame}

\begin{frame}{PV-Finder Before}
    \vspace{5pt}
    \centering
    This network performs nicely but is limited by execution time
    \vspace{25pt}
    \begin{figure}
        \centering
        \includegraphics[width=1\linewidth]{APS2025/analyticalkde_pv.png}
        \label{fig:oldPVFinderArchitecture}
    \end{figure}
\end{frame}

\begin{frame}{PV-Finder After}
    \centering
    TVPD presentation on PV-Finder results has more details! (\href{https://indico.cern.ch/event/1534661/}{Link})
    \begin{figure}
        \centering
        \includegraphics[width=1\linewidth]{april4th/updated_tracks_to_pv.png}
        \label{fig:oldPVFinderArchitecture}
    \end{figure}
\end{frame}

\begin{frame}{Example End-to-end tracks-to-hist Output}
    \begin{figure}
        \centering
        \includegraphics[width=0.8\linewidth]{june25th/example_end_to_end_with_class.png}
        \label{fig:example_Hists_output_1}
    \end{figure}
\end{frame}

\begin{frame}{\(\Delta z\) Resolution and \(\sigma_{\text{vtx-vtx}}\) Calculation}
    \begin{columns}
        \begin{column}{0.6\linewidth}
            \begin{figure}
                \centering
                \includegraphics[width=0.9\linewidth]{april4th/deltaz(2).png}
            \end{figure}
        \end{column}

        \begin{column}{0.4\linewidth}
            \begin{figure}
                \centering
                \includegraphics[width=1\linewidth]{april4th/sigma_vtx_vtx_dropout.png}
            \end{figure}
        \end{column}
    \end{columns}
\end{frame}

\begin{frame}{UNet vs UNet++ Efficiency and FPR with MLP}
    \centering
    UNet and UNet++ perform basically identically!
    \begin{figure}
        \centering
        \includegraphics[width=1\linewidth]{april4th/droput_unet_unetplusplus_eff.png}
    \end{figure}
\end{frame}

\begin{frame}{Intermission}
    \begin{itemize}
        \item Great that PV-Finder can reconstruct more clean/merged PV!
        \item Tracks to Vertex Association?
        \begin{itemize}
            \item Necessary for hard-scatter vertex calculations
            \item Allows for proper comparison with AMVF
        \end{itemize}
        \item \textbf{GNN for TTVA}
        \begin{itemize}
            \item Graph Construction
            \item AMVF-Stype Vertex Classification
            \item Preliminary Performance
            \item Case Study
        \end{itemize}
    \end{itemize}
\end{frame}

\begin{frame}{PV-Finder + GNN Fitter}
    \begin{figure}
        \centering
        \includegraphics[width=0.55\linewidth]{april4th/pipeline_pvfinder_gnn.png}
    \end{figure}
\end{frame}

\begin{frame}{GNN Concept}
    \begin{figure}
        \centering
        \includegraphics[width=0.55\linewidth]{april4th/gnn_concept.png}
    \end{figure}
\end{frame}

\begin{frame}{Forming Graphs}
    \begin{figure}
        \centering
        \includegraphics[width=0.6\linewidth]{april4th/forming_edges_gnn.png}
    \end{figure}
\end{frame}

\begin{frame}{Scoring Edges in Graphs}
    \begin{figure}
        \centering
        \includegraphics[width=0.6\linewidth]{april4th/scoring_edges_gnn.png}
    \end{figure}
\end{frame}

\begin{frame}{Choosing Edges in Graphs}
    \begin{figure}
        \centering
        \includegraphics[width=0.6\linewidth]{april4th/choosing_edges_gnn.png}
    \end{figure}
\end{frame}

\begin{frame}{GNN Pipeline}
    \begin{figure}
        \centering
        \includegraphics[width=0.9\linewidth]{april15th/endtoend_gnn_pipeline.png}
    \end{figure}
\end{frame}

\begin{frame}[plain]
    \centering
    AMVF-Style Vertex Classification
\end{frame}

\begin{frame}{Vertex Classification (1)}
    \begin{figure}
        \centering
        \includegraphics[width=0.7\linewidth]{june20th/reco_vertex_classification_1.png}
    \end{figure}
\end{frame}

\begin{frame}{Vertex Classification (2)}
    \begin{figure}
        \centering
        \includegraphics[width=0.7\linewidth]{june20th/reco_vertex_classification_2.png}
    \end{figure}
\end{frame}

\begin{frame}{Vertex Classification (3)}
    \begin{figure}
        \centering
        \includegraphics[width=0.7\linewidth]{june20th/reco_vertex_classification_3.png}
    \end{figure}
\end{frame}

\begin{frame}{Vertex Classification (4)}
    \begin{figure}
        \centering
        \includegraphics[width=0.7\linewidth]{june20th/reco_vertex_classification_4.png}
    \end{figure}
\end{frame}

\begin{frame}{Vertex Classification (5)}
    \begin{figure}
        \centering
        \includegraphics[width=0.7\linewidth]{june20th/reco_vertex_classification_5.png}
    \end{figure}
\end{frame}

\begin{frame}{Vertex Classification (6)}
    \begin{figure}
        \centering
        \includegraphics[width=0.7\linewidth]{june20th/reco_vertex_classification_6.png}
    \end{figure}
\end{frame}

\begin{frame}{Vertex Classification (7)}
    \begin{figure}
        \centering
        \includegraphics[width=0.7\linewidth]{june20th/reco_vertex_classification_7.png}
    \end{figure}
\end{frame}

\begin{frame}{Vertex Classification (8)}
    \begin{figure}
        \centering
        \includegraphics[width=0.7\linewidth]{june25th/reco_vertex_classification_8.png}
    \end{figure}
\end{frame}

\begin{frame}[plain]
    \centering
    Evaluating Reco Vertices with AMVF-Style Vertex Classification
\end{frame}

\begin{frame}{Steps Outlined}
    \begin{itemize}
        \item Train PV-Finder
        \begin{itemize}
            \item Using current best PV-Finder run (UNet with Dropout Model)
            \item (Old Metrics) Efficiency: 94.56\%
            \item (Old Metrics) FPR/Average Number of Fake Reco PV: 0.72
        \end{itemize}
        \item Calculate predicted primary vertex location
        \item Construct graphs from reconstructed PV information
        \item Evaluate using GNN model trained on truth information
        \begin{itemize}
            \item Using current best GNN run (GraphConv Model)
            \item Achieves results shown in prior slides
        \end{itemize}
    \end{itemize}
\end{frame}

% \begin{frame}{Re: Changes and Methods for Evaluation}
%     \begin{itemize}
%         \item Every track will now have a potential edge with every PV
%         \begin{itemize}
%             \item \textbf{No more} \( \abs{z_{0, \text{track}} - z_{\text{pv}}} <= 5\text{mm} \text{ or } 10\text{mm}\)
%         \end{itemize}
%         \item Outputs of the GNN can be interpreted as weights and evaluated in many ways
%         \item Method 1 (Threshold): Threshold of having a score be greater than 0.7, 0.85, or 0.9
%         \item Method 2 (MaxScore): Highest scoring edge for each track is 1, rest is 0
%     \end{itemize}
% \end{frame}

\begin{frame}{Reco Vertices with Classification at \(\left< \mu \right> =60\) (Thr=0.95)}
    \begin{columns}
        \begin{column}{0.5\textwidth}
            \begin{itemize}
                \item Total Reconstructed PVs: 62457
                \item Total Truth PVs: 80498
                \item Total Clean: 42880, Rate: 68.66\%
                \item Total Merged: 13437, Rate: 21.51\%
                \item Total Split: 3838, Rate: 9.83\%
                \item Total Fake: 3, Rate: 0\%
            \end{itemize}
        \end{column}

        \begin{column}{0.5\textwidth}
            \begin{figure}
                \centering
                \includegraphics[width=1\linewidth]{june25th/reco_pv_avg_bar_chart_95thr_recovertex.png}
                \caption{Number of reco vertices at \(\left< \mu\right> = 60\)}
            \end{figure}
        \end{column}
    \end{columns}
\end{frame}

% \begin{frame}{Reco Vertices with Classification at \(\left< \mu \right> =60\) MaxScore}
%     \begin{columns}
%         \begin{column}{0.55\textwidth}
%             \begin{itemize}
%                 \item Total Reconstructed PVs: 62457
%                 \item Total Clean: 49500, Rate: 0.7925
%                 \item Total Merged: 8243, Rate: 0.1320
%                 \item Total Split: 3838, Rate: 0.0615
%                 \item Total Fake: 876, Rate: 0.0140
%             \end{itemize}
%         \end{column}

%         \begin{column}{0.45\textwidth}
%             \begin{figure}
%                 \centering
%                 \includegraphics[width=1\linewidth]{example-image-a}
%             \end{figure}
%         \end{column}
%     \end{columns}
% \end{frame}

\begin{frame}{Comparison of Reco Vertex Classification: Thresholds (1)}
    \begin{columns}
        \begin{column}{0.5\textwidth}
            \begin{figure}
                \centering
                \includegraphics[width=1\linewidth]{june25th/stacked_pv_classification_5_80_85thr_recovertex.png}
                \caption{Thr = 0.85}
            \end{figure}
        \end{column}

        \begin{column}{0.5\textwidth}
            \begin{figure}
                \centering
                \includegraphics[width=1\linewidth]{june25th/stacked_pv_classification_5_80_95thr_recovertex.png}
                \caption{Thr = 0.95}
            \end{figure}
        \end{column}
    \end{columns}
\end{frame}

\begin{frame}{Comparison of Reco Vertex Classification: Thresholds (2)}
    \begin{columns}
        \begin{column}{0.5\textwidth}
            \begin{figure}
                \centering
                \includegraphics[width=1\linewidth]{june5th/Screenshot 2025-06-10 at 11.33.19 PM.png}
                \caption{AMVF Pub Note: \href{https://cds.cern.ch/record/2670380/files/ATL-PHYS-PUB-2019-015.pdf}{Link}}
            \end{figure}
        \end{column}

        \begin{column}{0.5\textwidth}
            \begin{figure}
                \centering
                \includegraphics[width=1\linewidth]{june25th/stacked_pv_classification_5_80_95thr_recovertex.png}
                \caption{Thr = 0.95}
            \end{figure}
        \end{column}
    \end{columns}
\end{frame}

% \begin{frame}{Comparison of Reco Vertex Classification: MaxScore}
%     \begin{columns}
%         \begin{column}{0.5\textwidth}
%             \begin{figure}
%                 \centering
%                 \includegraphics[width=1\linewidth]{example-image-a}
%                 \caption{MaxScore}
%             \end{figure}
%         \end{column}

%         \begin{column}{0.5\textwidth}
%             \begin{figure}
%                 \centering
%                 \includegraphics[width=1\linewidth]{june5th/stacked_pv_classification_5_80_thr95_recovertex.png}
%                 \caption{Thr = 0.95}
%             \end{figure}
%         \end{column}
%     \end{columns}
% \end{frame}

\begin{frame}{A Look into GNN Associations and Thresholds}
    \begin{itemize}
        \item Classification rates skew less clean with AMVF-style definition but are similar
        \item How do the track-vertex associations look?
        \item What is causing the increase in merged vertices?
        \item Are there tracks that are not being used due to threshold?
        \item Lets look at test cases!
    \end{itemize}
    \begin{figure}
        \centering
        \includegraphics[width=0.4\linewidth]{this_is_fine.jpg}
    \end{figure}
\end{frame}

\begin{frame}{Case Study: Looking at Merged/Split Reco PVs}
    \begin{columns}
        \begin{column}{0.5\textwidth}
            \begin{itemize}
                \item Total Reco PVs: 42
                \item Total Truth PVs: 59
                \item Thr=0.85
                \begin{itemize}
                    \item Num Clean: 22
                    \item Num Merged: 12
                    \item Num Split: 8
                    \item Num Fake: 0
                \end{itemize}
                \item Thr=0.95
                \begin{itemize}
                    \item Num Clean: 30
                    \item Num Merged: 7
                    \item Num Split: 5
                    \item Num Fake: 0
                \end{itemize}
            \end{itemize}
        \end{column}

        \begin{column}{0.5\textwidth}
            \begin{figure}
                \centering
                \includegraphics[width=1\linewidth]{june25th/casestudy_example_reco_truth_hist_dist_31_38_51.png}
            \end{figure}
        \end{column}
    \end{columns}
\end{frame}

\begin{frame}{Case Study: Three Nearby Vertices Thr=0.85 (1)}
    \begin{columns}
        \begin{column}{0.4\textwidth}
            \begin{itemize}
                \item Reco PV 24: 33 Assoc tracks
                \begin{itemize}
                    \item Old Class: Merged
                    \item \(W_{\text{reco},24} = 33\)
                    \item \(\sum p_T^2 = 4e7\) (GeV\(^2\))
                \end{itemize}
                \item Truth PV 31: 13 Assoc tracks
                \begin{itemize}
                    \item \(W_{\text{truth},31} = 12\)
                    \item \(W_{\text{truth},31}/W_{\text{reco},24} = 0.4\)
                \end{itemize}
                \item Truth PV 38: 10 Assoc tracks
                \begin{itemize}
                    \item \(W_{\text{truth},38} = 10\)
                    \item \(W_{\text{truth},38}/W_{\text{reco},24} = 0.3\)
                \end{itemize}
                \item Truth PV 51: 6 Assoc tracks
                \begin{itemize}
                    \item \(W_{\text{truth},51} = 5\)
                    \item \(W_{\text{truth},51}/W_{\text{reco},24} = 0.15\)
                \end{itemize}
                \item AMVF-style Class: \textbf{Split}
            \end{itemize}
        \end{column}

        \begin{column}{0.6\textwidth}
            \begin{figure}
                \centering
                \includegraphics[width=1\linewidth]{june25th/casestudy_example_reco_truth_hist_dist_24.png}
            \end{figure}
        \end{column}
    \end{columns}
\end{frame}

\begin{frame}{Case Study: Three Nearby Vertices Thr=0.85 (2)}
    \begin{columns}
        \begin{column}{0.4\textwidth}
            \begin{itemize}
                \item Reco PV 25: 42 Assoc tracks
                \begin{itemize}
                    \item Old Class: Clean
                    \item \(\sum p_T^2 = 6e7\) (GeV\(^2\))
                    \item \(W_{\text{reco},25} = 42\)
                \end{itemize}
                \item Truth PV 31: 13 Assoc tracks
                \begin{itemize}
                    \item \(W_{\text{truth},31} = 12\)
                    \item \(W_{\text{truth},31}/W_{\text{reco},25} = 0.4\)
                \end{itemize}
                \item Truth PV 38: 10 Assoc tracks
                \begin{itemize}
                    \item \(W_{\text{truth},38} = 10\)
                    \item \(W_{\text{truth},38}/W_{\text{reco},25} = 0.3\)
                \end{itemize}
                \item Truth PV 51: 6 Assoc tracks
                \begin{itemize}
                    \item \(W_{\text{truth},51} = 5\)
                    \item \(W_{\text{truth},51}/W_{\text{reco},25} = 0.15\)
                \end{itemize}
                \item AMVF-style Class: \textbf{Merged}
            \end{itemize}
        \end{column}

        \begin{column}{0.6\textwidth}
            \begin{figure}
                \centering
                \includegraphics[width=1\linewidth]{june25th/casestudy_example_reco_truth_hist_dist_25.png}
            \end{figure}
        \end{column}
    \end{columns}
\end{frame}

\begin{frame}{Case Study: Three Nearby Vertices Thr=0.95 (1)}
    \begin{columns}
        \begin{column}{0.4\textwidth}
            \begin{itemize}
                \item Reco PV 24: 24 Assoc tracks
                \begin{itemize}
                    \item \(W_{\text{reco},24} = 24\)
                    \item \(\sum p_T^2 = 4e7\) (GeV\(^2\))
                \end{itemize}
                \item Truth PV 31: 13 Assoc tracks
                \begin{itemize}
                    \item \(W_{\text{truth},31} = 12\)
                    \item \(W_{\text{truth},31}/W_{\text{reco},24} = 0.5\)
                \end{itemize}
                \item Truth PV 38: 10 Assoc tracks
                \begin{itemize}
                    \item \(W_{\text{truth},38} = 6\)
                    \item \(W_{\text{truth},38}/W_{\text{reco},24} = 0.25\)
                \end{itemize}
                \item Truth PV 51: 6 Assoc tracks
                \begin{itemize}
                    \item \(W_{\text{truth},51} = 5\)
                    \item \(W_{\text{truth},51}/W_{\text{reco},24} = 0.21\)
                \end{itemize}
                \item AMVF-style Class: \textbf{Merged}
            \end{itemize}
        \end{column}

        \begin{column}{0.6\textwidth}
            \begin{figure}
                \centering
                \includegraphics[width=1\linewidth]{june25th/casestudy_example_reco_truth_hist_dist_24.png}
            \end{figure}
        \end{column}
    \end{columns}
\end{frame}

\begin{frame}{Case Study: Three Nearby Vertices Thr=0.95 (2)}
    \begin{columns}
        \begin{column}{0.4\textwidth}
            \begin{itemize}
                \item Reco PV 25: 22 Assoc tracks
                \begin{itemize}
                    \item \(W_{\text{reco},25} = 22\)
                    \item \(\sum p_T^2 = 3e7\) (GeV\(^2\))
                \end{itemize}
                \item Truth PV 31: 13 Assoc tracks
                \begin{itemize}
                    \item \(W_{\text{truth},31} = 9\)
                    \item \(W_{\text{truth},31}/W_{\text{reco},25} = 0.4\)
                \end{itemize}
                \item Truth PV 38: 10 Assoc tracks
                \begin{itemize}
                    \item \(W_{\text{truth},38} = 9\)
                    \item \(W_{\text{truth},38}/W_{\text{reco},25} = 0.3\)
                \end{itemize}
                \item Truth PV 51: 6 Assoc tracks
                \begin{itemize}
                    \item \(W_{\text{truth},51} = 0\)
                \end{itemize}
                \item AMVF-style Class: \textbf{Merged}
            \end{itemize}
        \end{column}

        \begin{column}{0.6\textwidth}
            \begin{figure}
                \centering
                \includegraphics[width=1\linewidth]{june25th/casestudy_example_reco_truth_hist_dist_25.png}
            \end{figure}
        \end{column}
    \end{columns}
\end{frame}

\begin{frame}{Number of Associations Per Tracks Varying Threshold}
    \begin{columns}
        \begin{column}{0.5\textwidth}
            \begin{figure}
                \centering
                \includegraphics[width=1\linewidth]{june25th/redo_casestudy_num_track_assocs_thr85.png}
            \end{figure}
        \end{column}

        \begin{column}{0.5\textwidth}
            \begin{figure}
                \centering
                \includegraphics[width=1\linewidth]{june25th/redo_casestudy_num_track_assocs_thr95.png}
            \end{figure}
        \end{column}
    \end{columns}
\end{frame}

\begin{frame}{Re: A Look into GNN Associations and Thresholds}
    \begin{itemize}
        \item What is causing the increase in merged/split vertices?
        \begin{itemize}
            \item Seems to be that distinguishing tracks for nearby vertices are tough
        \end{itemize}
        \item Are there tracks that are not being used due to threshold?
        \begin{itemize}
            \item Tracks are being used even in high thresholds
        \end{itemize}
        \item Are there other ways on evaluating performance?
        \begin{itemize}
            \item \textbf{Hard Scatter Vertex Classification}
        \end{itemize}
    \end{itemize}
    \begin{figure}
        \centering
        \includegraphics[width=0.4\linewidth]{this_is_fine.jpg}
    \end{figure}
\end{frame}

\begin{frame}{Concluding Remarks}
    \begin{itemize}
        \item Presented a ML tool to find and fit primary vertices
        \item How to continue improving performance...? What is the next direction...?
        \begin{itemize}
            \item \textbf{Modifying Graphs}
            \begin{itemize}
                \item Edge Attributes
                \item Binary Cross Entropy to Cross Entropy
            \end{itemize}
            \item Start applying cuts to the output of GNN
            \item Limit 1 edge per track
            \item \textbf{Using Reco PVs to train GNN}
        \end{itemize}
        \item Developing tools to continue analyzing PV-Finder + GNN outputs
        \begin{itemize}
            \item \textbf{HS Vertex Quality Metrics}
        \end{itemize}
        \item Abstract accepted for ACAT 2025!
    \end{itemize}
\end{frame}

\begin{frame}[plain]
    \centering
    July 2nd LHCb and Rocky/Lauren Meeting
\end{frame}

\begin{frame}{Intro Musings/Ramblings}
    \centering
    \textit{Starting To Converge!}
    \begin{itemize}
        \item Studies:
        \begin{itemize}
            \item (Delayed) GNN Training on Reco Vertices
            \item \textbf{Edge Attributes}
            \item \textbf{Hard-Scatter Vertex Metrics}
            \begin{itemize}
                \item HS Vertex Quality
                \item HS Vertex Reco Efficiency
                \item HS Vertex Selection Efficiency
            \end{itemize}
        \end{itemize}
    \end{itemize}
\end{frame}

\begin{frame}[plain]
    \centering
    Edge Attributes Studies
\end{frame}

\begin{frame}{Re: Forming Graphs}
    \begin{figure}
        \centering
        \includegraphics[width=0.6\linewidth]{april4th/forming_edges_gnn.png}
    \end{figure}
\end{frame}

\begin{frame}{Adding Edge Attributes: Longitudinal Significance}
    \begin{figure}
        \centering
        \includegraphics[width=0.6\linewidth]{july2nd/forming_edge_attr_1.png}
    \end{figure}
\end{frame}

\begin{frame}{Adding Edge Attributes: Horizontal Significance}
    \begin{figure}
        \centering
        \includegraphics[width=0.6\linewidth]{july2nd/forming_edge_attr_2.png}
    \end{figure}
\end{frame}

\begin{frame}{Calculating Width of Reco PVs from Histograms}
    \begin{itemize}
        \item Now need to calculate variance of each PV histogram peak
        \begin{align*}
            \sigma^2_z &= \mathbb{E} \left(z^2 \right) - \mathbb{E} \left(z \right)^2 \\
            \mathbb{E} \left( z \right) &= \frac{\sum_i w_i z_i}{\sum w_i} \\
            \mathbb{E} \left( z^2 \right) &= \frac{\sum_i w_i z_i^2}{\sum w_i} \\
            w_i &= \text{histogram value at bin i} \\
            z_i &= \text{bin position i}
        \end{align*}
        \item Note: \(\left[\sigma_z^2\right] = \left[\text{mm}^2\right]\)
        \item To Do: Try this on PV-Finder output instead of \(\sigma_{\text{vtx-vtx}}\)?
    \end{itemize}
\end{frame}

\begin{frame}{Sanity Check on \(\sigma_z\) Calculation}
    \begin{figure}
        \centering
        \includegraphics[width=0.6\linewidth]{july2nd/sanity_check_sigma_z_recovtx.png}
    \end{figure}
\end{frame}

\begin{frame}{New Steps Outlined}
    \begin{itemize}
        \item Train PV-Finder
        \begin{itemize}
            \item Using current best PV-Finder run (UNet with Dropout Model)
        \end{itemize}
        \item \textbf{(New!)} Calculate predicted primary vertex location and \(\sigma_z\) from distribution
        \item Construct graphs from reco PV information
        \item Evaluate reco PV using GNN model trained on truth PV information
        \begin{itemize}
            \item \textbf{(New!)} Using new GNN layer type (\href{pytorch-geometric.readthedocs.io/en/2.5.3/generated/torch_geometric.nn.conv.GATConv.html}{GATConv}) that can incorporate edge attributes
            \item Will make slides later on explaining the layer!
        \end{itemize}
        \item Calculate Metrics!
    \end{itemize}
\end{frame}

\begin{frame}{Training with Edge Attr Classification at \(\left< \mu \right> =60\) (Thr=0.95)}
    \begin{columns}
        \begin{column}{0.55\textwidth}
            \begin{itemize}
                \item Total Reconstructed PVs: 61245
                \item Total Truth PVs: 72189
                \item Total Clean: 48871, Rate: 79.80\%
                \item Total Merged: 8627, Rate: 14.09\%
                \item Total Split: 3393, Rate: 5.54\%
                \item Total Fake: 354, Rate: 0.58\%
            \end{itemize}
        \end{column}

        \begin{column}{0.45\textwidth}
            \begin{figure}
                \centering
                \includegraphics[width=1\linewidth]{july2nd/reco_pv_avg_bar_chart_thr95_recovertex.png}
            \end{figure}
        \end{column}
    \end{columns}
\end{frame}

\begin{frame}{Training with Edge Attr Classification at \(\left< \mu \right> =60\) (MaxScore)}
    \begin{columns}
        \begin{column}{0.55\textwidth}
            \begin{itemize}
                \item Total Reconstructed PVs: 61245
                \item Total Truth PVs: 72189
                \item Total Clean: 49147, Rate: 80.25\%
                \item Total Merged: 9972, Rate: 16.28\%
                \item Total Split: 2073, Rate: 3.38\%
                \item Total Fake: 53, Rate: 0.09\%
            \end{itemize}
        \end{column}

        \begin{column}{0.45\textwidth}
            \begin{figure}
                \centering
                \includegraphics[width=1\linewidth]{july2nd/reco_pv_avg_bar_chart_maxscore_recovertex.png}
            \end{figure}
        \end{column}
    \end{columns}
\end{frame}

\begin{frame}{Comparison of Quality Classification Old vs New}
    \begin{columns}
        \begin{column}{0.5\textwidth}
            \begin{figure}
                \centering
                \includegraphics[width=1\linewidth]{june5th/reco_pv_avg_bar_chart_thr95_recovertex.png}
                \caption{(Old) Thr=0.95}
            \end{figure}
        \end{column}

        \begin{column}{0.5\textwidth}
            \begin{figure}
                \centering
                \includegraphics[width=1\linewidth]{july2nd/reco_pv_avg_bar_chart_maxscore_recovertex.png}
                \caption{\textbf{(New!)} MaxScore}
            \end{figure}
        \end{column}
    \end{columns}
\end{frame}

\begin{frame}{Comparison of Reco Vertex Classification: Thresholds (1)}
    \begin{columns}
        \begin{column}{0.5\textwidth}
            \begin{figure}
                \centering
                \includegraphics[width=1\linewidth]{june25th/stacked_pv_classification_5_80_95thr_recovertex.png}
                \caption{(Old) Thr = 0.95}
            \end{figure}
        \end{column}

        \begin{column}{0.5\textwidth}
            \begin{figure}
                \centering
                \includegraphics[width=1\linewidth]{july2nd/stacked_pv_classification_39_80_95thr_recovertex_GATCONV.png}
                \caption{\textbf{(New!)} Thr = 0.95}
            \end{figure}
        \end{column}
    \end{columns}
\end{frame}

\begin{frame}{Comparison of Reco Vertex Classification: Thresholds (2)}
    \begin{columns}
        \begin{column}{0.5\textwidth}
            \begin{figure}
                \centering
                \includegraphics[width=1\linewidth]{june5th/Screenshot 2025-06-10 at 11.33.19 PM.png}
                \caption{AMVF Pub Note (Fig. 7): \href{https://cds.cern.ch/record/2670380/files/ATL-PHYS-PUB-2019-015.pdf}{Link}}
            \end{figure}
        \end{column}

        \begin{column}{0.5\textwidth}
            \begin{figure}
                \centering
                \includegraphics[width=1\linewidth]{july2nd/stacked_pv_classification_39_80_95thr_recovertex_GATCONV.png}
                \caption{\textbf{(New!)} Thr = 0.95}
            \end{figure}
        \end{column}
    \end{columns}
\end{frame}

\begin{frame}{Comparison of Reco Vertex Classification: Thresholds (3)}
    \begin{columns}
        \begin{column}{0.5\textwidth}
            \begin{figure}
                \centering
                \includegraphics[width=1\linewidth]{july2nd/stacked_pv_classification_39_80_maxscore_recovertex_GATCONV.png}
                \caption{\textbf{(New!)} MaxScore}
            \end{figure}
        \end{column}

        \begin{column}{0.5\textwidth}
            \begin{figure}
                \centering
                \includegraphics[width=1\linewidth]{july2nd/stacked_pv_classification_39_80_95thr_recovertex_GATCONV.png}
                \caption{\textbf{(New!)} Thr = 0.95}
            \end{figure}
        \end{column}
    \end{columns}
\end{frame}

\begin{frame}{Number of Associations Per Tracks Varying Threshold}
    \begin{columns}
        \begin{column}{0.5\textwidth}
            \begin{figure}
                \centering
                \includegraphics[width=1\linewidth]{july2nd/gat_casestudy_num_track_assocs_maxscore.png}
                \caption{(Sanity Check!) MaxScore}
            \end{figure}
        \end{column}

        \begin{column}{0.5\textwidth}
            \begin{figure}
                \centering
                \includegraphics[width=1\linewidth]{july2nd/gat_casestudy_num_track_assocs_thr95.png}
                \caption{Average < 1}
            \end{figure}
        \end{column}
    \end{columns}
\end{frame}

\begin{frame}[plain]
    \centering
    Hard Scatter Vertex Classification
\end{frame}

\begin{frame}{HS Vertex Quality}
    \begin{itemize}
        \item HS Vertex Quality Classification:
        \begin{itemize}
            \item \textbf{Clean/Matched}: \textbf{CLEAN} reco vertex matching true HS interaction AND \textbf{track weights don't contribute \(> 50 \%\) to other vertices}
            \item \textbf{LowPU}: \textbf{MERGED} reco vertex with \textbf{at least 50\%} of the accumulated \textbf{track weight} coming \textbf{from the simulated HS interaction}
            \item \textbf{HighPU}: \textbf{MERGED} reco vertex the HS interaction contributing between \textbf{1\% and 50\% of the accumulated track weight}"
            \item \textbf{PurePU}: No reco vertex with at least 1\% accumulated track weight from HS interaction"
        \end{itemize}
        \item Edge Case:
        \begin{itemize}
            \item What to do when, there is a split reco vertex?
            \item A reco vertex that correctly matches 60 tracks from HS vertex (Clean)
            \item Another reco vertex that correctly matches 59 tracks from HS vertex (Split)
            \item Due to this, have 67 events "unclassified" in 2550 total events
        \end{itemize}
    \end{itemize}
\end{frame}

\begin{frame}{HS Vertex Quality Category Plot (1)}
    \begin{columns}
        \begin{column}{0.5\textwidth}
            \begin{figure}
                \centering
                \includegraphics[width=1\linewidth]{june25th/fig_07_amvf.png}
                \caption{AMVF Pub Note: \href{https://cds.cern.ch/record/2670380/files/ATL-PHYS-PUB-2019-015.pdf}{Link}}
            \end{figure}
        \end{column}

        \begin{column}{0.5\textwidth}
            \begin{figure}
                \centering
                \includegraphics[width=1\linewidth]{july2nd/rhs_vertex_quality_category_GATCONV_maxscore.png}
                \caption{\textbf{(New!)} MaxScore}
            \end{figure}
        \end{column}
    \end{columns}
\end{frame}

\begin{frame}{HS Vertex Quality Category Plot (2)}
    \begin{columns}
        \begin{column}{0.5\textwidth}
            \begin{figure}
                \centering
                \includegraphics[width=1\linewidth]{july2nd/rhs_vertex_quality_category_GATCONV_thr.png}
                \caption{\textbf{(New!)} Thr = 0.95}
            \end{figure}
        \end{column}

        \begin{column}{0.5\textwidth}
            \begin{figure}
                \centering
                \includegraphics[width=1\linewidth]{july2nd/rhs_vertex_quality_category_GATCONV_maxscore.png}
                \caption{\textbf{(New!)} MaxScore}
            \end{figure}
        \end{column}
    \end{columns}
\end{frame}

\begin{frame}{HS Vertex Reconstruction Efficiency}
    \begin{columns}
        \begin{column}{0.3\textwidth}
            \begin{itemize}
                \item Fraction of events which the HS vertex is reconstructed and classified as Clean/Matched, LowPU, or HighPu
                \item Local pile-up density is the number of vertices within a \(\pm\) 2mm window of a HS vertex
            \end{itemize}
        \end{column}

        \begin{column}{0.7\textwidth}
            \begin{figure}
                \centering
                \includegraphics[width=0.85\linewidth]{july2nd/vertex_reconstruction_efficiency_comparison.png}
            \end{figure}
        \end{column}
    \end{columns}
\end{frame}

\begin{frame}{Comparison: HS Vertex Reconstruction Efficiency}
    \begin{columns}
        \begin{column}{0.5\textwidth}
            \begin{figure}
                \centering
                \includegraphics[width=1\linewidth]{june25th/fig_08a_amvf_reco.png}
                \caption{AMVF Pub Note (Fig. 8): \href{https://cds.cern.ch/record/2670380/files/ATL-PHYS-PUB-2019-015.pdf}{Link}}
            \end{figure}
        \end{column}

        \begin{column}{0.5\textwidth}
            \begin{figure}
                \centering
                \includegraphics[width=1\linewidth]{july2nd/vertex_reconstruction_efficiency_comparison.png}
                \caption{Thr = 0.95}
            \end{figure}
        \end{column}
    \end{columns}
\end{frame}

\begin{frame}{HS Vertex Selection Efficiency}
    \begin{columns}
        \begin{column}{0.3\textwidth}
            \begin{itemize}
                \item Fraction of events where the reconstructed vertex with highest \(\sum p_T^2\) is the one contianing the largest total weight from true HS tracks
                \item \textbf{Currently the only metric that needs improvements!}
            \end{itemize}
        \end{column}

        \begin{column}{0.7\textwidth}
            \begin{figure}
                \centering
                \includegraphics[width=0.85\linewidth]{july2nd/vertex_selection_efficiency_comparison.png}
            \end{figure}
        \end{column}
    \end{columns}
\end{frame}

\begin{frame}{Comparison: HS Vertex Selection Efficiency}
    \begin{columns}
        \begin{column}{0.5\textwidth}
            \begin{figure}
                \centering
                \includegraphics[width=1\linewidth]{june25th/fig_09a_amvf_selection.png}
                \caption{AMVF Pub Note (Fig. 9): \href{https://cds.cern.ch/record/2670380/files/ATL-PHYS-PUB-2019-015.pdf}{Link}}
            \end{figure}
        \end{column}

        \begin{column}{0.5\textwidth}
            \begin{figure}
                \centering
                \includegraphics[width=1\linewidth]{july2nd/vertex_selection_efficiency_comparison.png}
            \end{figure}
        \end{column}
    \end{columns}
\end{frame}

\begin{frame}{Concluding Remarks}
    \begin{itemize}
        \item Demonstrated improved performance using edge attributes and maxscore!
        \item Starting to converge now! Still want last bit of performance for HS selection!
        \item How to continue improving performance...? What is the next direction...?
        \begin{itemize}
            \item \textbf{Input Level?}
            \begin{itemize}
                \item Beam Spot Information/POCA Variables?
                \item Reco PV x and y?
            \end{itemize}
            \item \textbf{Output Level?}
            \begin{itemize}
                \item Modified MaxScore?
                \item Cuts on the output of GNN?
            \end{itemize}
            \item Continue trying to train GNN using Reco PVs
        \end{itemize}
        \item Developing tools to continue analyzing PV-Finder + GNN outputs?
        \item \textbf{Abstract accepted for ACAT 2025!}
    \end{itemize}
\end{frame}

\begin{frame}[plain]
    \centering
    July 9th Rocky Meeting
\end{frame}

\begin{frame}{Intro Musings/Ramblings}
    \centering
    \begin{itemize}
        \item Studies:
        \begin{itemize}
            \item Why Edge Attributes?
            \item \textbf{Edge Attributes Case Study}
            \item \textbf{Modified Max Score}
            \item GNN Training on Reco Vertices
            \item \textbf{(Delayed) GATv2Conv Layer}
        \end{itemize}
        \item Studies in Backup:
        \begin{itemize}
            \item Applying Cuts \textbf{(Setting thresholds essentually does same thing)}
            \item Training with Truth PV x and y \textbf{(Converges to about same point)}
            % \item TTVA Efficiency and Impurity
        \end{itemize}
    \end{itemize}
\end{frame}

\begin{frame}[plain]
    \centering
    Edge Attributes
\end{frame}

\begin{frame}{Why Edge Attributes?}
    \begin{itemize}
        \item Are the edge attributes the reason why the GNN performs better?
        \item Can GNN reach same results by just adding \(\sigma_z\) to PV information?
        \begin{itemize}
            \item GraphConv is simpler, faster, and uses less weights than GATConv
        \end{itemize}
        \item Did a test run with old graph layer but with added pv width information
        \begin{itemize}
            \item \textbf{Answer: No}... need to research why GATConv learns better
        \end{itemize}
    \end{itemize}
\end{frame}

\begin{frame}[plain]
    \centering
    Edge Attributes Case Study
\end{frame}

\begin{frame}{Edge Attributes Case Study: Three Nearby Vertices (1)}
    \begin{columns}
        \begin{column}{0.4\textwidth}
            \begin{itemize}
                \item Reco PV 24: 20 Assoc tracks
                \begin{itemize}
                    \item Old Class: Merged
                    \item \(W_{\text{reco},24} = 20\)
                    \item \(\sum p_T^2 = 3e7\) (MeV\(^2\))
                \end{itemize}
                \item Truth PV 31: 13 Assoc tracks
                \begin{itemize}
                    \item \(W_{\text{truth},11} = 11\)
                    \item \(W_{\text{truth},31}/W_{\text{reco},24} = \)
                \end{itemize}
                \item Truth PV 38: 10 Assoc tracks
                \begin{itemize}
                    \item \(W_{\text{truth},38} = 4\)
                    \item \(W_{\text{truth},38}/W_{\text{reco},24} = \)
                \end{itemize}
                \item Truth PV 51: 6 Assoc tracks
                \begin{itemize}
                    \item \(W_{\text{truth},51} = 5\)
                    \item \(W_{\text{truth},51}/W_{\text{reco},24} = \)
                \end{itemize}
                \item AMVF-style Class: \textbf{Merged}
            \end{itemize}
        \end{column}

        \begin{column}{0.6\textwidth}
            \begin{figure}
                \centering
                \includegraphics[width=1\linewidth]{june25th/casestudy_example_reco_truth_hist_dist_24.png}
            \end{figure}
        \end{column}
    \end{columns}
\end{frame}

\begin{frame}{Edge Attributes Case Study: Three Nearby Vertices (2)}
    \begin{columns}
        \begin{column}{0.4\textwidth}
            \begin{itemize}
                \item Reco PV 24: 6 Assoc tracks
                \begin{itemize}
                    \item Old Class: Clean
                    \item \(W_{\text{reco},24} = 6\)
                    \item \(\sum p_T^2 = 9e6\) (MeV\(^2\))
                \end{itemize}
                \item Truth PV 38: 10 Assoc tracks
                \begin{itemize}
                    \item \(W_{\text{truth},38} = 6\)
                    \item \(W_{\text{truth},38}/W_{\text{reco},25} = 1\)
                \end{itemize}
                \item AMVF-style Class: \textbf{Clean}
            \end{itemize}
        \end{column}

        \begin{column}{0.6\textwidth}
            \begin{figure}
                \centering
                \includegraphics[width=1\linewidth]{june25th/casestudy_example_reco_truth_hist_dist_25.png}
            \end{figure}
        \end{column}
    \end{columns}
\end{frame}

\begin{frame}{Edge Attributes Case Study: Split HS Vertex}
    \begin{columns}
        \begin{column}{0.3\textwidth}
            \begin{itemize}
                \item More of an issue with PV-Finder's output thanso GNN
                \item Causes GNN HS reconstruction to be undefined
            \end{itemize}
        \end{column}

        \begin{column}{0.7\textwidth}
            \begin{figure}
                \centering
                \includegraphics[width=0.9\linewidth]{july9th/casestudy_splitHS.png}
            \end{figure}
        \end{column}
    \end{columns}
\end{frame}

\begin{frame}{Edge Attributes Case Study: HS Selection at High PU Density}
    \begin{columns}
        \begin{column}{0.4\textwidth}
            \begin{itemize}
                \item Reco PV 14: 63 Assoc tracks
                \begin{itemize}
                    \item Old Class: Clean
                    \item \(W_{\text{reco},14} = 63\)
                    \item \(\sum p_T^2 = 55e7\) (MeV\(^2\))
                \end{itemize}
                \item Truth PV 2: 44 Assoc Tracks
                \begin{itemize}
                    \item \(W_{\text{truth},2} = 2\)
                    \item \(W_{\text{truth},2}/W_{\text{reco},14} = 0.03\)
                \end{itemize}
                \item Truth PV 14: 65 Assoc Tracks
                \begin{itemize}
                    \item \(W_{\text{truth},14} = 60\)
                    \item \(W_{\text{truth},2}/W_{\text{reco},14} = 0.95\)
                \end{itemize}
                \item AMVF-style Class: \textbf{Clean}
                \item Not selected as HS because another vertex has higher \(\sum p_T^2\)
            \end{itemize}
        \end{column}

        \begin{column}{0.6\textwidth}
            \begin{figure}
                \centering
                \includegraphics[width=1\linewidth]{july9th/casestudy_clean_notselected_vertex.png}
            \end{figure}
        \end{column}
    \end{columns}
\end{frame}

\begin{frame}[plain]
    \centering
    Modified Max Score
\end{frame}

\begin{frame}{Modified Max Score}
    \begin{itemize}
        \item Want to prevent fake tracks from being associated!
        \item (Old) Require every track to have one association based on highest score
        \item \textbf{(New)} Add an "if" statement where if all the scores are below a threshold, a track can potentially have zero associations
        \item Look at effects it has on HS vertex efficiency and selection
        \item Spoilers: Generally just slight improvements overall
    \end{itemize}
\end{frame}

\begin{frame}{Max Score vs Modified Max Score at Thr=0.5}
    \begin{columns}
        \begin{column}{0.5\linewidth}
            \begin{itemize}
                \item Max Score
                \begin{itemize}
                    \item Total Reconstructed PVs: 61245
                    \item Total Truth PVs: 72189
                    \item Total Clean: 49147, Rate: 80.25\%
                    \item Total Merged: 9972, Rate: 16.28\%
                    \item Total Split: 2073, Rate: 3.38\%
                    \item Total Fake: 53, Rate: 0.09\%
                \end{itemize}
            \end{itemize}
        \end{column}

        \begin{column}{0.5\linewidth}
            \begin{itemize}
                \item Modified Max Score (Thr=0.5)
                \begin{itemize}
                    \item Total Reconstructed PVs: 61245
                    \item Total Truth PVs: 72189
                    \item Total Clean: 49755, Rate: 81.24\%
                    \item Total Merged: 9357, Rate: 15.28\%
                    \item Total Split: 2080, Rate: 3.40\%
                    \item Total Fake: 53, Rate: 0.09\%
                \end{itemize}
            \end{itemize}
        \end{column}
    \end{columns}
\end{frame}

\begin{frame}{Modified Max Score at Different Thresholds}
    \begin{columns}
        \begin{column}{0.5\linewidth}
            \begin{itemize}
                \item Modified Max Score (Thr=0.5)
                \begin{itemize}
                    \item \textbf{Total Clean: 49755, Rate: 81.24\%}
                    \item Total Merged: 9357, Rate: 15.28\%
                    \item Total Split: 2080, Rate: 3.40\%
                    \item Total Fake: 53, Rate: 0.09\%
                \end{itemize}
                \item Modified Max Score (Thr=0.9)
                \begin{itemize}
                    \item \textbf{Total Clean: 51260, Rate: 83.70\%}
                    \item Total Merged: 7875, Rate: 12.86\%
                    \item Total Split: 2037, Rate: 3.33\%
                    \item Total Fake: 73, Rate: 0.12\%
                \end{itemize}
            \end{itemize}
        \end{column}

        \begin{column}{0.5\linewidth}
            \begin{itemize}
                \item Modified Max Score (Thr=0.7)
                \begin{itemize}
                    \item \textbf{Total Clean: 50001, Rate: 81.64\%}
                    \item Total Merged: 9106, Rate: 14.87\%
                    \item Total Split: 2084, Rate: 3.40\%
                    \item Total Fake: 54, Rate: 0.09\%
                \end{itemize}
                \item Modified Max Score (Thr=0.95)
                \begin{itemize}
                    \item \textbf{Total Clean: 52626, Rate: 85.93\%}
                    \item Total Merged: 6252, Rate: 10.21\%
                    \item Total Split: 1921, Rate: 3.14\%
                    \item Total Fake: 446, Rate: 0.73\%
                \end{itemize}
            \end{itemize}
        \end{column}
    \end{columns}
\end{frame}

\begin{frame}{HS Vertex Reconstruction Efficiency Comparisons}
    \begin{columns}
        \begin{column}{0.5\textwidth}
            \begin{figure}
                \centering
                \includegraphics[width=1\linewidth]{june25th/fig_08a_amvf_reco.png}
                \caption{AMVF Pub Note (Fig. 8): \href{https://cds.cern.ch/record/2670380/files/ATL-PHYS-PUB-2019-015.pdf}{Link}}
            \end{figure}
        \end{column}

        \begin{column}{0.5\textwidth}
            \begin{figure}
                \centering
                \includegraphics[width=1\linewidth]{july9th/vertex_reconstruction_efficiency_comparison_modmaxscore.png}
                \caption{Thr = 0.95}
            \end{figure}
        \end{column}
    \end{columns}
\end{frame}

\begin{frame}{HS Vertex Reconstruction Efficiency Comparisons}
    \begin{columns}
        \begin{column}{0.5\textwidth}
            \begin{figure}
                \centering
                \includegraphics[width=1\linewidth]{july2nd/vertex_reconstruction_efficiency_comparison.png}
            \end{figure}
        \end{column}

        \begin{column}{0.5\textwidth}
            \begin{figure}
                \centering
                \includegraphics[width=1\linewidth]{july9th/vertex_reconstruction_efficiency_comparison_modmaxscore.png}
            \end{figure}
        \end{column}
    \end{columns}
\end{frame}

\begin{frame}{HS Vertex Selection Efficiency Comparisons}
    \begin{columns}
        \begin{column}{0.5\textwidth}
            \begin{figure}
                \centering
                \includegraphics[width=1\linewidth]{june25th/fig_09a_amvf_selection.png}
                \caption{AMVF Pub Note (Fig. 9a): \href{https://cds.cern.ch/record/2670380/files/ATL-PHYS-PUB-2019-015.pdf}{Link}}
            \end{figure}
        \end{column}

        \begin{column}{0.5\textwidth}
            \begin{figure}
                \centering
                \includegraphics[width=1\linewidth]{july9th/vertex_selection_efficiency_comparison_modmaxscore.png}
                \caption{Thr = 0.95}
            \end{figure}
        \end{column}
    \end{columns}
\end{frame}

\begin{frame}{HS Vertex Selection Efficiency Comparisons}
    \begin{columns}
        \begin{column}{0.5\textwidth}
            \begin{figure}
                \centering
                \includegraphics[width=1\linewidth]{july2nd/vertex_selection_efficiency_comparison.png}
            \end{figure}
        \end{column}

        \begin{column}{0.5\textwidth}
            \begin{figure}
                \centering
                \includegraphics[width=1\linewidth]{july9th/vertex_selection_efficiency_comparison_modmaxscore.png}
            \end{figure}
        \end{column}
    \end{columns}
\end{frame}

\begin{frame}[plain]
    \centering
    Training GNN on Reco Vertices
\end{frame}

\begin{frame}{Recap: Steps Outlined}
    \begin{itemize}
        \item Train PV-Finder
        \begin{itemize}
            \item Using current best PV-Finder run (UNet with Dropout Model)
        \end{itemize}
        \item Calculate predicted primary vertex location
        \item Construct graphs from reconstructed PV information
        \item \textbf{Current:} Evaluate using GNN model trained on truth information
        \begin{itemize}
            \item Using current best GNN run (GraphConv Model)
            \item Use a modified max score create associations
        \end{itemize}
        \item \textbf{Idea:} Train a GNN model using reco vertices matched to truth vertices
        \begin{itemize}
            \item Allows GNN to learn directly from PV-Finder outputs and associations
            \item Use a modified max score create associations
        \end{itemize}
    \end{itemize}
\end{frame}

\begin{frame}{Constructing Trainable Data from Reco Vertices}
    \begin{itemize}
        \item Based on Original PV-Finder Classification:
        \begin{itemize}
            \item Clean: Reco tracks (from matched truth vertex) \(\rightarrow\) reco vertex
            \item Merged: Reco tracks (from both matched truth vertex) \(\rightarrow\) reco vertex
            \item Fake: No reco tracks \(\rightarrow\) reco vertex
            \item Split: Remove reco vertices from training
            \begin{itemize}
                \item Don't want GNN to learn features of doubly associating reco tracks
                \item Found a pesky bug where I didn't deal with vertex case where it is both a split and merged
            \end{itemize}
        \end{itemize}
        \item Then you would do vertex classification on test data comparing to truth vertex
        \item Questions on constructing graphs?
        \begin{itemize}
            \item Should fake reco vertices be included in training?
        \end{itemize}
    \end{itemize}
\end{frame}

\begin{frame}{HS Vertex Reconstruction Efficiency Comparisons}
    \begin{columns}
        \begin{column}{0.5\textwidth}
            \begin{figure}
                \centering
                \includegraphics[width=1\linewidth]{july9th/vertex_reconstruction_efficiency_comparison_modmaxscore.png}
            \end{figure}
        \end{column}

        \begin{column}{0.5\textwidth}
            \begin{figure}
                \centering
                \includegraphics[width=1\linewidth]{july9th/vertex_reconstruction_efficiency_comparison_recovtx_modmaxscore.pdf}
            \end{figure}
        \end{column}
    \end{columns}
\end{frame}

\begin{frame}{HS Vertex Selection Efficiency Comparisons}
    \begin{columns}
        \begin{column}{0.5\textwidth}
            \begin{figure}
                \centering
                \includegraphics[width=1\linewidth]{july9th/vertex_selection_efficiency_comparison_modmaxscore.png}
            \end{figure}
        \end{column}

        \begin{column}{0.5\textwidth}
            \begin{figure}
                \centering
                \includegraphics[width=1\linewidth]{july9th/vertex_selection_efficiency_comparison_recovtx_maxscore.pdf}
            \end{figure}
        \end{column}
    \end{columns}
\end{frame}

\begin{frame}{Concluding Remarks (1)}
    \begin{itemize}
        \item How to continue improving performance...? What is the next direction...?
        \begin{itemize}
            \item \textbf{Input Level?}
            \begin{itemize}
                \item Beam Spot Information/POCA Variables? \textbf{(Don't think it'll add much...)}
                \item Truth PV x and y? \textbf{(Converges to about same point)}
                \item Note: Adding \(p_T\) or longitudinal significance doesn't affect output
            \end{itemize}
            \item \textbf{Output Level?}
            \begin{itemize}
                \item Modified MaxScore? \textbf{(Better than Thr and Max Score)}
                \item Cuts on the output of GNN? \textbf{(Setting thresholds essentially does same thing)}
            \end{itemize}
            \item Training GNN using Reco PVs
            \begin{itemize}
                \item \textbf{Results in similar/slightly worse performance than training on truth information}
                \item Might need to rethink how to construct graphs?
            \end{itemize}
            \item TTVA Efficiency and Impurity for HS Vertex
        \end{itemize}
        \item Alternative Ideas:
        \begin{itemize}
            \item Weighted \(\sum p_T^2\) for Vertex Selection \href{https://arxiv.org/abs/1910.07681}{(Link)}
            \item Talk with Prof. Ariel Schwartzman
        \end{itemize}
    \end{itemize}
\end{frame}

\begin{frame}{Concluding Remarks (2)}
    \begin{itemize}
        \item Seems to be starting to hit bottom of barrel... should I start shifting focus on?
        \begin{itemize}
            \item Pushing everything onto GitLab
            \item Start writing rest of ACAT 2025 paper
        \end{itemize}
        \item TVPD Presentation this Friday?
        \item Am I attending/presenting ACAT 2025?
    \end{itemize}
\end{frame}

\begin{frame}[plain]
    \centering
    Applying Cuts
\end{frame}

\begin{frame}{Current: Steps Outlined}
    \begin{itemize}
        \item Train PV-Finder
        \begin{itemize}
            \item Using current best PV-Finder run (UNet with Dropout Model)
        \end{itemize}
        \item Calculate predicted primary vertex location
        \item Construct graphs from reconstructed PV information
        \item Evaluate using GNN model trained on truth information
        \begin{itemize}
            \item Using current best GNN run (GraphConv Model)
            \item Use a modified max score create associations
        \end{itemize}
        \item \textbf{Idea:} Apply cuts to track associations
    \end{itemize}
\end{frame}

\begin{frame}{Run 2 Cuts}
    \begin{itemize}
        \item Cut on impact variable parameters:
        \begin{itemize}
            \item Loose:
            \begin{itemize}
                \item Track was used in a vertex fit
                \item If not used in ANY vertex fit and \(|\Delta z_0 \sin{\theta}| < 3 \text{mm}\)
            \end{itemize}
            \item Nominal:
            \begin{itemize}
                \item \(| d_0 | < 2 \text{mm}\)
                \item \(|\Delta z_0 \sin{\theta}| < 3 \text{mm}\)
            \end{itemize}
            \item Tight:
            \begin{itemize}
                \item \(| d_0 | < 0.5 \text{mm}\)
                \item \(|\Delta z_0 \sin{\theta}| < 0.5 \text{mm}\)
            \end{itemize}
        \end{itemize}
        \item \href{https://twiki.cern.ch/twiki/bin/viewauth/AtlasProtected/TrackingCPRecsRun2Final}{Link}
    \end{itemize}
\end{frame}

\begin{frame}{HS Vertex Selection Efficiency: Nominal}
    \begin{columns}
        \begin{column}{0.5\textwidth}
            \begin{figure}
                \centering
                \includegraphics[width=1\linewidth]{july9th/vertex_selection_efficiency_comparison_modmaxscore.png}
                \caption{ModMaxScore without Cuts}
            \end{figure}
        \end{column}

        \begin{column}{0.5\textwidth}
            \begin{figure}
                \centering
                \includegraphics[width=1\linewidth]{july9th/vertex_selection_efficiency_comparison_maxscore_cuts.pdf}
                \caption{ModMaxScore with Cuts}
            \end{figure}
        \end{column}
    \end{columns}
\end{frame}

\begin{frame}{HS Vertex Reconstruction Efficiency: Nominal vs Tight}
    \begin{columns}
        \begin{column}{0.5\textwidth}
            \begin{figure}
                \centering
                \includegraphics[width=1\linewidth]{july9th/vertex_selection_efficiency_comparison_maxscore_cuts.pdf}
                \caption{ModMaxScore with Tight Cuts}
            \end{figure}
        \end{column}

        \begin{column}{0.5\textwidth}
            \begin{figure}
                \centering
                \includegraphics[width=1\linewidth]{july9th/vertex_selection_efficiency_comparison_maxscore_cuts_tight.pdf}
                \caption{ModMaxScore with Nominal Cuts}
            \end{figure}
        \end{column}
    \end{columns}
\end{frame}



\begin{frame}[plain]
    \centering
    Tompkins' Group Meeting July 15th
\end{frame}

\begin{frame}[plain]
    \centering
    Group Meeting Slide
\end{frame}

\begin{frame}{Group Meeting Round Table}
    \begin{itemize}
        \item Learned: \(\sum p_T^2\) isn't always the best method to select a hard-scatter vertex
        \item Accomplished: Updated GNN architecture to also learn hard-scatter vertex selection!
        \item To Do: Improve selection efficiency of hard-scatter vertex
    \end{itemize}
    \begin{figure}
        \centering
        \includegraphics[width=0.6\linewidth]{july16th/pipeline_hspv.png}
    \end{figure}
\end{frame}

\begin{frame}[plain]
    \centering
    Tompkins' and Rocky Weekly Meeting July 16th
\end{frame}

\begin{frame}{Intro Musings/Ramblings}
    \centering
    \begin{itemize}
        \item Studies for this Presentation:
        \begin{itemize}
            \item \textbf{Multi-Task Learning}
            \item \textbf{Multi-Task Learning Overfitting}
        \end{itemize}
        \item Studies in Backup:
        \begin{itemize}
            \item Applying Cuts \textbf{(Setting thresholds essentually does same thing)}
            \item Training with Truth PV x and y \textbf{(Need to test again)}
            \item Why Edge Attributes?
            \item GNN Training on Reco Vertices \textbf{(Need to test again)}
            \item \textbf{GATv2Conv Layer}
            \begin{itemize}
                \item \textbf{Using more "expensive" layers causes model to overfit on truth graph}
                \item \textbf{Does not generalize well evaluating on reco vertices}
            \end{itemize}
            \item \textbf{Edge Attributes Case Study}
            \item \textbf{Modified Max Score}
            % \item TTVA Efficiency and Impurity
        \end{itemize}
    \end{itemize}
\end{frame}

\begin{frame}[plain]
    \centering
    Multi-Task Learning
\end{frame}

\begin{frame}{Idea for Multi-Task Learning}
    \begin{itemize}
        \item \textbf{(Current)} Select HS vertex to be vertex with highest \(\sum p_T^2\)
        \item \textbf{(Idea)} Have GNN also output a classification of what PV is the HS vertex
        \begin{itemize}
            \item \(\sum p_T^2\) isn’t always the optimal method for selecting HS Vertex
            \item Interesting work has been done using GN2X for HS vertex selection: \href{indico.cern.ch/event/1374927/contributions/5942568/attachments/2882952/5052279/20240624_hsgnn.pdf}{Link}
            \item Our method would allow us to use \textbf{lower level representations} for classification
        \end{itemize}
        \item Notes:
        \begin{itemize}
            \item Update to Loss Function
        \end{itemize}
        \begin{align*}
            \Lagr = \Lagr_{\text{BCE, Edge}} + \Lagr_{\text{BCE, HSPV}}
        \end{align*}
    \end{itemize}
\end{frame}

\begin{frame}{Method 1: GNN Pipeline Updated After GCN}
    \begin{figure}
        \centering
        \includegraphics[width=0.8\linewidth]{july16th/pipeline_hspv.png}
    \end{figure}
\end{frame}

\begin{frame}{Method 1: MLP HS PV Predictor Layer}
    \begin{figure}
        \centering
        \includegraphics[width=0.8\linewidth]{july16th/pipeline_hspv_output.png}
    \end{figure}
\end{frame}

\begin{frame}{HS Vertex Selection Efficiency Comparisons}
    \begin{columns}
        \begin{column}{0.5\textwidth}
            \begin{figure}
                \centering
                \includegraphics[width=1\linewidth]{june25th/fig_09a_amvf_selection.png}
                \caption{AMVF Pub Note (Fig. 9a): \href{https://cds.cern.ch/record/2670380/files/ATL-PHYS-PUB-2019-015.pdf}{Link}}
            \end{figure}
        \end{column}

        \begin{column}{0.5\textwidth}
            \begin{figure}
                \centering
                \includegraphics[width=1\linewidth]{july16th/vertex_selection_efficiency_comparison_gnn_sumpT.pdf}
                \caption{New Results with Multi-Task Learning}
            \end{figure}
        \end{column}
    \end{columns}
\end{frame}

\begin{frame}{HS Vertex Selection Efficiency Comparisons}
    \begin{columns}
        \begin{column}{0.5\textwidth}
            \begin{figure}
                \centering
                \includegraphics[width=1\linewidth]{july2nd/vertex_selection_efficiency_comparison.png}
            \end{figure}
        \end{column}

        \begin{column}{0.5\textwidth}
            \begin{figure}
                \centering
                \includegraphics[width=1\linewidth]{july16th/vertex_selection_efficiency_comparison_gnn_sumpT.pdf}
            \end{figure}
        \end{column}
    \end{columns}
\end{frame}

\begin{frame}{HS Vertex Selection Efficiency Comparisons}
    \begin{columns}
        \begin{column}{0.5\textwidth}
            \begin{figure}
                \centering
                \includegraphics[width=1\linewidth]{july9th/vertex_selection_efficiency_comparison_modmaxscore.png}
            \end{figure}
        \end{column}

        \begin{column}{0.5\textwidth}
            \begin{figure}
                \centering
                \includegraphics[width=1\linewidth]{july16th/vertex_selection_efficiency_comparison_gnn_sumpT.pdf}
            \end{figure}
        \end{column}
    \end{columns}
\end{frame}

\begin{frame}{Method 2: Graph Pooling Layer}
    \begin{figure}
        \centering
        \includegraphics[width=0.6\linewidth]{july16th/potential_inbuilt_HS_class.png}
    \end{figure}
\end{frame}

\begin{frame}{Overfitting Musings}
    \begin{itemize}
        \item Combinations Tried:
        \begin{itemize}
            \item GATConv + Method 1 (results shown before)
            \item GATConv + Method 1 + Track \(p_T\) (overfits)
            \item GATConv + Method 2 (overfits)
            \item GATv2Conv + Method 2 (overfits)
        \end{itemize}
        \item Will try training on reco vertices again but with HSPV selection
        \begin{itemize}
            \item If event has split reco vertices on HS vertex, throw out event
            \item \textbf{Question:} Are we beginning to be limited by PV-Finder?
        \end{itemize}
        \item Will try training on just HSPV selection
    \end{itemize}
\end{frame}

\begin{frame}{Concluding Remarks (1)}
    \begin{itemize}
        \item How to continue improving performance...? What is the next direction...?
        \begin{itemize}
            \item \textbf{Input Level?}
            \begin{itemize}
                \item Beam Spot Information/POCA Variables? \textbf{(Don't think it'll add much...)}
                \item Truth PV x and y?
            \end{itemize}
            \item \textbf{Output Level?}
            \begin{itemize}
                \item Modified MaxScore? \textbf{(Better than Thr and Max Score)}
                \item Cuts on the output of GNN? \textbf{(Setting thresholds essentially does same thing)}
            \end{itemize}
            \item Training GNN using Reco PVs
            \begin{itemize}
                \item \textbf{Results in similar/slightly worse performance than training on truth information}
                \item \textbf{Need to retry with multi-task learning}
            \end{itemize}
            \item \textbf{Multi-Task Learning}
            \begin{itemize}
                \item Continue testing (seems to be a possible solution...)
                \item Provide more information into MLP HSPV Prediction Layer?
            \end{itemize}
        \end{itemize}
    \end{itemize}
\end{frame}

\begin{frame}{Concluding Remarks (2)}
    \begin{itemize}
        \item After we get a lovely selection efficiency improvement...
        \begin{itemize}
            \item TTVA Efficiency and Impurity Metrics
            \item Clean up code
            \item Push everything into GitLab
        \end{itemize}
        \item \textbf{...and we are done!}
        \item Paper
        \begin{itemize}
            \item What portions of GNN should I include? (Probably only up to edge attributes?)
            \item Go through edits and start drafting GNN portion
            \item Update current GNN plots to be ATLAS Style
        \end{itemize}
    \end{itemize}
\end{frame}

\begin{frame}[plain]
    \centering
    TVPD July 18th
\end{frame}

\begin{frame}{Recap: Connecting the Pipelines}
    \begin{figure}
        \centering
        \includegraphics[width=0.55\linewidth]{april4th/pipeline_pvfinder_gnn.png}
    \end{figure}
\end{frame}

\begin{frame}{Introduction}
    \begin{itemize}
        \item Prior PV-Finder + GNN results have demonstrated an end-to-end ML method for primary vertex reconstruction
        \item Lingering concerns over track associations and reco vertex merged/split rate
        \item This Presentation:
        \begin{itemize}
            \item Edge Attributes
            \item Preliminary Performance
            \item Case Study
            \item Hard-Scater Vertex Reconstruction and Selection
        \end{itemize}
    \end{itemize}
\end{frame}

\begin{frame}[plain]
    \centering
    Edge Attributes Studies
\end{frame}

\begin{frame}{Re: Forming Graphs}
    \begin{figure}
        \centering
        \includegraphics[width=0.6\linewidth]{april4th/forming_edges_gnn.png}
    \end{figure}
\end{frame}

\begin{frame}{Adding Edge Attributes: Longitudinal Significance}
    \begin{figure}
        \centering
        \includegraphics[width=0.6\linewidth]{july2nd/forming_edge_attr_1.png}
    \end{figure}
\end{frame}

% \begin{frame}{Adding Edge Attributes: Horizontal Significance}
%     \begin{figure}
%         \centering
%         \includegraphics[width=0.6\linewidth]{july2nd/forming_edge_attr_2.png}
%     \end{figure}
% \end{frame}

\begin{frame}{Calculating Width of Reco PVs from Histograms}
    \begin{itemize}
        \item Now need to calculate variance of each PV histogram peak
        \begin{align*}
            \sigma^2_z &= \mathbb{E} \left(z^2 \right) - \mathbb{E} \left(z \right)^2 \\
            \mathbb{E} \left( z \right) &= \frac{\sum_i w_i z_i}{\sum w_i} \\
            \mathbb{E} \left( z^2 \right) &= \frac{\sum_i w_i z_i^2}{\sum w_i} \\
            w_i &= \text{histogram value at bin i} \\
            z_i &= \text{bin position i}
        \end{align*}
        \item Note: \(\left[\sigma_z^2\right] = \left[\text{mm}^2\right]\)
        \item To Do: Try this on PV-Finder output instead of \(\sigma_{\text{vtx-vtx}}\)?
    \end{itemize}
\end{frame}

\begin{frame}{Sanity Check on \(\sigma_z\) Calculation}
    \begin{figure}
        \centering
        \includegraphics[width=0.6\linewidth]{july2nd/sanity_check_sigma_z_recovtx.png}
    \end{figure}
\end{frame}

\begin{frame}[plain]
    \centering
    Preliminary Performance
\end{frame}

\begin{frame}{New Steps Outlined}
    \begin{itemize}
        \item Train PV-Finder on truth level information
        \begin{itemize}
            \item Using current best PV-Finder run (UNet with Dropout Model)
        \end{itemize}
        \item \textbf{(New!)} Calculate predicted primary vertex location and \(\sigma_z\) from distribution
        \item Train GNN on truth level information
        \item Evaluate reco PV using GNN model trained on truth PV information
        \begin{itemize}
            \item \textbf{(New!)} Using new GNN layer type (\href{pytorch-geometric.readthedocs.io/en/2.5.3/generated/torch_geometric.nn.conv.GATConv.html}{GATConv}) that can incorporate edge attributes
            \item Will make slides later on explaining the layer!
        \end{itemize}
        \item Creating associations between tracks and PVs
        \begin{itemize}
            \item Threshold Method: All scores above a threshold associates a track and PV
            \item MaxScore Method: Each track has 1 association based on highest scoring PV
        \end{itemize}
        \item Calculate Metrics!
    \end{itemize}
\end{frame}

\begin{frame}{Training with Edge Attr Classification at \(\left< \mu \right> =60\) (Thr=0.95)}
    \begin{columns}
        \begin{column}{0.55\textwidth}
            \begin{itemize}
                \item Total Reconstructed PVs: 61245
                \item Total Truth PVs: 72189
                \item Total Clean: 48871, Rate: 79.80\%
                \item Total Merged: 8627, Rate: 14.09\%
                \item Total Split: 3393, Rate: 5.54\%
                \item Total Fake: 354, Rate: 0.58\%
            \end{itemize}
        \end{column}

        \begin{column}{0.45\textwidth}
            \begin{figure}
                \centering
                \includegraphics[width=1\linewidth]{july2nd/reco_pv_avg_bar_chart_thr95_recovertex.png}
            \end{figure}
        \end{column}
    \end{columns}
\end{frame}

\begin{frame}{Training with Edge Attr Classification at \(\left< \mu \right> =60\) (MaxScore)}
    \begin{columns}
        \begin{column}{0.55\textwidth}
            \begin{itemize}
                \item Total Reconstructed PVs: 61245
                \item Total Truth PVs: 72189
                \item Total Clean: 49147, Rate: 80.25\%
                \item Total Merged: 9972, Rate: 16.28\%
                \item Total Split: 2073, Rate: 3.38\%
                \item Total Fake: 53, Rate: 0.09\%
            \end{itemize}
        \end{column}

        \begin{column}{0.45\textwidth}
            \begin{figure}
                \centering
                \includegraphics[width=1\linewidth]{july2nd/reco_pv_avg_bar_chart_maxscore_recovertex.png}
            \end{figure}
        \end{column}
    \end{columns}
\end{frame}

\begin{frame}{Comparison of Quality Classification Old vs New}
    \begin{columns}
        \begin{column}{0.5\textwidth}
            \begin{figure}
                \centering
                \includegraphics[width=1\linewidth]{june5th/reco_pv_avg_bar_chart_thr95_recovertex.png}
                \caption{(No Edge Attr) Thr=0.95}
            \end{figure}
        \end{column}

        \begin{column}{0.5\textwidth}
            \begin{figure}
                \centering
                \includegraphics[width=1\linewidth]{july2nd/reco_pv_avg_bar_chart_maxscore_recovertex.png}
                \caption{\textbf{(New!)} MaxScore}
            \end{figure}
        \end{column}
    \end{columns}
\end{frame}

\begin{frame}{Comparison of Reco Vertex Classification: Thresholds (1)}
    \begin{columns}
        \begin{column}{0.5\textwidth}
            \begin{figure}
                \centering
                \includegraphics[width=1\linewidth]{june25th/stacked_pv_classification_5_80_95thr_recovertex.png}
                \caption{(No Edge Attr) Thr = 0.95}
            \end{figure}
        \end{column}

        \begin{column}{0.5\textwidth}
            \begin{figure}
                \centering
                \includegraphics[width=1\linewidth]{july2nd/stacked_pv_classification_39_80_95thr_recovertex_GATCONV.png}
                \caption{\textbf{(New!)} Thr = 0.95}
            \end{figure}
        \end{column}
    \end{columns}
\end{frame}

\begin{frame}{Comparison of Reco Vertex Classification: Thresholds (2)}
    \begin{columns}
        \begin{column}{0.5\textwidth}
            \begin{figure}
                \centering
                \includegraphics[width=1\linewidth]{june5th/Screenshot 2025-06-10 at 11.33.19 PM.png}
                \caption{AMVF Pub Note (Fig. 7): \href{https://cds.cern.ch/record/2670380/files/ATL-PHYS-PUB-2019-015.pdf}{Link}}
            \end{figure}
        \end{column}

        \begin{column}{0.5\textwidth}
            \begin{figure}
                \centering
                \includegraphics[width=1\linewidth]{july2nd/stacked_pv_classification_39_80_95thr_recovertex_GATCONV.png}
                \caption{\textbf{(New!)} Thr = 0.95}
            \end{figure}
        \end{column}
    \end{columns}
\end{frame}

\begin{frame}{Comparison of Reco Vertex Classification: Thresholds (3)}
    \begin{columns}
        \begin{column}{0.5\textwidth}
            \begin{figure}
                \centering
                \includegraphics[width=1\linewidth]{july2nd/stacked_pv_classification_39_80_maxscore_recovertex_GATCONV.png}
                \caption{\textbf{(New!)} MaxScore}
            \end{figure}
        \end{column}

        \begin{column}{0.5\textwidth}
            \begin{figure}
                \centering
                \includegraphics[width=1\linewidth]{july2nd/stacked_pv_classification_39_80_95thr_recovertex_GATCONV.png}
                \caption{\textbf{(New!)} Thr = 0.95}
            \end{figure}
        \end{column}
    \end{columns}
\end{frame}

\begin{frame}{Number of Associations Per Tracks Varying Threshold}
    \begin{columns}
        \begin{column}{0.5\textwidth}
            \begin{figure}
                \centering
                \includegraphics[width=1\linewidth]{july2nd/gat_casestudy_num_track_assocs_maxscore.png}
                \caption{(Sanity Check!) MaxScore}
            \end{figure}
        \end{column}

        \begin{column}{0.5\textwidth}
            \begin{figure}
                \centering
                \includegraphics[width=1\linewidth]{july2nd/gat_casestudy_num_track_assocs_thr95.png}
                \caption{Average < 1}
            \end{figure}
        \end{column}
    \end{columns}
\end{frame}

\begin{frame}[plain]
    \centering
    Edge Attributes Case Study
\end{frame}

\begin{frame}{Edge Attributes Case Study: Three Nearby Vertices (1)}
    \begin{columns}
        \begin{column}{0.4\textwidth}
            \begin{itemize}
                \item Reco PV 24: 20 Assoc tracks
                \begin{itemize}
                    \item Old Class: Merged
                    \item \(W_{\text{reco},24} = 20\)
                    \item \(\sum p_T^2 = 3e7\) (MeV\(^2\))
                \end{itemize}
                \item Truth PV 31: 13 Assoc tracks
                \begin{itemize}
                    \item \(W_{\text{truth},11} = 11\)
                    \item \(W_{\text{truth},31}/W_{\text{reco},24} = 0.55\)
                \end{itemize}
                \item Truth PV 38: 10 Assoc tracks
                \begin{itemize}
                    \item \(W_{\text{truth},38} = 4\)
                    \item \(W_{\text{truth},38}/W_{\text{reco},24} = 0.2\)
                \end{itemize}
                \item Truth PV 51: 6 Assoc tracks
                \begin{itemize}
                    \item \(W_{\text{truth},51} = 5\)
                    \item \(W_{\text{truth},51}/W_{\text{reco},24} = 0.25\)
                \end{itemize}
                \item AMVF-style Class: \textbf{Merged}
            \end{itemize}
        \end{column}

        \begin{column}{0.6\textwidth}
            \begin{figure}
                \centering
                \includegraphics[width=1\linewidth]{june25th/casestudy_example_reco_truth_hist_dist_24.png}
            \end{figure}
        \end{column}
    \end{columns}
\end{frame}

\begin{frame}{Edge Attributes Case Study: Three Nearby Vertices (2)}
    \begin{columns}
        \begin{column}{0.4\textwidth}
            \begin{itemize}
                \item Reco PV 24: 6 Assoc tracks
                \begin{itemize}
                    \item Old Class: Clean
                    \item \(W_{\text{reco},24} = 6\)
                    \item \(\sum p_T^2 = 9e6\) (MeV\(^2\))
                \end{itemize}
                \item Truth PV 38: 10 Assoc tracks
                \begin{itemize}
                    \item \(W_{\text{truth},38} = 6\)
                    \item \(W_{\text{truth},38}/W_{\text{reco},25} = 1\)
                \end{itemize}
                \item AMVF-style Class: \textbf{Clean}
            \end{itemize}
        \end{column}

        \begin{column}{0.6\textwidth}
            \begin{figure}
                \centering
                \includegraphics[width=1\linewidth]{june25th/casestudy_example_reco_truth_hist_dist_25.png}
            \end{figure}
        \end{column}
    \end{columns}
\end{frame}

\begin{frame}{Edge Attributes Case Study: Split HS Vertex}
    \begin{columns}
        \begin{column}{0.3\textwidth}
            \begin{itemize}
                \item More of an issue with PV-Finder's output thanso GNN
                \item Causes GNN HS reconstruction to be undefined
            \end{itemize}
        \end{column}

        \begin{column}{0.7\textwidth}
            \begin{figure}
                \centering
                \includegraphics[width=0.9\linewidth]{july9th/casestudy_splitHS.png}
            \end{figure}
        \end{column}
    \end{columns}
\end{frame}

\begin{frame}{Edge Attributes Case Study: HS Selection at High PU Density}
    \begin{columns}
        \begin{column}{0.4\textwidth}
            \begin{itemize}
                \item Reco PV 14: 63 Assoc tracks
                \begin{itemize}
                    \item Old Class: Clean
                    \item \(W_{\text{reco},14} = 63\)
                    \item \(\sum p_T^2 = 55e7\) (MeV\(^2\))
                \end{itemize}
                \item Truth PV 2: 44 Assoc Tracks
                \begin{itemize}
                    \item \(W_{\text{truth},2} = 2\)
                    \item \(W_{\text{truth},2}/W_{\text{reco},14} = 0.03\)
                \end{itemize}
                \item Truth PV 14: 65 Assoc Tracks
                \begin{itemize}
                    \item \(W_{\text{truth},14} = 60\)
                    \item \(W_{\text{truth},2}/W_{\text{reco},14} = 0.95\)
                \end{itemize}
                \item AMVF-style Class: \textbf{Clean}
                \item Not selected as HS because another vertex has higher \(\sum p_T^2\)
            \end{itemize}
        \end{column}

        \begin{column}{0.6\textwidth}
            \begin{figure}
                \centering
                \includegraphics[width=1\linewidth]{july9th/casestudy_clean_notselected_vertex.png}
            \end{figure}
        \end{column}
    \end{columns}
\end{frame}

\begin{frame}[plain]
    \centering
    Hard Scatter Vertex Classification
\end{frame}

\begin{frame}{HS Vertex Quality}
    \begin{itemize}
        \item HS Vertex Quality Classification:
        \begin{itemize}
            \item \textbf{Clean/Matched}: \textbf{CLEAN} reco vertex matching true HS interaction AND \textbf{track weights don't contribute \(> 50 \%\) to other vertices}
            \item \textbf{LowPU}: \textbf{MERGED} reco vertex with \textbf{at least 50\%} of the accumulated \textbf{track weight} coming \textbf{from the simulated HS interaction}
            \item \textbf{HighPU}: \textbf{MERGED} reco vertex the HS interaction contributing between \textbf{1\% and 50\% of the accumulated track weight}"
            \item \textbf{PurePU}: No reco vertex with at least 1\% accumulated track weight from HS interaction"
        \end{itemize}
    \end{itemize}
\end{frame}

\begin{frame}{HS Vertex Quality Category Plot (1)}
    \begin{columns}
        \begin{column}{0.5\textwidth}
            \begin{figure}
                \centering
                \includegraphics[width=1\linewidth]{june25th/fig_07_amvf.png}
                \caption{AMVF Pub Note: \href{https://cds.cern.ch/record/2670380/files/ATL-PHYS-PUB-2019-015.pdf}{Link}}
            \end{figure}
        \end{column}

        \begin{column}{0.5\textwidth}
            \begin{figure}
                \centering
                \includegraphics[width=1\linewidth]{july2nd/rhs_vertex_quality_category_GATCONV_maxscore.png}
                \caption{\textbf{(New!)} MaxScore}
            \end{figure}
        \end{column}
    \end{columns}
\end{frame}

\begin{frame}{HS Vertex Quality Category Plot (2)}
    \begin{columns}
        \begin{column}{0.5\textwidth}
            \begin{figure}
                \centering
                \includegraphics[width=1\linewidth]{july2nd/rhs_vertex_quality_category_GATCONV_thr.png}
                \caption{\textbf{(New!)} Thr = 0.95}
            \end{figure}
        \end{column}

        \begin{column}{0.5\textwidth}
            \begin{figure}
                \centering
                \includegraphics[width=1\linewidth]{july2nd/rhs_vertex_quality_category_GATCONV_maxscore.png}
                \caption{\textbf{(New!)} MaxScore}
            \end{figure}
        \end{column}
    \end{columns}
\end{frame}

\begin{frame}{HS Vertex Reconstruction Efficiency}
    \begin{columns}
        \begin{column}{0.3\textwidth}
            \begin{itemize}
                \item Fraction of events which the HS vertex is reconstructed and classified as Clean/Matched, LowPU, or HighPu
                \item Local pile-up density is the number of vertices within a \(\pm\) 2mm window of a HS vertex
            \end{itemize}
        \end{column}

        \begin{column}{0.7\textwidth}
            \begin{figure}
                \centering
                \includegraphics[width=0.85\linewidth]{july2nd/vertex_reconstruction_efficiency_comparison.png}
            \end{figure}
        \end{column}
    \end{columns}
\end{frame}

\begin{frame}{Comparison: HS Vertex Reconstruction Efficiency}
    \begin{columns}
        \begin{column}{0.5\textwidth}
            \begin{figure}
                \centering
                \includegraphics[width=1\linewidth]{june25th/fig_08a_amvf_reco.png}
                \caption{AMVF Pub Note (Fig. 8): \href{https://cds.cern.ch/record/2670380/files/ATL-PHYS-PUB-2019-015.pdf}{Link}}
            \end{figure}
        \end{column}

        \begin{column}{0.5\textwidth}
            \begin{figure}
                \centering
                \includegraphics[width=1\linewidth]{july2nd/vertex_reconstruction_efficiency_comparison.png}
                \caption{Thr = 0.95}
            \end{figure}
        \end{column}
    \end{columns}
\end{frame}

\begin{frame}{HS Vertex Selection Efficiency}
    \begin{columns}
        \begin{column}{0.3\textwidth}
            \begin{itemize}
                \item Fraction of events where the reconstructed vertex with highest \(\sum p_T^2\) is the one contianing the largest total weight from true HS tracks
                \item \textbf{Currently the only metric that needs improvements!}
            \end{itemize}
        \end{column}

        \begin{column}{0.7\textwidth}
            \begin{figure}
                \centering
                \includegraphics[width=0.85\linewidth]{july2nd/vertex_selection_efficiency_comparison.png}
            \end{figure}
        \end{column}
    \end{columns}
\end{frame}

\begin{frame}{Comparison: HS Vertex Selection Efficiency}
    \begin{columns}
        \begin{column}{0.5\textwidth}
            \begin{figure}
                \centering
                \includegraphics[width=1\linewidth]{june25th/fig_09a_amvf_selection.png}
                \caption{AMVF Pub Note (Fig. 9): \href{https://cds.cern.ch/record/2670380/files/ATL-PHYS-PUB-2019-015.pdf}{Link}}
            \end{figure}
        \end{column}

        \begin{column}{0.5\textwidth}
            \begin{figure}
                \centering
                \includegraphics[width=1\linewidth]{july2nd/vertex_selection_efficiency_comparison.png}
            \end{figure}
        \end{column}
    \end{columns}
\end{frame}

\begin{frame}{Concluding Remarks}
    \begin{itemize}
        \item Demonstrated improved performance using edge attributes and maxscore!
        \item Starting to converge now! Still want last bit of performance for HS selection!
        \item How to continue improving performance...? What is the next direction...?
        \begin{itemize}
            \item Inbuilt GNN HS Selection?
            \item Training GNN using Reco PVs?
        \end{itemize}
        \item \textbf{Conference paper draft in progress for ACAT 2025!}
    \end{itemize}
\end{frame}

\begin{frame}[plain]
    \centering
    Group Meeting Updates 07/23
\end{frame}

\begin{frame}{Schedule Musings (1)}
    \begin{itemize}
        \item Work Scheduling:
        \begin{itemize}
            \item Group Presentation July 29th
            \item TVPD August ???
            \item \textbf{Week of August 4th to 8th (Doctoral Summer Bridge Program)}
            \item \textbf{Week of August 25th to 29th (Vacation)}
            \begin{itemize}
                \item Take a break before graduate school fiesta starts
            \end{itemize}
            \item \textbf{Week of September 8th to 12th (ACAT 2025)}
            \item September 15th, Last Day of Post-Bacc Employment
            \item September 16th/23rd, Race with Rory (?)
            \item \textbf{September/October ???, IRIS-HEP Retreat}
            \item \textbf{November 10th to 14th, Connecting the Dots 2025}
            \begin{itemize}
                \item Accepted as a poster + lightning talk
            \end{itemize}
        \end{itemize}
    \end{itemize}
\end{frame}

\begin{frame}{Schedule Musings (2)}
    \begin{itemize}
        \item \textbf{Goals for Remaining Time:}
        \begin{itemize}
            \item (July/August) Push everything to Git repository
            \item (July/August) Finalize draft of paper and plots
            \item (August/September) Finalize remaining studies
            \item (August/September) Comparison with AMVF TTVA
        \end{itemize}
        \item Future?:
        \begin{itemize}
            \item \textbf{Have a demanding Fall workload of Mech + E\&M + QFT + 2 Courses to TA + ???}
            \item What work do we want to include for poster in Connecting the Dots?
            \item Am I doing a presentation for IRIS-HEP retreat? (Can it be virtual?)
        \end{itemize}
    \end{itemize}
\end{frame}

\begin{frame}{Schedule Musings (3)}
    \begin{itemize}
        \item Studies:
        \begin{itemize}
            \item kNN Graph Construction and Evaluation (same results but more stable)
        \end{itemize}
        \item Comments:
        \begin{itemize}
            \item \textbf{Have been preparing everything to be pushed into GitLab}
            \item \textbf{Will do final case studies in coming week...}
            \item Tried augmenting inputs, outputs, loss functions, models,...
            \item Hitting limit of performance and have to either modify original PV-Finder or dataset
            \item Lower level inputs? (detector silicon hits)
            \item PV x and y? (can give small improvements)
        \end{itemize}
    \end{itemize}
\end{frame}

\begin{frame}{Git Musings}
    \begin{itemize}
        \item \textbf{Main branch} will include work up until edge attributes
        \item \textbf{Experimental branch} will include work involving training reco PVs, different HSPV/pooling graph models, kNN graph construction, and much more
        \item \textit{To Do:}
        \begin{itemize}
            \item Push all new files
            \item Update README.md
            \item Give repository a full chain run through
        \end{itemize}
    \end{itemize}
\end{frame}

\begin{frame}[plain]
    \centering
    Tompkins' Group Meeting 07/29
\end{frame}

\begin{frame}{Recap: Connecting the Pipelines}
    \begin{figure}
        \centering
        \includegraphics[width=0.55\linewidth]{april4th/pipeline_pvfinder_gnn.png}
    \end{figure}
\end{frame}

\begin{frame}{(Draft) Recap: MLP + UNet Architecture}
    \begin{figure}
        \centering
        \includegraphics[width=0.9\linewidth]{july29th/draft_MLPandUNet(1).png}
    \end{figure}
\end{frame}

\begin{frame}{(Draft) Recap: GNN Architecture}
    \begin{figure}
        \centering
        \includegraphics[width=1\linewidth]{july29th/draft_graph_architecture.png}
    \end{figure}
\end{frame}

\begin{frame}{Introduction}
    \begin{itemize}
        \item Prior PV-Finder + GNN results have demonstrated an end-to-end ML method for primary vertex reconstruction
        \item Lingering concerns over track associations and reco vertex merged/split rate
        \item This Presentation:
        \begin{itemize}
            \item Edge Attributes
            \item Preliminary Performance
            \item Case Study
            \item Hard-Scater Vertex Reconstruction and Selection
        \end{itemize}
    \end{itemize}
\end{frame}

\begin{frame}[plain]
    \centering
    Edge Attributes Studies
\end{frame}

\begin{frame}{Re: Forming Graphs}
    \begin{figure}
        \centering
        \includegraphics[width=0.6\linewidth]{april4th/forming_edges_gnn.png}
    \end{figure}
\end{frame}

\begin{frame}{Adding Edge Attributes: Longitudinal Significance}
    \begin{figure}
        \centering
        \includegraphics[width=0.6\linewidth]{july2nd/forming_edge_attr_1.png}
    \end{figure}
\end{frame}

% \begin{frame}{Adding Edge Attributes: Horizontal Significance}
%     \begin{figure}
%         \centering
%         \includegraphics[width=0.6\linewidth]{july2nd/forming_edge_attr_2.png}
%     \end{figure}
% \end{frame}

\begin{frame}{Calculating Width of Reco PVs from Histograms}
    \begin{itemize}
        \item Now need to calculate variance of each PV histogram peak
        \begin{align*}
            \sigma^2_z &= \mathbb{E} \left(z^2 \right) - \mathbb{E} \left(z \right)^2 \\
            \mathbb{E} \left( z \right) &= \frac{\sum_i w_i z_i}{\sum w_i} \\
            \mathbb{E} \left( z^2 \right) &= \frac{\sum_i w_i z_i^2}{\sum w_i} \\
            w_i &= \text{histogram value at bin i} \\
            z_i &= \text{bin position i}
        \end{align*}
        \item Note: \(\left[\sigma_z^2\right] = \left[\text{mm}^2\right]\)
        \item To Do: Try this on PV-Finder output instead of \(\sigma_{\text{vtx-vtx}}\)?
    \end{itemize}
\end{frame}

\begin{frame}{Sanity Check on \(\sigma_z\) Calculation}
    \begin{figure}
        \centering
        \includegraphics[width=0.6\linewidth]{july2nd/sanity_check_sigma_z_recovtx.png}
    \end{figure}
\end{frame}

\begin{frame}[plain]
    \centering
    Preliminary Performance
\end{frame}

\begin{frame}{New Steps Outlined}
    \begin{itemize}
        \item Train PV-Finder on truth level information
        \begin{itemize}
            \item Using current best PV-Finder run (UNet with Dropout Model)
        \end{itemize}
        \item \textbf{(New!)} Calculate predicted primary vertex location and \(\sigma_z\) from distribution
        \item Train GNN on truth level information
        \item Evaluate reco PV using GNN model trained on truth PV information
        \begin{itemize}
            \item \textbf{(New!)} Using new GNN layer type (\href{pytorch-geometric.readthedocs.io/en/2.5.3/generated/torch_geometric.nn.conv.GATConv.html}{GATConv}) that can incorporate edge attributes
            \item Will make slides later on explaining the layer!
        \end{itemize}
        \item Creating associations between tracks and PVs
        \begin{itemize}
            \item Threshold Method: All scores above a threshold associates a track and PV
            \item MaxScore Method: Each track has 1 association based on highest scoring PV
        \end{itemize}
        \item Calculate Metrics!
    \end{itemize}
\end{frame}

\begin{frame}{Training with Edge Attr Classification at \(\left< \mu \right> =60\) (Thr=0.95)}
    \begin{columns}
        \begin{column}{0.45\textwidth}
            \begin{itemize}
                \item Total Reconstructed PVs: 61245
                \item Total Truth PVs: 72189
                \item Total Clean: 50988, Rate: 83.25\%
                \item Total Merged: 7711, Rate: 12.59\%
                \item Total Split: 198, Rate: 0.32\%
                \item Total Fake: 2348, Rate: 3.83\%
            \end{itemize}
        \end{column}

        \begin{column}{0.55\textwidth}
            \begin{figure}
                \centering
                \includegraphics[width=1\linewidth]{july29th/TEST_reco_pv_avg_bar_chart_Threshold_recovertex.pdf}
            \end{figure}
        \end{column}
    \end{columns}
\end{frame}

\begin{frame}{Training with Edge Attr Classification at \(\left< \mu \right> =60\) (MaxScore)}
    \begin{columns}
        \begin{column}{0.45\textwidth}
            \begin{itemize}
                \item Total Reconstructed PVs: 61245
                \item Total Truth PVs: 72189
                \item Total Clean: 51805, Rate: 84.59\%
                \item Total Merged: 7132, Rate: 11.65\%
                \item Total Split: 361, Rate: 0.59\%
                \item Total Fake: 1947, Rate: 3.18\%
            \end{itemize}
        \end{column}

        \begin{column}{0.55\textwidth}
            \begin{figure}
                \centering
                \includegraphics[width=1\linewidth]{july29th/TEST_reco_pv_avg_bar_chart_MaxScore_recovertex.pdf}
            \end{figure}
        \end{column}
    \end{columns}
\end{frame}

\begin{frame}{Comparison of Quality Classification Old vs New}
    \begin{columns}
        \begin{column}{0.5\textwidth}
            \begin{figure}
                \centering
                \includegraphics[width=1\linewidth]{june5th/reco_pv_avg_bar_chart_thr95_recovertex.png}
                \caption{(No Edge Attr) Thr=0.95}
            \end{figure}
        \end{column}

        \begin{column}{0.5\textwidth}
            \begin{figure}
                \centering
                \includegraphics[width=1\linewidth]{july29th/TEST_reco_pv_avg_bar_chart_MaxScore_recovertex.pdf}
                \caption{\textbf{(New!)} MaxScore}
            \end{figure}
        \end{column}
    \end{columns}
\end{frame}


\begin{frame}{Comparison of Quality Classification Old vs New}
    \begin{columns}
        \begin{column}{0.5\textwidth}
            \begin{figure}
                \centering
                \includegraphics[width=1\linewidth]{july29th/fig_14_amvf.pdf}
                \caption{(No Edge Attr) Thr=0.95}
            \end{figure}
        \end{column}

        \begin{column}{0.5\textwidth}
            \begin{figure}
                \centering
                \includegraphics[width=1\linewidth]{july29th/TEST_reco_pv_avg_bar_chart_MaxScore_recovertex.pdf}
                \caption{\textbf{(New!)} MaxScore}
            \end{figure}
        \end{column}
    \end{columns}
\end{frame}

\begin{frame}{Comparison of Reco Vertex Classification: Thresholds (1)}
    \begin{columns}
        \begin{column}{0.5\textwidth}
            \begin{figure}
                \centering
                \includegraphics[width=1\linewidth]{june25th/stacked_pv_classification_5_80_95thr_recovertex.png}
                \caption{(No Edge Attr) Thr = 0.95}
            \end{figure}
        \end{column}

        \begin{column}{0.5\textwidth}
            \begin{figure}
                \centering
                \includegraphics[width=1\linewidth]{july29th/TEST_stacked_pv_classification_0_60_maxscore_recovertex_Threshold.pdf}
                \caption{\textbf{(New!)} Thr = 0.95}
            \end{figure}
        \end{column}
    \end{columns}
\end{frame}

\begin{frame}{Comparison of Reco Vertex Classification: Thresholds (2)}
    \begin{columns}
        \begin{column}{0.5\textwidth}
            \begin{figure}
                \centering
                \includegraphics[width=1\linewidth]{june5th/Screenshot 2025-06-10 at 11.33.19 PM.png}
                \caption{AMVF Pub Note (Fig. 7): \href{https://cds.cern.ch/record/2670380/files/ATL-PHYS-PUB-2019-015.pdf}{Link}}
            \end{figure}
        \end{column}

        \begin{column}{0.5\textwidth}
            \begin{figure}
                \centering
                \includegraphics[width=1\linewidth]{july29th/TEST_stacked_pv_classification_39_80_maxscore_recovertex_Threshold.pdf}
                \caption{\textbf{(New!)} Thr = 0.95}
            \end{figure}
        \end{column}
    \end{columns}
\end{frame}

\begin{frame}{Comparison of Reco Vertex Classification: Thresholds (3)}
    \begin{columns}
        \begin{column}{0.5\textwidth}
            \begin{figure}
                \centering
                \includegraphics[width=1\linewidth]{july29th/TEST_stacked_pv_classification_0_60_maxscore_recovertex_MaxScore.pdf}
                \caption{\textbf{(New!)} MaxScore}
            \end{figure}
        \end{column}

        \begin{column}{0.5\textwidth}
            \begin{figure}
                \centering
                \includegraphics[width=1\linewidth]{july29th/TEST_stacked_pv_classification_0_60_maxscore_recovertex_Threshold.pdf}
                \caption{\textbf{(New!)} Thr = 0.95}
            \end{figure}
        \end{column}
    \end{columns}
\end{frame}

% \begin{frame}{Number of Associations Per Tracks Varying Threshold}
%     \begin{columns}
%         \begin{column}{0.5\textwidth}
%             \begin{figure}
%                 \centering
%                 \includegraphics[width=1\linewidth]{july2nd/gat_casestudy_num_track_assocs_maxscore.png}
%                 \caption{(Sanity Check!) MaxScore}
%             \end{figure}
%         \end{column}

%         \begin{column}{0.5\textwidth}
%             \begin{figure}
%                 \centering
%                 \includegraphics[width=1\linewidth]{july2nd/gat_casestudy_num_track_assocs_thr95.png}
%                 \caption{Average < 1}
%             \end{figure}
%         \end{column}
%     \end{columns}
% \end{frame}

\begin{frame}[plain]
    \centering
    Edge Attributes Case Study
\end{frame}

\begin{frame}{Edge Attributes Case Study: Three Nearby Vertices (1)}
    \begin{columns}
        \begin{column}{0.4\textwidth}
            \begin{itemize}
                \item Reco PV 24: 20 Assoc tracks
                \begin{itemize}
                    \item Old Class: Merged
                    \item \(W_{\text{reco},24} = 20\)
                    \item \(\sum p_T^2 = 3e7\) (MeV\(^2\))
                \end{itemize}
                \item Truth PV 31: 13 Assoc tracks
                \begin{itemize}
                    \item \(W_{\text{truth},11} = 11\)
                    \item \(W_{\text{truth},31}/W_{\text{reco},24} = 0.55\)
                \end{itemize}
                \item Truth PV 38: 10 Assoc tracks
                \begin{itemize}
                    \item \(W_{\text{truth},38} = 4\)
                    \item \(W_{\text{truth},38}/W_{\text{reco},24} = 0.2\)
                \end{itemize}
                \item Truth PV 51: 6 Assoc tracks
                \begin{itemize}
                    \item \(W_{\text{truth},51} = 5\)
                    \item \(W_{\text{truth},51}/W_{\text{reco},24} = 0.25\)
                \end{itemize}
                \item AMVF-style Class: \textbf{Merged}
            \end{itemize}
        \end{column}

        \begin{column}{0.6\textwidth}
            \begin{figure}
                \centering
                \includegraphics[width=1\linewidth]{june25th/casestudy_example_reco_truth_hist_dist_24.png}
            \end{figure}
        \end{column}
    \end{columns}
\end{frame}

\begin{frame}{Edge Attributes Case Study: Three Nearby Vertices (2)}
    \begin{columns}
        \begin{column}{0.4\textwidth}
            \begin{itemize}
                \item Reco PV 24: 6 Assoc tracks
                \begin{itemize}
                    \item Old Class: Clean
                    \item \(W_{\text{reco},24} = 6\)
                    \item \(\sum p_T^2 = 9e6\) (MeV\(^2\))
                \end{itemize}
                \item Truth PV 38: 10 Assoc tracks
                \begin{itemize}
                    \item \(W_{\text{truth},38} = 6\)
                    \item \(W_{\text{truth},38}/W_{\text{reco},25} = 1\)
                \end{itemize}
                \item AMVF-style Class: \textbf{Clean}
            \end{itemize}
        \end{column}

        \begin{column}{0.6\textwidth}
            \begin{figure}
                \centering
                \includegraphics[width=1\linewidth]{june25th/casestudy_example_reco_truth_hist_dist_25.png}
            \end{figure}
        \end{column}
    \end{columns}
\end{frame}

\begin{frame}{Edge Attributes Case Study: Split HS Vertex}
    \begin{columns}
        \begin{column}{0.3\textwidth}
            \begin{itemize}
                \item More of an issue with PV-Finder's output thanso GNN
                \item Causes GNN HS reconstruction to be undefined
            \end{itemize}
        \end{column}

        \begin{column}{0.7\textwidth}
            \begin{figure}
                \centering
                \includegraphics[width=0.9\linewidth]{july9th/casestudy_splitHS.png}
            \end{figure}
        \end{column}
    \end{columns}
\end{frame}

\begin{frame}{Edge Attributes Case Study: HS Selection at High PU Density}
    \begin{columns}
        \begin{column}{0.4\textwidth}
            \begin{itemize}
                \item Reco PV 14: 63 Assoc tracks
                \begin{itemize}
                    \item Old Class: Clean
                    \item \(W_{\text{reco},14} = 63\)
                    \item \(\sum p_T^2 = 55e7\) (MeV\(^2\))
                \end{itemize}
                \item Truth PV 2: 44 Assoc Tracks
                \begin{itemize}
                    \item \(W_{\text{truth},2} = 2\)
                    \item \(W_{\text{truth},2}/W_{\text{reco},14} = 0.03\)
                \end{itemize}
                \item Truth PV 14: 65 Assoc Tracks
                \begin{itemize}
                    \item \(W_{\text{truth},14} = 60\)
                    \item \(W_{\text{truth},2}/W_{\text{reco},14} = 0.95\)
                \end{itemize}
                \item AMVF-style Class: \textbf{Clean}
                \item Not selected as HS because another vertex has higher \(\sum p_T^2\)
            \end{itemize}
        \end{column}

        \begin{column}{0.6\textwidth}
            \begin{figure}
                \centering
                \includegraphics[width=1\linewidth]{july9th/casestudy_clean_notselected_vertex.png}
            \end{figure}
        \end{column}
    \end{columns}
\end{frame}

\begin{frame}[plain]
    \centering
    Hard Scatter Vertex Classification
\end{frame}

\begin{frame}{HS Vertex Quality}
    \begin{itemize}
        \item HS Vertex Quality Classification:
        \begin{itemize}
            \item \textbf{Clean/Matched}: \textbf{CLEAN} reco vertex matching true HS interaction AND \textbf{track weights don't contribute \(> 50 \%\) to other vertices}
            \item \textbf{LowPU}: \textbf{MERGED} reco vertex with \textbf{at least 50\%} of the accumulated \textbf{track weight} coming \textbf{from the simulated HS interaction}
            \item \textbf{HighPU}: \textbf{MERGED} reco vertex the HS interaction contributing between \textbf{1\% and 50\% of the accumulated track weight}"
            \item \textbf{PurePU}: No reco vertex with at least 1\% accumulated track weight from HS interaction"
        \end{itemize}
    \end{itemize}
\end{frame}

\begin{frame}{HS Vertex Quality Category Plot (1)}
    \begin{columns}
        \begin{column}{0.5\textwidth}
            \begin{figure}
                \centering
                \includegraphics[width=1\linewidth]{june25th/fig_07_amvf.png}
                \caption{AMVF Pub Note: \href{https://cds.cern.ch/record/2670380/files/ATL-PHYS-PUB-2019-015.pdf}{Link}}
            \end{figure}
        \end{column}

        \begin{column}{0.5\textwidth}
            \begin{figure}
                \centering
                \includegraphics[width=1\linewidth]{july2nd/rhs_vertex_quality_category_GATCONV_maxscore.png}
                \caption{\textbf{(New!)} MaxScore}
            \end{figure}
        \end{column}
    \end{columns}
\end{frame}

\begin{frame}{HS Vertex Quality Category Plot (2)}
    \begin{columns}
        \begin{column}{0.5\textwidth}
            \begin{figure}
                \centering
                \includegraphics[width=1\linewidth]{july2nd/rhs_vertex_quality_category_GATCONV_thr.png}
                \caption{\textbf{(New!)} Thr = 0.95}
            \end{figure}
        \end{column}

        \begin{column}{0.5\textwidth}
            \begin{figure}
                \centering
                \includegraphics[width=1\linewidth]{july2nd/rhs_vertex_quality_category_GATCONV_maxscore.png}
                \caption{\textbf{(New!)} MaxScore}
            \end{figure}
        \end{column}
    \end{columns}
\end{frame}

\begin{frame}{HS Vertex Reconstruction Efficiency}
    \begin{columns}
        \begin{column}{0.3\textwidth}
            \begin{itemize}
                \item Fraction of events which the HS vertex is reconstructed and classified as Clean/Matched, LowPU, or HighPu
                \item Local pile-up density is the number of vertices within a \(\pm\) 2mm window of a HS vertex
            \end{itemize}
        \end{column}

        \begin{column}{0.7\textwidth}
            \begin{figure}
                \centering
                \includegraphics[width=0.85\linewidth]{july2nd/vertex_reconstruction_efficiency_comparison.png}
            \end{figure}
        \end{column}
    \end{columns}
\end{frame}

\begin{frame}{Comparison: HS Vertex Reconstruction Efficiency}
    \begin{columns}
        \begin{column}{0.5\textwidth}
            \begin{figure}
                \centering
                \includegraphics[width=1\linewidth]{june25th/fig_08a_amvf_reco.png}
                \caption{AMVF Pub Note (Fig. 8): \href{https://cds.cern.ch/record/2670380/files/ATL-PHYS-PUB-2019-015.pdf}{Link}}
            \end{figure}
        \end{column}

        \begin{column}{0.5\textwidth}
            \begin{figure}
                \centering
                \includegraphics[width=1\linewidth]{july2nd/vertex_reconstruction_efficiency_comparison.png}
                \caption{Thr = 0.95}
            \end{figure}
        \end{column}
    \end{columns}
\end{frame}

\begin{frame}{HS Vertex Selection Efficiency}
    \begin{columns}
        \begin{column}{0.3\textwidth}
            \begin{itemize}
                \item Fraction of events where the reconstructed vertex with highest \(\sum p_T^2\) is the one contianing the largest total weight from true HS tracks
                \item \textbf{Currently the only metric that needs improvements!}
            \end{itemize}
        \end{column}

        \begin{column}{0.7\textwidth}
            \begin{figure}
                \centering
                \includegraphics[width=0.85\linewidth]{july2nd/vertex_selection_efficiency_comparison.png}
            \end{figure}
        \end{column}
    \end{columns}
\end{frame}

\begin{frame}{Comparison: HS Vertex Selection Efficiency}
    \begin{columns}
        \begin{column}{0.5\textwidth}
            \begin{figure}
                \centering
                \includegraphics[width=1\linewidth]{june25th/fig_09a_amvf_selection.png}
                \caption{AMVF Pub Note (Fig. 9): \href{https://cds.cern.ch/record/2670380/files/ATL-PHYS-PUB-2019-015.pdf}{Link}}
            \end{figure}
        \end{column}

        \begin{column}{0.5\textwidth}
            \begin{figure}
                \centering
                \includegraphics[width=1\linewidth]{july2nd/vertex_selection_efficiency_comparison.png}
            \end{figure}
        \end{column}
    \end{columns}
\end{frame}

\begin{frame}{Concluding Remarks}
    \begin{itemize}
        \item Demonstrated improved performance using edge attributes and maxscore!
        \item Starting to converge now! Still want last bit of performance for HS selection!
        \item How to continue improving performance...? What is the next direction...?
        \begin{itemize}
            \item Inbuilt GNN HS Selection?
            \item Training GNN using Reco PVs?
        \end{itemize}
        \item \textbf{Conference paper draft in progress for ACAT 2025!}
    \end{itemize}
\end{frame}

\begin{frame}[plain]
    \centering
    TVPD August 1st, 2025
\end{frame}

\begin{frame}{PV-Finder Before}
    \vspace{5pt}
    \centering
    This network performs nicely but is limited by execution time
    \vspace{25pt}
    \begin{figure}
        \centering
        \includegraphics[width=1\linewidth]{APS2025/analyticalkde_pv.png}
        \label{fig:oldPVFinderArchitecture}
    \end{figure}
\end{frame}

\begin{frame}{PV-Finder After}
    \centering
    TVPD presentation on PV-Finder results has more details! (\href{https://indico.cern.ch/event/1534661/}{TVPD, April 4th})
    \begin{figure}
        \centering
        \includegraphics[width=1\linewidth]{july29th/draft_MLPandUNet(1).png}
        \label{fig:oldPVFinderArchitecture}
    \end{figure}
\end{frame}

\begin{frame}{Example End-to-end tracks-to-hist Output}
    \begin{figure}
        \centering
        \includegraphics[width=0.8\linewidth]{june25th/example_end_to_end_with_class.png}
        \label{fig:example_Hists_output_1}
    \end{figure}
\end{frame}

\begin{frame}{\(\Delta z\) Resolution and \(\sigma_{\text{vtx-vtx}}\) Calculation}
    \begin{columns}
        \begin{column}{0.6\linewidth}
            \begin{figure}
                \centering
                \includegraphics[width=0.9\linewidth]{ACAT2025/draft_deltaz_ACAT2025.pdf}
            \end{figure}
        \end{column}

        \begin{column}{0.4\linewidth}
            \begin{table}
                \resizebox{\columnwidth}{!}{%
                \begin{tabular}{c c}
                    \hline
                    Method & \(\sigma_{\text{vtx-vtx}}\) (mm) \\
                    \hline
                    % (PV-Finder Prior) UNet & \(0.23 \pm 0.01\) \\
                    % (PV-Finder Prior) UNet++ & \(0.37 \pm 0.01\) \\
                    (PV-Finder) UNet++ & \(0.4 \pm 0.01\) \\
                    (PV-Finder) UNet & \(0.42 \pm 0.01\) \\
                    AMVF & \(0.76 \pm 0.02\) \\
                    \hline
                \end{tabular}%
                }
            \end{table}
        \end{column}
    \end{columns}
\end{frame}

\begin{frame}{UNet, UNet++, and AMVF: Vertex Classification}
    \begin{figure}
        \centering
        \includegraphics[width=0.5\linewidth]{ACAT2025/draft_avg_per_category_ACAT2025.pdf}
    \end{figure}
\end{frame}

\begin{frame}{UNet, UNet++, and AMVF: Efficiency and FPR}
    \centering
    UNet and UNet++ perform basically identically!
    \begin{columns}
        \begin{column}{0.6\linewidth}
            \begin{figure}
                \centering
                \includegraphics[width=0.9\linewidth]{ACAT2025/draft_ntrks_eff_ACAT2025.pdf}
            \end{figure}
        \end{column}

        \begin{column}{0.4\linewidth}
            \begin{table}
                \resizebox{\columnwidth}{!}{%
                \begin{tabular}{c c c}
                    \hline
                    Model & Efficiency & False Positive Rate (\# / event) \\
                    \hline
                    % (PV-Finder Prior) UNet & 88.7\% & 2.6 \\
                    % (PV-Finder Prior) UNet++ & 94.2\% & 1.5 \\
                    (PV-Finder) UNet++ & 94.2\% & 0.74 \\
                    (PV-Finder) UNet & 94.6\% & 0.72 \\
                    AMVF & 93.9\% & 0.8 \\
                    \hline
                \end{tabular}%
                }
            \end{table}
        \end{column}
    \end{columns}
\end{frame}

% \begin{frame}{Intermission}
%     \begin{itemize}
%         \item Great that PV-Finder can reconstruct more clean/merged PV!
%         \item Tracks to Vertex Association?
%         \begin{itemize}
%             \item Necessary for hard-scatter vertex calculations
%             \item Allows for proper comparison with AMVF
%         \end{itemize}
%         \item \textbf{GNN for TTVA}
%         \begin{itemize}
%             \item Graph Construction
%             \item AMVF-Stype Vertex Classification
%             \item Preliminary Performance
%             \item Case Study
%         \end{itemize}
%     \end{itemize}
% \end{frame}

\begin{frame}{PV-Finder + GNN Fitter}
    \begin{figure}
        \centering
        \includegraphics[width=0.55\linewidth]{april4th/pipeline_pvfinder_gnn.png}
    \end{figure}
\end{frame}

\begin{frame}{GNN Fitter}
    \begin{figure}
        \centering
        \includegraphics[width=1\linewidth]{july29th/draft_graph_architecture.png}
    \end{figure}
\end{frame}

\begin{frame}{Number of Reco Vertices with Quality Classification}
    \begin{columns}
        \begin{column}{0.4\textwidth}
            \begin{itemize}
                \item Using AMVF-style vertex classification
                \item GNN is evaluated using a MaxScore method
            \end{itemize}
        \end{column}
        \begin{column}{0.6\textwidth}
            \begin{figure}
                \centering
                \includegraphics[width=1\linewidth]{ACAT2025/draft_reco_pv_overlaid_comparison_ACAT2025.pdf}
            \end{figure}
        \end{column}
    \end{columns}
\end{frame}

\begin{frame}{Num Reco Vertices vs Num Truth Vertices 20 to 70}
    \begin{columns}
        \begin{column}{0.5\textwidth}
            \begin{figure}
                \centering
                \includegraphics[width=1\linewidth]{ACAT2025/draft_stacked_pv_classification_20_70_maxscore_recovertex_MaxScore_ACAT2025.pdf}
            \end{figure}
        \end{column}

        \begin{column}{0.5\textwidth}
            \begin{figure}
                \centering
                \includegraphics[width=1\linewidth]{ACAT2025/draft_stacked_pv_classification_20_70_maxscore_recovertex_AMVF_ACAT2025.pdf}
            \end{figure}
        \end{column}
    \end{columns}
\end{frame}

\begin{frame}{Concluding Remarks}
    \begin{itemize}
        \item Demonstrated improved performance using edge attributes and maxscore!
        \item Starting to converge now! Still want last bit of performance for HS selection!
        \begin{itemize}
            \item Should I have HS plots for ACAT?
        \end{itemize}
        \item How to continue improving performance...? What is the next direction...?
        \begin{itemize}
            \item Inbuilt GNN HS Selection?
            \item Training GNN using Reco PVs?
        \end{itemize}
        \item \textbf{Plot circulation slides draft for ACAT!}
        \begin{itemize}
            \item Attached in Indico along with these slides!
        \end{itemize}
        \item \textbf{Conference paper draft in progress for ACAT!}
    \end{itemize}
\end{frame}

\begin{frame}[plain]
    \centering
    Inner Tracking CP Plenary, August 25th, 2025
\end{frame}

\begin{frame}{Introduction}
    \begin{itemize}
        \item PV-Finder (\href{https://cds.cern.ch/record/2858348?ln=en}{ATL-PHYS-PUB-2023-011}) is a deep learning algorithm for primary vertex identification
        \item New figures present an updated version of PV-Finder and a new GNN, trained to provide track-to-vertex associations (\href{https://cds.cern.ch/record/2939446/}{ATL-COM-INDET-2025-009})
        \item Have addressed all comments on CDS! Thanks to all the reviewers!
        \item \textbf{Today:} Go over the circulation plots and comments/updates
    \end{itemize}
\end{frame}

\includepdf[pages=-]{PVFinder_ACAT2025_ApprovalPlots-5.pdf}

\end{document}
